% ======================================================================
%  Chapter: Hertz's Forceless Mechanics and the Double Pendulum
%  Include with % ======================================================================
%  Chapter: Hertz's Forceless Mechanics and the Double Pendulum
%  Include with % ======================================================================
%  Chapter: Hertz's Forceless Mechanics and the Double Pendulum
%  Include with % ======================================================================
%  Chapter: Hertz's Forceless Mechanics and the Double Pendulum
%  Include with \include{hertz_chapter} from a `book' preamble that loads
%    amsmath, amssymb, amsthm, mathtools, enumitem, booktabs, graphicx,
%    tikz (+ arrows.meta library), hyperref.
%  Figure file needed in the same directory:  jacobi_curvature.pdf
%  Theorem-type environments are defined below only if not yet defined.
% ======================================================================

\makeatletter
\@ifundefined{theorem}{\theoremstyle{plain}\newtheorem{theorem}{Theorem}[chapter]}{}
\@ifundefined{proposition}{\theoremstyle{plain}\newtheorem{proposition}[theorem]{Proposition}}{}
\@ifundefined{corollary}{\theoremstyle{plain}\newtheorem{corollary}[theorem]{Corollary}}{}
\@ifundefined{definition}{\theoremstyle{definition}\newtheorem{definition}[theorem]{Definition}}{}
\@ifundefined{example}{\theoremstyle{definition}\newtheorem{example}[theorem]{Example}}{}
\@ifundefined{remark}{\theoremstyle{remark}\newtheorem{remark}[theorem]{Remark}}{}
\makeatother

\providecommand{\Rb}{\mathbb{R}}
\providecommand{\Tor}{\mathbb{T}}
\providecommand{\Sph}{\mathbb{S}}
\providecommand{\ud}{\mathrm{d}}
\providecommand{\Lap}{\Delta}
\setlength{\emergencystretch}{2.5em}

\chapter{Hertz's Forceless Mechanics and the Double Pendulum}
\label{ch:hertz-double-pendulum}

\section{A mechanics without forces}
\label{sec:hz-intro}

Newton's \emph{Principia} takes force as a primitive. The concept is so
deeply built into the textbook presentation of mechanics that one forgets
how uneasy its status has always been: a force is neither a displacement
nor a mass, it cannot be seen, and the classical examples---gravitation
across empty space, the tension of a string, the reaction of a table---have
little in common except that each can be written on the right-hand side of
$m\ddot{\mathbf x}=\mathbf F$. In \emph{Die Prinzipien der Mechanik in neuem
Zusammenhange dargestellt}, which appeared in 1894 after Hertz's death
and with a preface by Helmholtz, Heinrich Hertz proposed to remove force
from the foundations altogether \cite{hdp:hertz1894}. His own ``image''
of mechanics rests on three notions only---time, space and mass---and on
a single law, which in the Jones--Walley translation of \S309 reads:
``Every free system persists in its state of rest or of uniform motion in
a straightest path.''

The sentence has two moving parts. A \emph{free} system is one subject to
no forces whatever; it is subject only to rigid connections between the
positions of its parts. A \emph{straightest path} is a purely geometric
notion, the path of least curvature compatible with those connections. What
the Newtonian calls a force is then, in Hertz's account, of one of two
kinds. It is either the reaction of a rigid connection among the visible
bodies, or it is the visible trace of \emph{hidden masses}, rigidly
connected to the visible ones and moving in ways we do not observe
\cite{hdp:hertz1894,hdp:lutzen2005}. Helmholtz's preface to the book already
voices reservations on whether hidden systems of the required kind can
actually be found for the forces of nature, and the reception of the
book was, for that reason, largely respectful and largely unconvinced.

This chapter takes Hertz's program seriously as \emph{mathematics} and
asks what it does when applied to the smallest system in which all its
ingredients meet: the planar double pendulum. The double pendulum has
rigid connections, which are Hertz's native material (the two rods); it has
a force, gravity, which Hertz must explain away; and it has chaos, which
any mechanics must be able to represent. The plan is as follows.

\begin{itemize}[nosep]
  \item Section~\ref{sec:hz-geometry} states Hertz's law precisely and shows
        that, for holonomic connections, it is the statement that motion is
        a geodesic of the mass-weighted configuration space.
  \item Section~\ref{sec:hz-free-dp} shows that the double pendulum
        \emph{without} gravity is a free Hertz system, computes the
        curvature of its configuration space, and finds a hidden Killing
        field.
  \item Section~\ref{sec:hz-hidden} switches gravity on and eliminates it in
        Hertz's way: by a representation theorem that turns any positive
        potential into the kinetic energy of one hidden cyclic coordinate.
  \item Section~\ref{sec:hz-jacobi} gives the second geometrization of the
        same dynamics, the Jacobi metric at fixed energy, and uses it to ask
        where the chaos of the double pendulum sits in the geometry.
  \item Sections~\ref{sec:hz-compare} and~\ref{sec:hz-assess} compare the
        geometrizations and assess what the Hertzian program does and does
        not deliver, in the light of the theme of this book: the equations of
        motion, not the formulations built around them, are the real-world
        content of mechanics.
\end{itemize}

\section{Hertz's law and the geometry of configuration space}
\label{sec:hz-geometry}

\subsection{Systems, connections, free motion}

\begin{definition}[Hertzian material system]\label{def:hz-system}
A \emph{material system} consists of $N$ particles with masses $m_i>0$ and
positions $\mathbf x_i\in\Rb^3$. Its \emph{mass-weighted ambient space} is
$(\Rb^{3N},g)$ with
\[
  g(\xi,\eta)\;=\;\sum_{i=1}^{N}m_i\,\xi_i\!\cdot\eta_i .
\]
A \emph{rigid connection} is a relation that is linear and homogeneous in
the differentials of the coordinates, $\omega_\alpha(x)(\ud x)=0$,
$\alpha=1,\dots,k$, where the $\omega_\alpha$ are Pfaffian forms depending
on position. Write
$D_x=\bigcap_\alpha\ker\omega_\alpha(x)\subset\Rb^{3N}$ for the space of
\emph{admissible velocities}. The connections are \emph{holonomic} if $D$
is integrable, i.e.\ if there is a submanifold $\mathcal N\subset\Rb^{3N}$
through every admissible position with $D=T\mathcal N$; otherwise they are
\emph{nonholonomic} (both terms are Hertz's). A system is \emph{free} if
it is subject to rigid connections only, with no active forces.
\end{definition}

Hertz measures path elements by the same mass-weighted expression
$g(\ud x,\ud x)=\sum m_i\,\ud x_i^2$ (he divides by the total mass, which is
immaterial here). Differentiating an admissible motion $x(t)$, with
$\dot x\in D_x$, along the connections gives, for the acceleration $a=\ddot x$,
\[
  \omega_\alpha(x)(a)\;=\;c_\alpha(x,\dot x),
  \qquad
  c_\alpha:=-\,\partial_\beta\omega_{\alpha,\gamma}\,\dot x^\beta\dot x^\gamma ,
\]
so that at a given state $(x,\dot x)$ the admissible accelerations form an
affine space $A_{(x,\dot x)}=a_0+D_x$.

\begin{definition}[Straightest path]\label{def:hz-straightest}
An admissible path is \emph{straightest} if at every point its curvature
is the least among all admissible paths having the same position and the
same tangent. For a path parametrized by constant speed, curvature is
proportional to $|a|_g$, so this means
\[
  a\;=\;\operatorname*{arg\,min}_{a'\in A_{(x,\dot x)}}\ g(a',a').
\]
\end{definition}

\begin{proposition}[Hertz's law in Newtonian language]\label{prop:hz-law}
Let a free system move along a straightest path. Then:
\begin{enumerate}[label=(\roman*),nosep]
  \item $g(a,w)=0$ for every $w\in D_x$, i.e.\ $a\in D_x^{\perp}$;
  \item the speed $g(\dot x,\dot x)^{1/2}$ is constant;
  \item in components, $m_i\ddot{\mathbf x}_i=\mathbf R_i$ with a reaction
        $R=(\mathbf R_i)$ satisfying $g^{-1}\!R\perp D_x$, that is, the
        equations of d'Alembert--Lagrange with zero active forces;
  \item if the connections are holonomic, $D=T\mathcal N$, then $x(t)$ is a
        geodesic of the induced Riemannian manifold $(\mathcal N,g|_{\mathcal N})$.
\end{enumerate}
\end{proposition}

\begin{proof}
(i) Minimizing $g(a',a')$ over the affine space $a_0+D_x$ yields the
orthogonal projection of the origin onto it, which is orthogonal to the
direction space $D_x$. (ii) $\frac{\ud}{\ud t}g(\dot x,\dot x)=2g(a,\dot x)=0$
because $\dot x\in D_x$. (iii) Write the reaction as $R:=g(a,\cdot)$; then
$R(w)=0$ on $D_x$ is (i). (iv) Let $\nabla$ be the flat connection of
$\Rb^{3N}$. For a curve in $\mathcal N$ the Gauss formula splits $a=\nabla_{\dot x}\dot x$
into a part tangent to $\mathcal N$, which equals the covariant acceleration
$\nabla^{\mathcal N}_{\dot x}\dot x$ of the induced metric, and a normal part.
By (i) the tangential part vanishes.
\end{proof}

\begin{remark}[Gauss, and what is new]
Gauss's principle of least constraint minimizes
\[
  \sum_i m_i\Bigl|\mathbf a_i-\frac{\mathbf F_i}{m_i}\Bigr|^2
\]
over admissible accelerations \cite{hdp:gauss1829}; Hertz's law is the case
$\mathbf F=0$, and for stationary connections without active forces the two
coincide \cite{hdp:eom-hertz}. The mathematical gain is therefore nil; the
gain is conceptual. In Gauss's principle forces are given data and
constraints are restrictions; in Hertz's the forces are to be \emph{derived},
and the whole of dynamics becomes a statement about the metric $g$ and the
distribution $D$.
\end{remark}

\begin{remark}[Nonholonomic connections]
For nonholonomic $D$ a straightest path is in general neither a geodesic of
any induced metric nor a minimizer of length among admissible paths; Hertz
stressed that Hamilton's principle and the principle of least action fail for
such systems \cite{hdp:lutzen2005}. The straightest-path condition (i) is the
Lagrange--d'Alembert equation, not the variational (``vakonomic'') one. The
double pendulum has holonomic connections only, so none of this affects the
rest of the chapter.
\end{remark}

\subsection{Holonomic Hertzian mechanics is geodesic flow}

For holonomic connections, part (iv) says that the entire dynamics of a free
system is the geodesic flow of $(\mathcal N,M)$, where in coordinates
$q^a$ on $\mathcal N$ the \emph{kinetic metric} is
\begin{equation}\label{eq:hz-kinetic-metric}
  M_{ab}(q)\;=\;\sum_i m_i\,\frac{\partial\mathbf x_i}{\partial q^a}\!\cdot\frac{\partial\mathbf x_i}{\partial q^b},
\end{equation}
and the equations of motion are
\begin{equation}\label{eq:hz-geodesic}
  \ddot q^a+\Gamma^a_{bc}(q)\,\dot q^b\dot q^c=0 ,
  \qquad
  \Gamma^a_{bc}=\tfrac12M^{ad}\bigl(\partial_bM_{dc}+\partial_cM_{db}-\partial_dM_{bc}\bigr).
\end{equation}
The ``centrifugal'' and ``Coriolis'' terms of textbook mechanics are the
Christoffel symbols $\Gamma$: they are not forces but the coordinate
expression of the connection of $M$. What remains outside this scheme, and what
Hertz must account for in order to keep his claim of generality, is every
term that is \emph{not} of this form, notably a potential force on the
right-hand side of~\eqref{eq:hz-geodesic}.

\section{The double pendulum without gravity is a free Hertz system}
\label{sec:hz-free-dp}

\begin{figure}[t]
\centering
\begin{tikzpicture}[scale=1.15,>=Stealth]
  % ceiling
  \draw[thick] (-1.0,0) -- (1.0,0);
  \foreach \x in {-0.9,-0.6,...,0.9}{\draw (\x,0) -- (\x-0.15,0.18);}
  \coordinate (O) at (0,0);
  \coordinate (A) at ({2*sin(28)},{-2*cos(28)});
  \coordinate (B) at ({2*sin(28)+1.6*sin(62)},{-2*cos(28)-1.6*cos(62)});
  % verticals
  \draw[dashed,gray] (O) -- (0,-3.2);
  \draw[dashed,gray] (A) -- ++(0,-1.9);
  % rods
  \draw[very thick] (O) -- (A) node[midway,above right=-2pt] {$l_1$};
  \draw[very thick] (A) -- (B) node[midway,above right=-1pt] {$l_2$};
  % bobs
  \fill (O) circle (1.6pt);
  \filldraw[fill=white,thick] (A) circle (0.13) node[right=5pt] {$m_1$};
  \filldraw[fill=white,thick] (B) circle (0.13) node[right=5pt] {$m_2$};
  % angles
  \draw[->] (0,-0.85) arc[start angle=-90,end angle=-62,radius=0.85];
  \node at (0.22,-1.12) {\small$\theta_1$};
  \draw[->] (A)++(0,-0.7) arc[start angle=-90,end angle=-28,radius=0.7];
  \node at ({2*sin(28)+0.42},{-2*cos(28)-0.98}) {\small$\theta_2$};
  % gravity
  \draw[->,thick] (3.2,-0.4) -- (3.2,-1.4) node[midway,right] {$g$};
\end{tikzpicture}
\caption{The planar double pendulum. Both angles are measured from the
downward vertical, so the configuration manifold is the torus
$\Tor^2=\Sph^1\times\Sph^1$ (the rods may swing through the upright position).}
\label{fig:hz-dp}
\end{figure}

Let the two bobs have masses $m_1,m_2$ and positions $\mathbf x_1,\mathbf x_2$
in a vertical plane, joined by massless rigid rods of lengths $l_1,l_2$ to the
pivot and to each other (Figure~\ref{fig:hz-dp}). The mass-weighted ambient
space is $\Rb^4$ with $g=\operatorname{diag}(m_1,m_1,m_2,m_2)$, and the
connections are the two holonomic relations
$|\mathbf x_1|=l_1$, $|\mathbf x_2-\mathbf x_1|=l_2$. They define
$\mathcal N=\Tor^2$, parametrized by
\[
  \mathbf x_1=l_1(\sin\theta_1,-\cos\theta_1),\qquad
  \mathbf x_2=\mathbf x_1+l_2(\sin\theta_2,-\cos\theta_2).
\]

\begin{proposition}[Kinetic metric]\label{prop:hz-dp-metric}
With $\delta:=\theta_1-\theta_2$, the induced metric~\eqref{eq:hz-kinetic-metric} is
\[
  M\;=\;\begin{pmatrix}(m_1+m_2)\,l_1^2 & m_2\,l_1l_2\cos\delta\\[2pt]
                       m_2\,l_1l_2\cos\delta & m_2\,l_2^2\end{pmatrix},
  \qquad
  \det M=m_2\,l_1^2l_2^2\,\bigl(m_1+m_2\sin^2\delta\bigr).
\]
\end{proposition}

\begin{proof}
Differentiate the parametrization: $\partial_{\theta_1}\mathbf x_1=l_1(\cos\theta_1,\sin\theta_1)$,
$\partial_{\theta_1}\mathbf x_2=\partial_{\theta_1}\mathbf x_1$,
$\partial_{\theta_2}\mathbf x_1=0$, $\partial_{\theta_2}\mathbf x_2=l_2(\cos\theta_2,\sin\theta_2)$.
Then $M_{11}=(m_1+m_2)l_1^2$, $M_{22}=m_2l_2^2$ and
$M_{12}=m_2l_1l_2(\cos\theta_1\cos\theta_2+\sin\theta_1\sin\theta_2)=m_2l_1l_2\cos\delta$.
\end{proof}

The metric depends on $\theta_1,\theta_2$ only through $\delta$. Hertz's
law for the gravity-free pendulum is therefore the geodesic
equation~\eqref{eq:hz-geodesic} of $(\Tor^2,M)$, which reads explicitly
\begin{equation}\label{eq:hz-dp-free}
\begin{aligned}
  \ddot\theta_1&=-\frac{m_2\,l_1\dot\theta_1^2\sin\delta\cos\delta+m_2\,l_2\dot\theta_2^2\sin\delta}{l_1\,(m_1+m_2\sin^2\delta)},\\[2pt]
  \ddot\theta_2&=\frac{(m_1+m_2)\,l_1\dot\theta_1^2\sin\delta+m_2\,l_2\dot\theta_2^2\sin\delta\cos\delta}{l_2\,(m_1+m_2\sin^2\delta)} .
\end{aligned}
\end{equation}

\begin{proposition}[Curvature and a hidden symmetry]\label{prop:hz-dp-curvature}
The Gaussian curvature of $(\Tor^2,M)$ is
\begin{equation}\label{eq:hz-KM}
  K_M\;=\;\frac{m_1\cos\delta}{l_1l_2\,\bigl(m_1+m_2\sin^2\delta\bigr)^{2}} .
\end{equation}
The vector field $\partial_{\theta_1}+\partial_{\theta_2}$ is Killing, with Noether
integral
\[
  J=\bigl[(m_1+m_2)l_1^2+m_2l_1l_2\cos\delta\bigr]\dot\theta_1
   +\bigl[m_2l_2^2+m_2l_1l_2\cos\delta\bigr]\dot\theta_2 ,
\]
the total angular momentum about the pivot. Together with the
speed, $J$ makes the geodesic flow of $(\Tor^2,M)$ Liouville integrable.
\end{proposition}

\begin{proof}
Formula~\eqref{eq:hz-KM} follows from $K=R_{1212}/\det M$ with the Christoffel
symbols of Proposition~\ref{prop:hz-dp-metric} (the computation was checked
symbolically). Since $M$ depends on $\delta=\theta_1-\theta_2$ only, it is
invariant under the simultaneous rotation $\theta_i\mapsto\theta_i+s$, whose
generator is $\partial_{\theta_1}+\partial_{\theta_2}$; the conserved quantity
$M(\partial_{\theta_1}+\partial_{\theta_2},\dot\theta)$ is $J$. A $2$-degree-of-freedom
system with two independent commuting integrals (here the speed and $J$) is
integrable.
\end{proof}

\begin{remark}[Negative curvature without chaos]\label{rem:hz-gb-free}
By~\eqref{eq:hz-KM}, $K_M>0$ for $|\delta|<\pi/2$ and $K_M<0$ for $|\delta|>\pi/2$.
This is forced: the Gauss--Bonnet theorem for the torus gives
$\int_{\Tor^2}K_M\,\ud A=2\pi\chi(\Tor^2)=0$, so a non-flat metric on $\Tor^2$ must
have regions of both signs (check directly: $\ud A=\sqrt{\det M}\,\ud\theta_1\ud\theta_2$
and $K_M\sqrt{\det M}\propto\cos\delta\,(m_1+m_2\sin^2\delta)^{-3/2}$ integrates to zero
over $\delta\in[0,2\pi]$). Yet the geodesic flow is integrable. Negative curvature somewhere is thus
\emph{not} an indicator of chaos; we shall need this caveat in
Section~\ref{sec:hz-jacobi}.
\end{remark}

\section{Switching on gravity: hidden masses}
\label{sec:hz-hidden}

With gravity the double pendulum is the Lagrangian system with
$L=\tfrac12M_{ab}\dot\theta^a\dot\theta^b-V$,
\begin{equation}\label{eq:hz-V}
  V(\theta_1,\theta_2)=-(m_1+m_2)\,g\,l_1\cos\theta_1-m_2\,g\,l_2\cos\theta_2 ,
\end{equation}
and its equations of motion are
\begin{equation}\label{eq:hz-dp-grav}
  \ddot\theta^a+\Gamma^a_{bc}\dot\theta^b\dot\theta^c=-M^{ab}\partial_bV ,
\end{equation}
that is, \eqref{eq:hz-dp-free} with the additional terms
\begin{align*}
  \ddot\theta_1&\ni\frac{m_2\,g\sin\theta_2\cos\delta-(m_1+m_2)\,g\sin\theta_1}{l_1\,(m_1+m_2\sin^2\delta)},\\
  \ddot\theta_2&\ni\frac{(m_1+m_2)\,g\,(\sin\theta_1\cos\delta-\sin\theta_2)}{l_2\,(m_1+m_2\sin^2\delta)} .
\end{align*}
The right-hand side of~\eqref{eq:hz-dp-grav} is a force. For a Hertzian it
must be an appearance.

\subsection{Hertz's response and a representation theorem}

Hertz's answer is that forces are the effect of motions we do not see. The
visible system is coupled by rigid connections to a hidden system; the
combined system is free; and the visible variables, which alone are
observed, move as if acted on by forces \cite{hdp:hertz1894,hdp:lutzen2005}.
In the language of later analytical mechanics the mechanism is familiar: a
hidden coordinate that is \emph{cyclic} can be eliminated by the Routh
procedure \cite{hdp:routh1877}, leaving a visible system with an extra
potential term (Helmholtz's monocyclic systems are the thermodynamic
instance \cite{hdp:helmholtz1884}). The following theorem shows that this
device is not merely available but \emph{universal} for potentials.

\begin{theorem}[Cyclic representation of a potential]\label{thm:hz-cyclic}
Let $(Q,M)$ be a Riemannian manifold, $V\in C^\infty(Q)$ with $V>0$, and
$p>0$ a constant. On $Q\times\Sph^1$ with coordinates $(q^a,\varphi)$ put
\begin{equation}\label{eq:hz-warped}
  G=M_{ab}(q)\,\ud q^a\ud q^b+f(q)^2\,\ud\varphi^2,\qquad f=\frac{p}{\sqrt{2V}} .
\end{equation}
\begin{enumerate}[label=(\alph*),nosep]
  \item If $(q(t),\varphi(t))$ is a geodesic of $G$ with $f^2\dot\varphi=p$, then
        $q(t)$ solves $\ddot q^a+\Gamma^a_{bc}\dot q^b\dot q^c=-M^{ab}\partial_bV$ and
        $\tfrac12M_{ab}\dot q^a\dot q^b+V=\tfrac12G(\dot\gamma,\dot\gamma)$ is constant.
  \item Conversely, if $q(t)$ solves the equation in (a), then
        $\varphi(t)=\varphi_0+\int_0^t 2V(q(s))/p\;\ud s$ makes
        $(q,\varphi)$ a geodesic of $G$ with $f^2\dot\varphi=p$.
\end{enumerate}
\end{theorem}

\begin{proof}
Since $G$ is block diagonal and $\varphi$ does not appear in it, the only
Christoffel symbols of $G$ that differ from those of $M$ and from zero are
\[
  \Gamma^a_{\varphi\varphi}=-\tfrac12M^{ab}\partial_b(f^2),\qquad
  \Gamma^\varphi_{a\varphi}=\Gamma^\varphi_{\varphi a}=\partial_af/f ,
\]
while $\Gamma^a_{b\varphi}=\Gamma^\varphi_{ab}=\Gamma^\varphi_{\varphi\varphi}=0$.
The geodesic equations are therefore
\[
  \ddot q^a+\Gamma^a_{bc}\dot q^b\dot q^c-\tfrac12M^{ab}\partial_b(f^2)\,\dot\varphi^2=0,
  \qquad
  \ddot\varphi+2\,\frac{\partial_af}{f}\,\dot q^a\dot\varphi=0 .
\]
The second equation is $\frac{\ud}{\ud t}(f^2\dot\varphi)=0$, the conservation of
the cyclic momentum. Insert $\dot\varphi=p/f^2$ in the first:
\[
  \tfrac12M^{ab}\partial_b(f^2)\,\frac{p^2}{f^4}
  =-\tfrac12M^{ab}\partial_b\!\Bigl(\frac{p^2}{f^2}\Bigr)
  =-M^{ab}\partial_b V ,
\]
because $p^2/(2f^2)=V$. This proves (a); the energy statement follows from
$\tfrac12f^2\dot\varphi^2=p^2/(2f^2)=V$. Part (b) reverses the steps: define
$\varphi$ by $\dot\varphi=p/f^2=2V/p$, so the second geodesic equation holds, and
the first reduces to the given equation.
\end{proof}

The theorem is Routh reduction run backwards: $R=\tfrac12M\dot q\dot q-p^2/(2f^2)$
is the Routhian of $G$, and $p^2/(2f^2)$ is its effective potential.
(A metric of the form $g+f^2\ud\varphi^2$ is also what Kaluza--Klein
theory calls a circle compactification with a scalar $f$ and no gauge field;
the potential $p^2/(2f^2)$ is the usual momentum-induced one.)
The physical reading, in Hertz's terms, is that \emph{potential energy is the
kinetic energy of the hidden rotor}: $V=\tfrac12f^2\dot\varphi^2$.

\begin{example}[The simple pendulum as a Clairaut problem]\label{ex:hz-simple}
For a pendulum of mass $m$ and length $l$, $Q=\Sph^1$ with $M=ml^2\ud\theta^2$ and
$V=mgl(1-\cos\theta)+\varepsilon$, $\varepsilon>0$ (a constant shift changes only
the value of the conserved energy). Then \eqref{eq:hz-warped} is a metric on a
torus of revolution with profile $f(\theta)=p/\sqrt{2V(\theta)}$,
and the pendulum's motion is a geodesic of that surface. The cyclic
momentum $p=f^2\dot\varphi$ is Clairaut's constant: it equals $f$ times the
component of the velocity along the parallel. In terms of arclength
$s=\sqrt m\,l\,\theta$ the Gaussian curvature of the surface is $-f_{ss}/f$.
The oscillating and rotating regimes of the pendulum are the geodesics that
oscillate about, or wind around, the neck of the surface.
\end{example}

\begin{figure}[t]
\centering
\begin{tikzpicture}[scale=1.0]
  \draw[->] (-0.3,-2.0) -- (7.3,-2.0) node[right] {$q\in Q$};
  \draw[thick,domain=0.2:6.8,samples=80,smooth] plot (\x,{0.85+0.4*sin(58*\x)});
  \draw[thick,domain=0.2:6.8,samples=80,smooth] plot (\x,{-(0.85+0.4*sin(58*\x))});
  \foreach \x in {0.6,1.6,2.6,3.6,4.6,5.6,6.6}{
    \pgfmathsetmacro{\r}{0.85+0.4*sin(58*\x)}
    \draw (\x,0) ellipse ({0.17} and {\r});
    \draw[dotted] (\x,-2.0) -- (\x,{-\r});
  }
  \node at (0.6,1.55) {$f(q)$};
  \draw[->] (3.0,1.25) -- (3.0,0.02);
  \node[right] at (3.0,1.35) {\small hidden circle, angle $\varphi$};
\end{tikzpicture}
\caption{Schematic of Theorem~\ref{thm:hz-cyclic}: over each configuration
$q$ of the visible system sits a hidden circle of radius $f(q)=p/\sqrt{2V(q)}$.
Free (geodesic) motion on the total space projects to motion in the potential
$V$; the hidden circle spins faster where the radius is smaller, that is,
where the potential is higher.}
\label{fig:hz-warped}
\end{figure}

\subsection{Application to the double pendulum}

\begin{corollary}[Hertzian double pendulum]\label{cor:hz-dp}
Let $V_c=V+V_0$ with $V$ given by~\eqref{eq:hz-V} and a constant
$V_0>(m_1+m_2)gl_1+m_2gl_2$, and let $M$ be as in
Proposition~\ref{prop:hz-dp-metric}. On $\Tor^3=\Tor^2\times\Sph^1$ the metric
\[
  G=M_{ab}\,\ud\theta^a\ud\theta^b+\frac{p^2}{2V_c(\theta)}\,\ud\varphi^2
\]
makes Hertz's law (geodesic motion with constant speed) equivalent to the
double pendulum~\eqref{eq:hz-dp-grav} in the sense of
Theorem~\ref{thm:hz-cyclic}, with $\tfrac12G(\dot\gamma,\dot\gamma)=E+V_0$,
where $E=\tfrac12M\dot\theta\dot\theta+V$ is the pendulum's energy.
The hidden rotor turns with angular velocity
$\dot\varphi=2V_c(\theta(t))/p$.
\end{corollary}

The constant $V_0$ is necessary only because $f=p/\sqrt{2V_c}$ requires
$V_c>0$ and $V$ attains negative values; it shifts the zero of potential
energy and so, by Theorem~\ref{thm:hz-cyclic}, only the numerical value of the
total geodesic energy. All the double-pendulum quantities in the following
section are unchanged by it.

\subsection{How literal are the hidden masses?}

Hertz demands more than a Riemannian metric: hidden \emph{masses}, with
rigid connections to the visible ones. Three observations clarify what the
representation theorem does and does not give.

\paragraph{(a) A literal rotor, with an inertia correction.}
Let a hidden particle of unit mass be constrained to the position
$\mathbf y=(r(\theta)\cos\varphi,\,r(\theta)\sin\varphi,\,0)$, a point on a circle
whose radius is a prescribed function of the pendulum angles (a cam or linkage
in a real mechanism). The total ambient metric restricted to the constraint manifold is
\[
  G=\bigl(M^{\rm bare}_{ab}+\partial_ar\,\partial_br\bigr)\ud\theta^a\ud\theta^b+r^2\ud\varphi^2 ,
\]
so the hidden mass also contributes to the inertia of the visible variables,
and by Theorem~\ref{thm:hz-cyclic} (its proof uses only $V=p^2/(2f^2)$) the visible motion has effective metric
$M^{\rm bare}+\ud r\otimes\ud r$ and potential $p^2/(2r^2)$. Choosing $r=p/\sqrt{2V_c}$ gives
$\ud r\otimes\ud r=\dfrac{p^2}{8V_c^{3}}\,\ud V\otimes\ud V$, which is $O(V_0^{-3})$ at fixed $p$.
Thus the literal rotor reproduces the double pendulum with a bare inertia that differs from
the observed one by a correction that vanishes as the hidden energy $V_c$ grows relative to
the variation of $V$. Exactness is recovered only if one declares the \emph{effective}
inertia $M$ to be the observed one. The metric $G$ of Corollary~\ref{cor:hz-dp}
encodes precisely this declaration.

\paragraph{(b) Nash embedding: realizability at the level of geometry.}
By Nash's theorem any compact Riemannian manifold embeds isometrically in
some Euclidean space \cite{hdp:nash1956}. Applied to $(\Tor^3,G)$ this exhibits
the Hertzian configuration space as the constraint manifold of a system of
unit-mass coordinates with holonomic rigid connections. But in such a
realization the pendulum angles $\theta^a$ are merely functions on the
constraint manifold, not the positions of two identifiable bobs. The theorem
shows that Hertz's \emph{geometry} is always realizable; it does not show that
the hidden masses can be chosen compatibly with the visible ones.

\paragraph{(c) Non-uniqueness.}
Nothing in the visible motion selects the hidden structure. Replacing $V_c$ by
a different positive function changes only the energy bookkeeping; adding more
hidden rotors gives $V=\sum_kp_k^2/(2f_k^2)$ with many solutions; and the Jacobi
and Eisenhart constructions of Sections~\ref{sec:hz-jacobi}
and~\ref{sec:hz-compare} reproduce the same trajectories with entirely
different hidden content. The Hertzian explanation of gravity is therefore not
testable by observing the pendulum.

\section{A second geometrization: the Jacobi metric}
\label{sec:hz-jacobi}

Hertz enlarges the configuration space. There is also a way to keep its
dimension and to absorb the potential into the metric by a conformal
rescaling, at the price of fixing the energy.

\begin{theorem}[Jacobi]\label{thm:hz-jacobi}
On the Hill region $U_E=\{q\in Q: V(q)<E\}$ let
$J_E=(E-V)\,M$. The solutions of Newton's equations $\ddot q^a+\Gamma^a_{bc}\dot q^b\dot q^c=-M^{ab}\partial_bV$
of energy $E$ are, after the reparametrization $\ud\tau=(E-V)\,\ud t$, the
geodesics of $(U_E,J_E)$.
\end{theorem}

This is the geometric content of Maupertuis's principle: the abbreviated action
\[
  \int\sqrt{2(E-V)}\;\ud s_M
\]
is, up to the factor $\sqrt2$, the $J_E$-length of the path
\cite{hdp:jacobi1866,hdp:arnold1989}. (The relation of this principle to Hamilton's and
Hertz's is treated in the chapter on variational principles.)

\subsection{Curvature of the Jacobi metric}

In two dimensions the stability of nearby geodesics is governed by the Gaussian
curvature: for a unit-speed geodesic of a surface, the normal deviation $\xi$
between neighbouring geodesics obeys $\xi''+K\,\xi=0$ (Jacobi--Levi-Civita
equation), exponentially growing where $K<0$.

\begin{proposition}[Conformal formula]\label{prop:hz-conformal}
For $J_E=\Omega\,M$ with $\Omega=E-V>0$ on a surface $(Q,M)$,
\begin{equation}\label{eq:hz-KJ}
  K_{J_E}\;=\;\frac{1}{\Omega}\Bigl(K_M-\tfrac12\,\Lap_M\ln\Omega\Bigr),
\end{equation}
where $\Lap_M$ is the Laplace--Beltrami operator of $M$. Near a regular point of the
Hill boundary $\{V=E\}$, $K_{J_E}\simeq|\nabla\Omega|_M^2/(2\Omega^3)\to+\infty$.
\end{proposition}

\begin{proof}
For $\hat g=e^{2u}g$ on a surface, $\hat K=e^{-2u}(K-\Lap_gu)$ \cite{hdp:docarmo1992}.
Put $e^{2u}=\Omega$. For the boundary behaviour,
$\Lap\ln\Omega=\Lap\Omega/\Omega-|\nabla\Omega|^2/\Omega^2$, and the second term dominates.
(Formula~\eqref{eq:hz-KJ} was also checked symbolically against the Riemann tensor of $J_E$ for the double pendulum.)
\end{proof}

For the double pendulum $K_M$ is~\eqref{eq:hz-KM}, $\Omega=E-V$ with $V$
from~\eqref{eq:hz-V}, and~\eqref{eq:hz-KJ} gives $K_{J_E}$ in closed form.

\begin{proposition}[Negative curvature is unavoidable above the separatrix energy]\label{prop:hz-gb}
Let $E>V_{\max}=(m_1+m_2)gl_1+m_2gl_2$. Then $U_E=\Tor^2$, $J_E$ is a smooth
Riemannian metric on the torus, and $\int_{\Tor^2}K_{J_E}\,\ud A_{J_E}=0$. Hence, unless
$K_{J_E}\equiv0$, it takes negative values somewhere.
\end{proposition}

\begin{proof}
$V<V_{\max}<E$ everywhere, so $\Omega>0$ on all of $\Tor^2$. Apply Gauss--Bonnet
with $\chi(\Tor^2)=0$. (For $m_1=m_2=l_1=l_2=g=1$ one finds numerically
$K_{J_4}\in[-1.0,\,0.30]$, so $K_{J_E}\not\equiv0$.)
\end{proof}

\subsection{Where the negative curvature sits}

Figure~\ref{fig:hz-curvature} shows the sign of $K_{J_E}$ for the
unit-parameter pendulum ($m_1=m_2=l_1=l_2=g=1$, so $V=-2\cos\theta_1-\cos\theta_2$,
$-3\le V\le3$) at three energies, and Table~\ref{tab:hz-curv} the fraction of
the (Jacobi-)area where it is negative.

\begin{figure}[t]
\centering
\includegraphics[width=\textwidth]{jacobi_curvature.pdf}
\caption{Sign of the Gaussian curvature $K_{J_E}$ of the Jacobi metric of the
double pendulum on the torus $(\theta_1,\theta_2)\in[-\pi,\pi)^2$, for
$m_1=m_2=l_1=l_2=g=1$. Blue: $K_{J_E}>0$; red: $K_{J_E}<0$; grey: inaccessible
region $V\ge E$. At $E=-2$ the accessible region is a disc about the stable
equilibrium and $K_{J_E}>0$ throughout. At $E=1$ a red region has opened; at
$E=4>V_{\max}=3$ the whole torus is accessible and, as
Proposition~\ref{prop:hz-gb} requires, the curvature has both signs.}
\label{fig:hz-curvature}
\end{figure}

\begin{table}[t]
\centering
\begin{tabular}{@{}rcc@{}}
\toprule
$E$ & accessible part of the torus & fraction of $J_E$-area with $K_{J_E}<0$\\
\midrule
$-2$ & $0.13$ & $0$\\
$1$  & $0.69$ & $0.20$\\
$4$  & $1$    & $0.44$\\
\bottomrule
\end{tabular}
\caption{Accessible fraction of the coordinate torus and fraction of the Jacobi area
$\sqrt{\det J_E}\,\ud\theta_1\ud\theta_2$ on which the curvature is negative, for the
unit-parameter double pendulum (grid evaluation, two significant digits).}
\label{tab:hz-curv}
\end{table}

\begin{remark}[What the picture does and does not say]\label{rem:hz-caveat}
The pattern is the expected one. At low energy the motion is a perturbed
small oscillation and the whole accessible region has $K_{J_E}>0$; as energy
rises, negative-curvature regions open and grow, and above the energy at which
the rods can pass through the upright position the topology forces them to exist.
But three cautions are required.
(i) $K<0$ is sufficient for local exponential divergence of nearby geodesics, not necessary:
oscillating positive curvature can produce instability by parametric resonance, which is the
mechanism of the Hannay-angle chapter, and this is how chaos arises in many
systems whose average curvature is positive \cite{hdp:casetti1996,hdp:pettini2007}.
(ii) Remark~\ref{rem:hz-gb-free} shows that a torus metric with negative-curvature regions can
have integrable geodesic flow; the sign of $K$ locates \emph{possible} instability,
it does not by itself establish chaos.
(iii) $J_E$ degenerates at the Hill boundary and its curvature diverges there
(Proposition~\ref{prop:hz-conformal}), so the picture below the separatrix energy is a statement
about the interior of $U_E$ only.
\end{remark}

\section{Three geometrizations side by side}
\label{sec:hz-compare}

A third construction due to Eisenhart lifts the dynamics to a \emph{Lorentzian}
space. On $Q\times\Rb_t\times\Rb_z$ take
\begin{equation}\label{eq:hz-eisenhart}
  G_{\rm E}=M_{ab}\,\ud q^a\ud q^b+2\,\ud t\,\ud z-2V(q)\,\ud t^2 .
\end{equation}
The $t$-equation of the geodesic flow gives $\ddot t=0$, so $t$ is an affine
parameter, and the $q$-equation is $\ddot q^a+\Gamma^a_{bc}\dot q^b\dot q^c=-M^{ab}\partial_bV$ (with
$\Gamma^a_{tt}=M^{ab}\partial_bV$) \cite{hdp:eisenhart1929}. Table~\ref{tab:hz-compare} compares the three.

\begin{table}[t]
\centering
\small
\begin{tabular}{@{}lccc@{}}
\toprule
 & Hertz (hidden rotor) & Jacobi & Eisenhart\\
\midrule
Signature & Riemannian & Riemannian & Lorentzian\\
Dimension & $n+1$ & $n$ & $n+2$\\
Fixed data & hidden momentum $p$ & energy $E$ & none\\
Restriction on $V$ & $V+V_0>0$ & Hill region $V<E$ & none\\
Geodesic parameter & physical time $t$ & $\tau$, $\ud\tau=(E-V)\ud t$ & physical time $t$\\
Where $V$ sits & hidden kinetic energy & conformal factor & $G_{tt}$\\
Hertzian reading & hidden masses & none & none\\
\bottomrule
\end{tabular}
\caption{Three ways to turn the potential force into geometry. Only Hertz's
construction has a mechanical reading in terms of positive masses.}
\label{tab:hz-compare}
\end{table}

% ======================================================================
%  Section insert: Eisenhart's lifts in detail  (+ comparison table)
%  Drop into hertz_chapter.tex in place of the old "Three geometrizations
%  side by side" section (i.e. before \section{Assessment...}).
%
%  Needs: amsmath, amsthm (definition/theorem/proposition/remark/example
%  environments as defined in the chapter), enumitem, booktabs.
%  Cross-references used: thm:hz-cyclic, sec:hz-hidden, sec:hz-jacobi,
%  prop:hz-dp-metric, eq:hz-V, eq:hz-KM, eq:hz-warped.
%  Bibliography keys used: hdp:eisenhart1929, hdp:duval1985,
%  hdp:duval1991, hdp:casetti2000.
% ======================================================================

\section{Eisenhart's lifts in detail}
\label{sec:hz-eisenhart}

In 1928--29 L.~P.~Eisenhart showed that the trajectories of a natural
Lagrangian system can be obtained as geodesics of a metric on a larger space
\cite{hdp:eisenhart1929}. His paper contains \emph{two} constructions,
and they stand in different relations to Hertz. The first is
Lorentzian, works for every potential (including time-dependent ones) and
needs two extra coordinates; the second is Riemannian, needs a positive
potential and one extra coordinate. The second turns out to be
the metric of Theorem~\ref{thm:hz-cyclic}. The first was rediscovered
in 1984--85 by Duval, Burdet, K\"unzle and Perrin, who called the resulting
spaces \emph{Bargmann structures} \cite{hdp:duval1985,hdp:duval1991}, and
the construction is now often called the Eisenhart--Duval lift.

\subsection{The Lorentzian (null) lift}

\begin{definition}[Eisenhart lift]\label{def:hz-eis}
Let $(Q,M)$ be a Riemannian manifold and $V\in C^\infty(Q)$ arbitrary. On
$\widehat Q=Q\times\Rb_t\times\Rb_z$ put
\begin{equation}\label{eq:hz-eis-metric}
  G_{\rm E}=M_{ab}(q)\,\ud q^a\ud q^b+2\,\ud t\,\ud z-2V(q)\,\ud t^2 .
\end{equation}
\end{definition}

The $(t,z)$ block $\bigl(\begin{smallmatrix}-2V&1\\1&0\end{smallmatrix}\bigr)$
has determinant $-1$ whatever $V$ is, so $G_{\rm E}$ is a non-degenerate metric
of Lorentzian signature $(n{+}1,1)$. Its inverse has the components
\begin{equation}\label{eq:hz-eis-inverse}
  G^{ab}=M^{ab},\qquad G^{tz}=1,\qquad G^{zz}=2V,\qquad G^{tt}=G^{ta}=G^{za}=0 .
\end{equation}
The only Christoffel symbols that differ from those of $M$ and from zero are
\begin{equation}\label{eq:hz-eis-christoffel}
  \Gamma^a_{tt}=M^{ab}\partial_bV,\qquad
  \Gamma^z_{at}=\Gamma^z_{ta}=-\partial_aV .
\end{equation}
(For the double pendulum this was checked symbolically for a flat base and
confirmed for the curved base $M$ through the Ricci tensor below.)

\begin{theorem}[Eisenhart]\label{thm:hz-eis}
A curve $\gamma(s)=(q(s),t(s),z(s))$ is a geodesic of $G_{\rm E}$ if and only if
\begin{equation}\label{eq:hz-eis-geod}
  \ddot t=0,\qquad
  \ddot q^a+\Gamma^a_{bc}\dot q^b\dot q^c=-M^{ab}\partial_bV\,\dot t^{\,2},\qquad
  \ddot z=2\,\dot t\,\partial_aV\,\dot q^a .
\end{equation}
Consequently, if $\dot t=1$ (so $s=t$ is physical time):
\begin{enumerate}[label=(\alph*),nosep]
  \item $q(t)$ solves Newton's equations $\ddot q^a+\Gamma^a_{bc}\dot q^b\dot q^c=-M^{ab}\partial_bV$;
  \item $\dot z=2V(q)-\varepsilon$ with a constant $\varepsilon$, so $p_t:=G(\partial_t,\dot\gamma)=-\varepsilon$;
  \item $G_{\rm E}(\dot\gamma,\dot\gamma)=2(E-\varepsilon)$, where $E=\tfrac12M\dot q\dot q+V$ is the energy of $q(t)$;
  \item $\gamma$ is a \emph{null} geodesic exactly when $\varepsilon=E$, and then
        $z(t)=z_0-\int_0^t L\,\ud t'=z_0-S[q]$ with $L=\tfrac12M\dot q\dot q-V$.
\end{enumerate}
Conversely every solution $q(t)$ of Newton's equations lifts, for every
choice of $\varepsilon$, $z_0$ and $t_0$, to a geodesic with $\dot t=1$.
\end{theorem}

\begin{proof}
Equations~\eqref{eq:hz-eis-geod} are the geodesic equations with the symbols
\eqref{eq:hz-eis-christoffel}: $\Gamma^t_{AB}=0$ gives $\ddot t=0$;
$\ddot q^a+\Gamma^a_{bc}\dot q^b\dot q^c+\Gamma^a_{tt}\dot t^2=0$ is the second;
$\ddot z+2\Gamma^z_{at}\dot q^a\dot t=0$ is the third. For $\dot t=1$ the third says
$\ddot z=2\,\ud V/\ud t$, hence $\dot z=2V-\varepsilon$, and $p_t=-2V\dot t+\dot z=-\varepsilon$.
Then
\[
  G(\dot\gamma,\dot\gamma)=M\dot q\dot q+2\dot z-2V
  =2T+2(2V-\varepsilon)-2V=2(T+V-\varepsilon)=2(E-\varepsilon),
\]
which gives (c) and the first half of (d). If $\varepsilon=E$ then
$\dot z=2V-E=V-T=-L$. The converse is immediate by defining
$z$ through $\dot z=2V-\varepsilon$.
\end{proof}

\begin{remark}[The extra coordinate is the action]\label{rem:hz-z-action}
Statement (d) is the geometric content of the lift: along the physical
(null) geodesics the new coordinate $z$ is minus Hamilton's principal
function evaluated along the path. This is no accident. With $\dot t=1$ the
geodesic Lagrangian is $\tfrac12G(\dot\gamma,\dot\gamma)=L+\dot z$, so
$\int\tfrac12G(\dot\gamma,\dot\gamma)\,\ud t=S[q]+z(t_2)-z(t_1)$:
the energy functional of the lift is Hamilton's action up to endpoint terms.
\end{remark}

\subsection{Hamiltonian form and conserved quantities}

The geodesic Hamiltonian $H_{\rm E}=\tfrac12G^{AB}p_Ap_B$ is, by~\eqref{eq:hz-eis-inverse},
\begin{equation}\label{eq:hz-eis-ham}
  H_{\rm E}=\tfrac12M^{ab}p_ap_b+p_tp_z+V(q)\,p_z^{\,2}.
\end{equation}
Both $t$ and $z$ are cyclic for time-independent $V$, so $p_t$ and $p_z$ are conserved.
With $p_z=G(\partial_z,\dot\gamma)=\dot t=1$ and
$H(q,p)=\tfrac12M^{ab}p_ap_b+V$ the physical Hamiltonian, \eqref{eq:hz-eis-ham} becomes
\[
  H_{\rm E}=H(q,p)+p_t ,
\]
and the null condition $H_{\rm E}=0$ is $p_t=-H=-E$: Hamilton's equations of $H_{\rm E}$
are the equations of the \emph{extended phase space}, in which energy is the momentum conjugate
to time and the dynamics is a constraint. The constant $p_z$ plays the role of a mass
(in the Bargmann picture it is the mass or the electric charge of the particle), and the
fixed value $p_z=1$ is what selects Newtonian dynamics with unit
inertia. Hamilton's equations give $\dot t=\partial H_{\rm E}/\partial p_t=p_z=1$ and
$\dot z=\partial H_{\rm E}/\partial p_z=p_t+2V=2V-E$, as found above.

\subsection{Geometry of the lift: a pp-wave}

\begin{proposition}[Curvature of the Eisenhart metric]\label{prop:hz-eis-curv}
\leavevmode
\begin{enumerate}[label=(\roman*),nosep]
  \item $\nabla\,\partial_z=0$: the null vector field $\partial_z$ is covariantly constant.
  \item The Ricci tensor of $G_{\rm E}$ has only the components
        \[
          \operatorname{Ric}_{ab}=\operatorname{Ric}^{M}_{ab},\qquad
          \operatorname{Ric}_{tt}=\Lap_MV .
        \]
\end{enumerate}
\end{proposition}

\begin{proof}
(i) is $\Gamma^A_{Bz}=0$, read off from \eqref{eq:hz-eis-christoffel}. For (ii) I computed the Ricci
tensor symbolically for a flat base with arbitrary $V(x,y)$ and, for the
unit-parameter double pendulum with its curved metric $M$, at random points, where the
only non-vanishing components were $\operatorname{Ric}_{ab}=K_MM_{ab}$ and
$\operatorname{Ric}_{tt}=\Lap_MV$ to rounding error; the general statement is the standard curvature
property of Eisenhart metrics \cite{hdp:casetti2000}.
\end{proof}

A Lorentzian manifold with a covariantly constant null vector whose metric has this
block form is a (generalized) \emph{pp-wave}, and the structure
$(\widehat Q,G_{\rm E},\partial_z)$ is a Bargmann structure \cite{hdp:duval1985,hdp:duval1991}.
Two consequences are worth stating.
\begin{itemize}[nosep]
  \item \emph{Gravity as curvature.} With $\Lap_MV=4\pi G\rho$ for a Newtonian
        potential, $\operatorname{Ric}_{tt}=\Lap_MV$ is the Poisson equation, and Ricci-flatness of the lift
        (flat $M$, harmonic $V$) is Newtonian gravity in vacuum. The lift makes Newtonian
        gravity \emph{look} like general relativity; it is a statement about a Lorentzian
        space of one dimension more than spacetime, not about spacetime.
  \item \emph{Tidal tensor.} The geodesic deviation equation of $G_{\rm E}$ along a lifted
        trajectory, restricted to the $q$-directions, is the usual variational equation of the
        dynamics,
        \[
          \frac{D^2\xi^a}{\ud t^2}+R^{M\,a}{}_{bcd}\,\dot q^b\xi^c\dot q^d
          +M^{ab}\nabla_b\nabla_cV\,\xi^c=0,
        \]
        and the trace of the tidal tensor is $\Lap_MV$ (the origin of $\operatorname{Ric}_{tt}$)
        \cite{hdp:casetti2000}. In contrast to the Jacobi metric of
        Section~\ref{sec:hz-jacobi}, no factor $E-V$ enters and there is no Hill-boundary singularity;
        but the price is that the criterion for exponential instability is the standard
        linear-stability computation, not a purely geometric sign condition.
\end{itemize}

\subsection{Examples}

\begin{example}[Harmonic oscillator]\label{ex:hz-eis-ho}
For $Q=\Rb$, $M=\ud q^2$ and $V=\tfrac12\omega^2q^2$,
$G_{\rm E}=\ud q^2+2\,\ud t\,\ud z-\omega^2q^2\ud t^2$ is a plane wave with $\operatorname{Ric}_{tt}=\omega^2$.
The null geodesics with $\dot t=1$ are $q=A\cos\theta$, $\theta=\omega t+\phi$, and
\[
  z=z_0+\tfrac14A^2\omega\,\sin2\theta ,
\]
which is $z_0-\int L\,\ud t$ with $L=-\tfrac12A^2\omega^2\cos2\theta$. The $z$-motion is a
bounded oscillation at twice the frequency.
\end{example}

\begin{example}[Double pendulum]\label{ex:hz-eis-dp}
With $M$ from Proposition~\ref{prop:hz-dp-metric} and $V$ from~\eqref{eq:hz-V},
\[
\begin{aligned}
  G_{\rm E}={}&(m_1+m_2)\,l_1^2\,\ud\theta_1^2+2m_2l_1l_2\cos(\theta_1-\theta_2)\,\ud\theta_1\ud\theta_2
  +m_2l_2^2\,\ud\theta_2^2\\
  &+2\,\ud t\,\ud z+2g\bigl[(m_1+m_2)l_1\cos\theta_1+m_2l_2\cos\theta_2\bigr]\ud t^2 .
\end{aligned}
\]
This is a four-dimensional Lorentzian manifold $\Tor^2\times\Rb^2$ with
$\operatorname{Ric}_{ab}=K_MM_{ab}$ (with $K_M$ from~\eqref{eq:hz-KM}) and $\operatorname{Ric}_{tt}=\Lap_MV$.
Its Killing fields include $\partial_z$ (null, parallel) and $\partial_t$; the rotation
$\partial_{\theta_1}+\partial_{\theta_2}$ that was a symmetry of the free pendulum is
broken by $V$. The null geodesics with $\dot t=1$ project to the chaotic trajectories of the pendulum,
and $z$ along them is $-S[\theta]$.
\end{example}

\subsection{The Riemannian lift, and its identity with the hidden rotor}

Eisenhart's second construction, for \emph{time-independent potentials that are positive} on $Q$,
uses one extra coordinate and a Riemannian metric of the form
$M+A(q)\,\ud z^2$ with $A$ built from $V$
\cite{hdp:eisenhart1929,hdp:casetti2000}. Taking $A=1/V_c$ with $V_c=V+c>0$,
\begin{equation}\label{eq:hz-eis-riem}
  G_{\rm R}=M_{ab}\,\ud q^a\ud q^b+\frac{1}{V_c(q)}\,\ud z^2 ,
\end{equation}
the Hamiltonian is $\tfrac12M^{ab}p_ap_b+\tfrac12V_c\,p_z^2$, which equals
the physical Hamiltonian $\tfrac12M^{ab}p_ap_b+V_c$ for $p_z=\sqrt2$. This is exactly
the metric~\eqref{eq:hz-warped} of Theorem~\ref{thm:hz-cyclic} with $f^2=p^2/(2V_c)$, $p=\sqrt2$,
and $\varphi=z$. The theorem therefore \emph{is} Eisenhart's Riemannian lift, with a different proof and with
a physical interpretation of the extra coordinate added: Eisenhart's $z$ is Hertz's hidden
cyclic coordinate, and $V=\tfrac12f^2\dot z^2$ is the kinetic energy of the hidden rotor.
(Hertz did not write down the warped metric; the idea that a potential is the kinetic energy of
a cyclic hidden motion is his, and the exact geometric statement is Eisenhart's.)

\subsection{Hertzian or not?}

Does the Lorentzian lift satisfy Hertz's requirements? Only in part. Its dynamics is
free in the sense that the law is geodesic, and no force appears. But Hertz's mass metric is positive
definite: kinetic energy is $\tfrac12\sum m\dot x^2\ge0$, and the straightest path is defined by minimizing
a norm of acceleration. In signature $(n{+}1,1)$ there is no such norm, the physical
trajectories are \emph{null}, and the extra coordinate $z$ is a null direction, not a hidden mass.
The Lorentzian lift is thus forceless but not Hertzian; the Riemannian lift is both.
That the two lifts are available side by side underlines the point of
Section~\ref{sec:hz-hidden}(c): the geometric content of the visible dynamics does not fix its higher-dimensional
realization.

\section{Four geometrizations side by side}
\label{sec:hz-compare}

\begin{table}[t]
\centering
\small
\setlength{\tabcolsep}{4pt}
\begin{tabular}{@{}lccc@{}}
\toprule
 & Hertz / Eisenhart--Riem. & Jacobi & Eisenhart--Lorentz\\
\midrule
Signature & Riemannian & Riemannian & Lorentzian $(n{+}1,1)$\\
Dimension & $n+1$ & $n$ & $n+2$\\
Fixed data & $p_z$ (hidden momentum) & energy $E$ & $p_z=1$, null condition\\
Restriction on $V$ & $V+c>0$ & Hill region $V<E$ & none\\
Geodesic parameter & physical time $t$ & $\tau$, $\ud\tau=(E-V)\ud t$ & physical time $t$\\
Where $V$ sits & hidden kinetic energy & conformal factor & $G_{tt}$\\
Role of extra coordinates & hidden angle $z$ & none & $t$ and null $z=-S$\\
Time-dependent $V$ & no & no & yes\\
Hertzian reading & hidden masses & none & none (null)\\
\bottomrule
\end{tabular}
\caption{Three ways (four constructions, since the first column is Hertz's idea and Eisenhart's
Riemannian metric at once) to turn the potential force into geometry. Only the first has a
mechanical reading in terms of positive masses.}
\label{tab:hz-compare}
\end{table}

Table~\ref{tab:hz-compare} summarizes the comparison. The Lorentzian lift is the only one that
treats time-dependent potentials uniformly, since $t$ is a coordinate of the lifted
space; the Jacobi construction requires conservation of energy and the Riemannian lift
requires a time-independent positive $V$.

\section{Assessment: what Hertz's program delivers}
\label{sec:hz-assess}

\paragraph{What it gains.}
Applied to the double pendulum the program does what it promises, in one
respect: it removes every force from the statement of the law. The tension in
the rods is a reaction of the rigid connection, the centrifugal-looking terms
are Christoffel symbols, and gravity is the kinetic energy of a rotor. The
geometry also yields something new that is invisible in the force picture:
the Gauss--Bonnet argument of Proposition~\ref{prop:hz-gb} explains, without any
numerics, why the Jacobi geometry of the double pendulum above the upright
energy must contain regions in which nearby motions diverge.

\paragraph{What it costs.}
The cost is the hidden structure, and the analysis of Section~\ref{sec:hz-hidden}
shows that it is both inessential and unconstrained. The visible motion
of the pendulum determines the ordinary differential equation~\eqref{eq:hz-dp-grav}
and nothing else; the hidden rotor, the Jacobi metric and the Eisenhart
lift are three unrelated ways of writing that same equation as a geodesic problem,
each in a different dimension and with different auxiliary data. To the extent
that observation of the pendulum is all there is, the hidden masses
are idle wheels. This is the pattern argued throughout this book: the equations
of motion are the primary, real-world content of mechanics, and variational and
geometric formulations are ideal-world devices for finding and organizing them.
The three geometrizations are in this sense charts of one object, each valid
on its own domain and each introducing singularities of its own: the positivity
requirement $V+V_0>0$, the Hill boundary, the Lorentzian signature.

\paragraph{What it cannot reach.}
Hertz's systems are closed and conservative: all forces, including friction,
must be traced to hidden conservative motion. Dissipation therefore enters
only as a statement that the hidden energy is out of sight, and genuine open
systems, such as the damped double pendulum, are outside the scheme as stated;
they are the subject of the chapters on damping and on open systems.
The nonholonomic extension, in which straightest paths differ from shortest
paths, is the other place where Hertz's mechanics is not simply geodesic flow
(Section~\ref{sec:hz-geometry}).

\paragraph{Historical coda.}
The reception was divided between those who found the hidden-mass hypothesis a
barrier to physical content and those who appreciated the geometric form
\cite{hdp:lutzen2005}. What survived was less the philosophy than the geometry:
the representation of constrained mechanics as geodesic flow on a Riemannian
manifold, which is the standard language of the modern geometric treatments.

\section*{Problems}
\addcontentsline{toc}{section}{Problems}

\begin{enumerate}[label=\textbf{\thechapter.\arabic*}.]
  \item Prove Proposition~\ref{prop:hz-law}(i) by Lagrange multipliers: minimize
        $g(a,a)$ subject to $\omega_\alpha(a)=c_\alpha$ and show the multipliers
        are the constraint reactions.
  \item Derive~\eqref{eq:hz-dp-free} directly from Proposition~\ref{prop:hz-law}(iii) with
        the constraints $|\mathbf x_1|=l_1$, $|\mathbf x_2-\mathbf x_1|=l_2$, eliminating the
        multipliers (the rod tensions).
  \item Identify the second integral of the free double pendulum with
        $J$ of Proposition~\ref{prop:hz-dp-curvature}, and find the dependence of the
        period of the $\delta$-motion on $J$ and the speed.
  \item In Example~\ref{ex:hz-simple} show that the surface of revolution is
        smooth, that $K=-f_{ss}/f$ does not depend on $p$, and determine where $K<0$
        (compare Gauss--Bonnet for the torus).
  \item Verify the claim of Section~\ref{sec:hz-hidden}(a) that $\ud r\otimes\ud r=O(V_0^{-3})$
        at fixed $p$, and compute the leading correction to the pendulum's small-oscillation frequency.
  \item Verify that the geodesics of~\eqref{eq:hz-eisenhart} with $\dot t=1$ project to
        solutions of~\eqref{eq:hz-dp-grav}, and find the quantity conserved by the
        null Killing vector $\partial_z$.
  \item Show that for $g\neq0$ the vector field $\partial_{\theta_1}+\partial_{\theta_2}$ is
        \emph{not} a Killing field of $J_E$, so the angular momentum $J$ of
        Proposition~\ref{prop:hz-dp-curvature} is no longer conserved. Why does this loss of
        symmetry not by itself prove that the flow is non-integrable?
\end{enumerate}

\section*{Notes and references}
\addcontentsline{toc}{section}{Notes and references}

\emph{On verification.} The kinetic metric, the geodesic-versus-Lagrange form of the equations of
motion, the explicit accelerations~\eqref{eq:hz-dp-free}--\eqref{eq:hz-dp-grav}, the Christoffel symbols
of the warped metric, the curvature~\eqref{eq:hz-KM} and the conformal
formula~\eqref{eq:hz-KJ} were checked symbolically (SymPy); the curvature map and the table were computed
numerically from~\eqref{eq:hz-KJ} on a $900\times900$ grid, and the Gauss--Bonnet integral at $E=4$
vanishes to rounding error, as it should.

\begingroup
\let\chapter\section
\renewcommand{\bibname}{References}
\begin{thebibliography}{99}
\bibitem{hdp:hertz1894} H.~Hertz, \emph{Die Prinzipien der Mechanik in neuem Zusammenhange dargestellt},
  Gesammelte Werke, Vol.~III (J.~A.~Barth, Leipzig, 1894), with a preface by H.~von Helmholtz.
  English translation: \emph{The Principles of Mechanics Presented in a New Form}, trans.\ D.~E.~Jones and
  J.~T.~Walley (Macmillan, London, 1899; repr.\ Dover, New York, 1956).
\bibitem{hdp:lutzen2005} J.~L\"utzen, \emph{Mechanistic Images in Geometric Form: Heinrich Hertz's
  ``Principles of Mechanics''} (Oxford University Press, Oxford, 2005).
\bibitem{hdp:gauss1829} C.~F.~Gauss, \"Uber ein neues allgemeines Grundgesetz der Mechanik,
  \emph{J.\ reine angew.\ Math.}\ \textbf{4} (1829) 232--235.
\bibitem{hdp:eom-hertz} V.~V.~Rumyantsev, Hertz principle, in \emph{Encyclopedia of Mathematics} (EMS Press).
\bibitem{hdp:routh1877} E.~J.~Routh, \emph{A Treatise on the Stability of a Given State of Motion}
  (Macmillan, London, 1877).
\bibitem{hdp:helmholtz1884} H.~von Helmholtz, Principien der Statik monocyklischer Systeme,
  \emph{J.\ reine angew.\ Math.}\ \textbf{97} (1884) 111--140, 317--336.
\bibitem{hdp:nash1956} J.~Nash, The imbedding problem for Riemannian manifolds,
  \emph{Ann.\ of Math.}\ \textbf{63} (1956) 20--63.
\bibitem{hdp:jacobi1866} C.~G.~J.~Jacobi, \emph{Vorlesungen \"uber Dynamik}, ed.\ A.~Clebsch
  (G.~Reimer, Berlin, 1866).
\bibitem{hdp:arnold1989} V.~I.~Arnold, \emph{Mathematical Methods of Classical Mechanics},
  2nd ed., Graduate Texts in Mathematics \textbf{60} (Springer, New York, 1989).
\bibitem{hdp:docarmo1992} M.~P.~do Carmo, \emph{Riemannian Geometry} (Birkh\"auser, Boston, 1992).
\bibitem{hdp:eisenhart1929} L.~P.~Eisenhart, Dynamical trajectories and geodesics,
  \emph{Ann.\ of Math.}\ \textbf{30} (1929) 591--606.
\bibitem{hdp:casetti1996} L.~Casetti, C.~Clementi and M.~Pettini, Riemannian theory of Hamiltonian
  chaos and Lyapunov exponents, \emph{Phys.\ Rev.\ E} \textbf{54} (1996) 5969--5984.
\bibitem{hdp:pettini2007} M.~Pettini, \emph{Geometry and Topology in Hamiltonian Dynamics and
  Statistical Mechanics} (Springer, New York, 2007).
\end{thebibliography}
\endgroup
 from a `book' preamble that loads
%    amsmath, amssymb, amsthm, mathtools, enumitem, booktabs, graphicx,
%    tikz (+ arrows.meta library), hyperref.
%  Figure file needed in the same directory:  jacobi_curvature.pdf
%  Theorem-type environments are defined below only if not yet defined.
% ======================================================================

\makeatletter
\@ifundefined{theorem}{\theoremstyle{plain}\newtheorem{theorem}{Theorem}[chapter]}{}
\@ifundefined{proposition}{\theoremstyle{plain}\newtheorem{proposition}[theorem]{Proposition}}{}
\@ifundefined{corollary}{\theoremstyle{plain}\newtheorem{corollary}[theorem]{Corollary}}{}
\@ifundefined{definition}{\theoremstyle{definition}\newtheorem{definition}[theorem]{Definition}}{}
\@ifundefined{example}{\theoremstyle{definition}\newtheorem{example}[theorem]{Example}}{}
\@ifundefined{remark}{\theoremstyle{remark}\newtheorem{remark}[theorem]{Remark}}{}
\makeatother

\providecommand{\Rb}{\mathbb{R}}
\providecommand{\Tor}{\mathbb{T}}
\providecommand{\Sph}{\mathbb{S}}
\providecommand{\ud}{\mathrm{d}}
\providecommand{\Lap}{\Delta}
\setlength{\emergencystretch}{2.5em}

\chapter{Hertz's Forceless Mechanics and the Double Pendulum}
\label{ch:hertz-double-pendulum}

\section{A mechanics without forces}
\label{sec:hz-intro}

Newton's \emph{Principia} takes force as a primitive. The concept is so
deeply built into the textbook presentation of mechanics that one forgets
how uneasy its status has always been: a force is neither a displacement
nor a mass, it cannot be seen, and the classical examples---gravitation
across empty space, the tension of a string, the reaction of a table---have
little in common except that each can be written on the right-hand side of
$m\ddot{\mathbf x}=\mathbf F$. In \emph{Die Prinzipien der Mechanik in neuem
Zusammenhange dargestellt}, which appeared in 1894 after Hertz's death
and with a preface by Helmholtz, Heinrich Hertz proposed to remove force
from the foundations altogether \cite{hdp:hertz1894}. His own ``image''
of mechanics rests on three notions only---time, space and mass---and on
a single law, which in the Jones--Walley translation of \S309 reads:
``Every free system persists in its state of rest or of uniform motion in
a straightest path.''

The sentence has two moving parts. A \emph{free} system is one subject to
no forces whatever; it is subject only to rigid connections between the
positions of its parts. A \emph{straightest path} is a purely geometric
notion, the path of least curvature compatible with those connections. What
the Newtonian calls a force is then, in Hertz's account, of one of two
kinds. It is either the reaction of a rigid connection among the visible
bodies, or it is the visible trace of \emph{hidden masses}, rigidly
connected to the visible ones and moving in ways we do not observe
\cite{hdp:hertz1894,hdp:lutzen2005}. Helmholtz's preface to the book already
voices reservations on whether hidden systems of the required kind can
actually be found for the forces of nature, and the reception of the
book was, for that reason, largely respectful and largely unconvinced.

This chapter takes Hertz's program seriously as \emph{mathematics} and
asks what it does when applied to the smallest system in which all its
ingredients meet: the planar double pendulum. The double pendulum has
rigid connections, which are Hertz's native material (the two rods); it has
a force, gravity, which Hertz must explain away; and it has chaos, which
any mechanics must be able to represent. The plan is as follows.

\begin{itemize}[nosep]
  \item Section~\ref{sec:hz-geometry} states Hertz's law precisely and shows
        that, for holonomic connections, it is the statement that motion is
        a geodesic of the mass-weighted configuration space.
  \item Section~\ref{sec:hz-free-dp} shows that the double pendulum
        \emph{without} gravity is a free Hertz system, computes the
        curvature of its configuration space, and finds a hidden Killing
        field.
  \item Section~\ref{sec:hz-hidden} switches gravity on and eliminates it in
        Hertz's way: by a representation theorem that turns any positive
        potential into the kinetic energy of one hidden cyclic coordinate.
  \item Section~\ref{sec:hz-jacobi} gives the second geometrization of the
        same dynamics, the Jacobi metric at fixed energy, and uses it to ask
        where the chaos of the double pendulum sits in the geometry.
  \item Sections~\ref{sec:hz-compare} and~\ref{sec:hz-assess} compare the
        geometrizations and assess what the Hertzian program does and does
        not deliver, in the light of the theme of this book: the equations of
        motion, not the formulations built around them, are the real-world
        content of mechanics.
\end{itemize}

\section{Hertz's law and the geometry of configuration space}
\label{sec:hz-geometry}

\subsection{Systems, connections, free motion}

\begin{definition}[Hertzian material system]\label{def:hz-system}
A \emph{material system} consists of $N$ particles with masses $m_i>0$ and
positions $\mathbf x_i\in\Rb^3$. Its \emph{mass-weighted ambient space} is
$(\Rb^{3N},g)$ with
\[
  g(\xi,\eta)\;=\;\sum_{i=1}^{N}m_i\,\xi_i\!\cdot\eta_i .
\]
A \emph{rigid connection} is a relation that is linear and homogeneous in
the differentials of the coordinates, $\omega_\alpha(x)(\ud x)=0$,
$\alpha=1,\dots,k$, where the $\omega_\alpha$ are Pfaffian forms depending
on position. Write
$D_x=\bigcap_\alpha\ker\omega_\alpha(x)\subset\Rb^{3N}$ for the space of
\emph{admissible velocities}. The connections are \emph{holonomic} if $D$
is integrable, i.e.\ if there is a submanifold $\mathcal N\subset\Rb^{3N}$
through every admissible position with $D=T\mathcal N$; otherwise they are
\emph{nonholonomic} (both terms are Hertz's). A system is \emph{free} if
it is subject to rigid connections only, with no active forces.
\end{definition}

Hertz measures path elements by the same mass-weighted expression
$g(\ud x,\ud x)=\sum m_i\,\ud x_i^2$ (he divides by the total mass, which is
immaterial here). Differentiating an admissible motion $x(t)$, with
$\dot x\in D_x$, along the connections gives, for the acceleration $a=\ddot x$,
\[
  \omega_\alpha(x)(a)\;=\;c_\alpha(x,\dot x),
  \qquad
  c_\alpha:=-\,\partial_\beta\omega_{\alpha,\gamma}\,\dot x^\beta\dot x^\gamma ,
\]
so that at a given state $(x,\dot x)$ the admissible accelerations form an
affine space $A_{(x,\dot x)}=a_0+D_x$.

\begin{definition}[Straightest path]\label{def:hz-straightest}
An admissible path is \emph{straightest} if at every point its curvature
is the least among all admissible paths having the same position and the
same tangent. For a path parametrized by constant speed, curvature is
proportional to $|a|_g$, so this means
\[
  a\;=\;\operatorname*{arg\,min}_{a'\in A_{(x,\dot x)}}\ g(a',a').
\]
\end{definition}

\begin{proposition}[Hertz's law in Newtonian language]\label{prop:hz-law}
Let a free system move along a straightest path. Then:
\begin{enumerate}[label=(\roman*),nosep]
  \item $g(a,w)=0$ for every $w\in D_x$, i.e.\ $a\in D_x^{\perp}$;
  \item the speed $g(\dot x,\dot x)^{1/2}$ is constant;
  \item in components, $m_i\ddot{\mathbf x}_i=\mathbf R_i$ with a reaction
        $R=(\mathbf R_i)$ satisfying $g^{-1}\!R\perp D_x$, that is, the
        equations of d'Alembert--Lagrange with zero active forces;
  \item if the connections are holonomic, $D=T\mathcal N$, then $x(t)$ is a
        geodesic of the induced Riemannian manifold $(\mathcal N,g|_{\mathcal N})$.
\end{enumerate}
\end{proposition}

\begin{proof}
(i) Minimizing $g(a',a')$ over the affine space $a_0+D_x$ yields the
orthogonal projection of the origin onto it, which is orthogonal to the
direction space $D_x$. (ii) $\frac{\ud}{\ud t}g(\dot x,\dot x)=2g(a,\dot x)=0$
because $\dot x\in D_x$. (iii) Write the reaction as $R:=g(a,\cdot)$; then
$R(w)=0$ on $D_x$ is (i). (iv) Let $\nabla$ be the flat connection of
$\Rb^{3N}$. For a curve in $\mathcal N$ the Gauss formula splits $a=\nabla_{\dot x}\dot x$
into a part tangent to $\mathcal N$, which equals the covariant acceleration
$\nabla^{\mathcal N}_{\dot x}\dot x$ of the induced metric, and a normal part.
By (i) the tangential part vanishes.
\end{proof}

\begin{remark}[Gauss, and what is new]
Gauss's principle of least constraint minimizes
\[
  \sum_i m_i\Bigl|\mathbf a_i-\frac{\mathbf F_i}{m_i}\Bigr|^2
\]
over admissible accelerations \cite{hdp:gauss1829}; Hertz's law is the case
$\mathbf F=0$, and for stationary connections without active forces the two
coincide \cite{hdp:eom-hertz}. The mathematical gain is therefore nil; the
gain is conceptual. In Gauss's principle forces are given data and
constraints are restrictions; in Hertz's the forces are to be \emph{derived},
and the whole of dynamics becomes a statement about the metric $g$ and the
distribution $D$.
\end{remark}

\begin{remark}[Nonholonomic connections]
For nonholonomic $D$ a straightest path is in general neither a geodesic of
any induced metric nor a minimizer of length among admissible paths; Hertz
stressed that Hamilton's principle and the principle of least action fail for
such systems \cite{hdp:lutzen2005}. The straightest-path condition (i) is the
Lagrange--d'Alembert equation, not the variational (``vakonomic'') one. The
double pendulum has holonomic connections only, so none of this affects the
rest of the chapter.
\end{remark}

\subsection{Holonomic Hertzian mechanics is geodesic flow}

For holonomic connections, part (iv) says that the entire dynamics of a free
system is the geodesic flow of $(\mathcal N,M)$, where in coordinates
$q^a$ on $\mathcal N$ the \emph{kinetic metric} is
\begin{equation}\label{eq:hz-kinetic-metric}
  M_{ab}(q)\;=\;\sum_i m_i\,\frac{\partial\mathbf x_i}{\partial q^a}\!\cdot\frac{\partial\mathbf x_i}{\partial q^b},
\end{equation}
and the equations of motion are
\begin{equation}\label{eq:hz-geodesic}
  \ddot q^a+\Gamma^a_{bc}(q)\,\dot q^b\dot q^c=0 ,
  \qquad
  \Gamma^a_{bc}=\tfrac12M^{ad}\bigl(\partial_bM_{dc}+\partial_cM_{db}-\partial_dM_{bc}\bigr).
\end{equation}
The ``centrifugal'' and ``Coriolis'' terms of textbook mechanics are the
Christoffel symbols $\Gamma$: they are not forces but the coordinate
expression of the connection of $M$. What remains outside this scheme, and what
Hertz must account for in order to keep his claim of generality, is every
term that is \emph{not} of this form, notably a potential force on the
right-hand side of~\eqref{eq:hz-geodesic}.

\section{The double pendulum without gravity is a free Hertz system}
\label{sec:hz-free-dp}

\begin{figure}[t]
\centering
\begin{tikzpicture}[scale=1.15,>=Stealth]
  % ceiling
  \draw[thick] (-1.0,0) -- (1.0,0);
  \foreach \x in {-0.9,-0.6,...,0.9}{\draw (\x,0) -- (\x-0.15,0.18);}
  \coordinate (O) at (0,0);
  \coordinate (A) at ({2*sin(28)},{-2*cos(28)});
  \coordinate (B) at ({2*sin(28)+1.6*sin(62)},{-2*cos(28)-1.6*cos(62)});
  % verticals
  \draw[dashed,gray] (O) -- (0,-3.2);
  \draw[dashed,gray] (A) -- ++(0,-1.9);
  % rods
  \draw[very thick] (O) -- (A) node[midway,above right=-2pt] {$l_1$};
  \draw[very thick] (A) -- (B) node[midway,above right=-1pt] {$l_2$};
  % bobs
  \fill (O) circle (1.6pt);
  \filldraw[fill=white,thick] (A) circle (0.13) node[right=5pt] {$m_1$};
  \filldraw[fill=white,thick] (B) circle (0.13) node[right=5pt] {$m_2$};
  % angles
  \draw[->] (0,-0.85) arc[start angle=-90,end angle=-62,radius=0.85];
  \node at (0.22,-1.12) {\small$\theta_1$};
  \draw[->] (A)++(0,-0.7) arc[start angle=-90,end angle=-28,radius=0.7];
  \node at ({2*sin(28)+0.42},{-2*cos(28)-0.98}) {\small$\theta_2$};
  % gravity
  \draw[->,thick] (3.2,-0.4) -- (3.2,-1.4) node[midway,right] {$g$};
\end{tikzpicture}
\caption{The planar double pendulum. Both angles are measured from the
downward vertical, so the configuration manifold is the torus
$\Tor^2=\Sph^1\times\Sph^1$ (the rods may swing through the upright position).}
\label{fig:hz-dp}
\end{figure}

Let the two bobs have masses $m_1,m_2$ and positions $\mathbf x_1,\mathbf x_2$
in a vertical plane, joined by massless rigid rods of lengths $l_1,l_2$ to the
pivot and to each other (Figure~\ref{fig:hz-dp}). The mass-weighted ambient
space is $\Rb^4$ with $g=\operatorname{diag}(m_1,m_1,m_2,m_2)$, and the
connections are the two holonomic relations
$|\mathbf x_1|=l_1$, $|\mathbf x_2-\mathbf x_1|=l_2$. They define
$\mathcal N=\Tor^2$, parametrized by
\[
  \mathbf x_1=l_1(\sin\theta_1,-\cos\theta_1),\qquad
  \mathbf x_2=\mathbf x_1+l_2(\sin\theta_2,-\cos\theta_2).
\]

\begin{proposition}[Kinetic metric]\label{prop:hz-dp-metric}
With $\delta:=\theta_1-\theta_2$, the induced metric~\eqref{eq:hz-kinetic-metric} is
\[
  M\;=\;\begin{pmatrix}(m_1+m_2)\,l_1^2 & m_2\,l_1l_2\cos\delta\\[2pt]
                       m_2\,l_1l_2\cos\delta & m_2\,l_2^2\end{pmatrix},
  \qquad
  \det M=m_2\,l_1^2l_2^2\,\bigl(m_1+m_2\sin^2\delta\bigr).
\]
\end{proposition}

\begin{proof}
Differentiate the parametrization: $\partial_{\theta_1}\mathbf x_1=l_1(\cos\theta_1,\sin\theta_1)$,
$\partial_{\theta_1}\mathbf x_2=\partial_{\theta_1}\mathbf x_1$,
$\partial_{\theta_2}\mathbf x_1=0$, $\partial_{\theta_2}\mathbf x_2=l_2(\cos\theta_2,\sin\theta_2)$.
Then $M_{11}=(m_1+m_2)l_1^2$, $M_{22}=m_2l_2^2$ and
$M_{12}=m_2l_1l_2(\cos\theta_1\cos\theta_2+\sin\theta_1\sin\theta_2)=m_2l_1l_2\cos\delta$.
\end{proof}

The metric depends on $\theta_1,\theta_2$ only through $\delta$. Hertz's
law for the gravity-free pendulum is therefore the geodesic
equation~\eqref{eq:hz-geodesic} of $(\Tor^2,M)$, which reads explicitly
\begin{equation}\label{eq:hz-dp-free}
\begin{aligned}
  \ddot\theta_1&=-\frac{m_2\,l_1\dot\theta_1^2\sin\delta\cos\delta+m_2\,l_2\dot\theta_2^2\sin\delta}{l_1\,(m_1+m_2\sin^2\delta)},\\[2pt]
  \ddot\theta_2&=\frac{(m_1+m_2)\,l_1\dot\theta_1^2\sin\delta+m_2\,l_2\dot\theta_2^2\sin\delta\cos\delta}{l_2\,(m_1+m_2\sin^2\delta)} .
\end{aligned}
\end{equation}

\begin{proposition}[Curvature and a hidden symmetry]\label{prop:hz-dp-curvature}
The Gaussian curvature of $(\Tor^2,M)$ is
\begin{equation}\label{eq:hz-KM}
  K_M\;=\;\frac{m_1\cos\delta}{l_1l_2\,\bigl(m_1+m_2\sin^2\delta\bigr)^{2}} .
\end{equation}
The vector field $\partial_{\theta_1}+\partial_{\theta_2}$ is Killing, with Noether
integral
\[
  J=\bigl[(m_1+m_2)l_1^2+m_2l_1l_2\cos\delta\bigr]\dot\theta_1
   +\bigl[m_2l_2^2+m_2l_1l_2\cos\delta\bigr]\dot\theta_2 ,
\]
the total angular momentum about the pivot. Together with the
speed, $J$ makes the geodesic flow of $(\Tor^2,M)$ Liouville integrable.
\end{proposition}

\begin{proof}
Formula~\eqref{eq:hz-KM} follows from $K=R_{1212}/\det M$ with the Christoffel
symbols of Proposition~\ref{prop:hz-dp-metric} (the computation was checked
symbolically). Since $M$ depends on $\delta=\theta_1-\theta_2$ only, it is
invariant under the simultaneous rotation $\theta_i\mapsto\theta_i+s$, whose
generator is $\partial_{\theta_1}+\partial_{\theta_2}$; the conserved quantity
$M(\partial_{\theta_1}+\partial_{\theta_2},\dot\theta)$ is $J$. A $2$-degree-of-freedom
system with two independent commuting integrals (here the speed and $J$) is
integrable.
\end{proof}

\begin{remark}[Negative curvature without chaos]\label{rem:hz-gb-free}
By~\eqref{eq:hz-KM}, $K_M>0$ for $|\delta|<\pi/2$ and $K_M<0$ for $|\delta|>\pi/2$.
This is forced: the Gauss--Bonnet theorem for the torus gives
$\int_{\Tor^2}K_M\,\ud A=2\pi\chi(\Tor^2)=0$, so a non-flat metric on $\Tor^2$ must
have regions of both signs (check directly: $\ud A=\sqrt{\det M}\,\ud\theta_1\ud\theta_2$
and $K_M\sqrt{\det M}\propto\cos\delta\,(m_1+m_2\sin^2\delta)^{-3/2}$ integrates to zero
over $\delta\in[0,2\pi]$). Yet the geodesic flow is integrable. Negative curvature somewhere is thus
\emph{not} an indicator of chaos; we shall need this caveat in
Section~\ref{sec:hz-jacobi}.
\end{remark}

\section{Switching on gravity: hidden masses}
\label{sec:hz-hidden}

With gravity the double pendulum is the Lagrangian system with
$L=\tfrac12M_{ab}\dot\theta^a\dot\theta^b-V$,
\begin{equation}\label{eq:hz-V}
  V(\theta_1,\theta_2)=-(m_1+m_2)\,g\,l_1\cos\theta_1-m_2\,g\,l_2\cos\theta_2 ,
\end{equation}
and its equations of motion are
\begin{equation}\label{eq:hz-dp-grav}
  \ddot\theta^a+\Gamma^a_{bc}\dot\theta^b\dot\theta^c=-M^{ab}\partial_bV ,
\end{equation}
that is, \eqref{eq:hz-dp-free} with the additional terms
\begin{align*}
  \ddot\theta_1&\ni\frac{m_2\,g\sin\theta_2\cos\delta-(m_1+m_2)\,g\sin\theta_1}{l_1\,(m_1+m_2\sin^2\delta)},\\
  \ddot\theta_2&\ni\frac{(m_1+m_2)\,g\,(\sin\theta_1\cos\delta-\sin\theta_2)}{l_2\,(m_1+m_2\sin^2\delta)} .
\end{align*}
The right-hand side of~\eqref{eq:hz-dp-grav} is a force. For a Hertzian it
must be an appearance.

\subsection{Hertz's response and a representation theorem}

Hertz's answer is that forces are the effect of motions we do not see. The
visible system is coupled by rigid connections to a hidden system; the
combined system is free; and the visible variables, which alone are
observed, move as if acted on by forces \cite{hdp:hertz1894,hdp:lutzen2005}.
In the language of later analytical mechanics the mechanism is familiar: a
hidden coordinate that is \emph{cyclic} can be eliminated by the Routh
procedure \cite{hdp:routh1877}, leaving a visible system with an extra
potential term (Helmholtz's monocyclic systems are the thermodynamic
instance \cite{hdp:helmholtz1884}). The following theorem shows that this
device is not merely available but \emph{universal} for potentials.

\begin{theorem}[Cyclic representation of a potential]\label{thm:hz-cyclic}
Let $(Q,M)$ be a Riemannian manifold, $V\in C^\infty(Q)$ with $V>0$, and
$p>0$ a constant. On $Q\times\Sph^1$ with coordinates $(q^a,\varphi)$ put
\begin{equation}\label{eq:hz-warped}
  G=M_{ab}(q)\,\ud q^a\ud q^b+f(q)^2\,\ud\varphi^2,\qquad f=\frac{p}{\sqrt{2V}} .
\end{equation}
\begin{enumerate}[label=(\alph*),nosep]
  \item If $(q(t),\varphi(t))$ is a geodesic of $G$ with $f^2\dot\varphi=p$, then
        $q(t)$ solves $\ddot q^a+\Gamma^a_{bc}\dot q^b\dot q^c=-M^{ab}\partial_bV$ and
        $\tfrac12M_{ab}\dot q^a\dot q^b+V=\tfrac12G(\dot\gamma,\dot\gamma)$ is constant.
  \item Conversely, if $q(t)$ solves the equation in (a), then
        $\varphi(t)=\varphi_0+\int_0^t 2V(q(s))/p\;\ud s$ makes
        $(q,\varphi)$ a geodesic of $G$ with $f^2\dot\varphi=p$.
\end{enumerate}
\end{theorem}

\begin{proof}
Since $G$ is block diagonal and $\varphi$ does not appear in it, the only
Christoffel symbols of $G$ that differ from those of $M$ and from zero are
\[
  \Gamma^a_{\varphi\varphi}=-\tfrac12M^{ab}\partial_b(f^2),\qquad
  \Gamma^\varphi_{a\varphi}=\Gamma^\varphi_{\varphi a}=\partial_af/f ,
\]
while $\Gamma^a_{b\varphi}=\Gamma^\varphi_{ab}=\Gamma^\varphi_{\varphi\varphi}=0$.
The geodesic equations are therefore
\[
  \ddot q^a+\Gamma^a_{bc}\dot q^b\dot q^c-\tfrac12M^{ab}\partial_b(f^2)\,\dot\varphi^2=0,
  \qquad
  \ddot\varphi+2\,\frac{\partial_af}{f}\,\dot q^a\dot\varphi=0 .
\]
The second equation is $\frac{\ud}{\ud t}(f^2\dot\varphi)=0$, the conservation of
the cyclic momentum. Insert $\dot\varphi=p/f^2$ in the first:
\[
  \tfrac12M^{ab}\partial_b(f^2)\,\frac{p^2}{f^4}
  =-\tfrac12M^{ab}\partial_b\!\Bigl(\frac{p^2}{f^2}\Bigr)
  =-M^{ab}\partial_b V ,
\]
because $p^2/(2f^2)=V$. This proves (a); the energy statement follows from
$\tfrac12f^2\dot\varphi^2=p^2/(2f^2)=V$. Part (b) reverses the steps: define
$\varphi$ by $\dot\varphi=p/f^2=2V/p$, so the second geodesic equation holds, and
the first reduces to the given equation.
\end{proof}

The theorem is Routh reduction run backwards: $R=\tfrac12M\dot q\dot q-p^2/(2f^2)$
is the Routhian of $G$, and $p^2/(2f^2)$ is its effective potential.
(A metric of the form $g+f^2\ud\varphi^2$ is also what Kaluza--Klein
theory calls a circle compactification with a scalar $f$ and no gauge field;
the potential $p^2/(2f^2)$ is the usual momentum-induced one.)
The physical reading, in Hertz's terms, is that \emph{potential energy is the
kinetic energy of the hidden rotor}: $V=\tfrac12f^2\dot\varphi^2$.

\begin{example}[The simple pendulum as a Clairaut problem]\label{ex:hz-simple}
For a pendulum of mass $m$ and length $l$, $Q=\Sph^1$ with $M=ml^2\ud\theta^2$ and
$V=mgl(1-\cos\theta)+\varepsilon$, $\varepsilon>0$ (a constant shift changes only
the value of the conserved energy). Then \eqref{eq:hz-warped} is a metric on a
torus of revolution with profile $f(\theta)=p/\sqrt{2V(\theta)}$,
and the pendulum's motion is a geodesic of that surface. The cyclic
momentum $p=f^2\dot\varphi$ is Clairaut's constant: it equals $f$ times the
component of the velocity along the parallel. In terms of arclength
$s=\sqrt m\,l\,\theta$ the Gaussian curvature of the surface is $-f_{ss}/f$.
The oscillating and rotating regimes of the pendulum are the geodesics that
oscillate about, or wind around, the neck of the surface.
\end{example}

\begin{figure}[t]
\centering
\begin{tikzpicture}[scale=1.0]
  \draw[->] (-0.3,-2.0) -- (7.3,-2.0) node[right] {$q\in Q$};
  \draw[thick,domain=0.2:6.8,samples=80,smooth] plot (\x,{0.85+0.4*sin(58*\x)});
  \draw[thick,domain=0.2:6.8,samples=80,smooth] plot (\x,{-(0.85+0.4*sin(58*\x))});
  \foreach \x in {0.6,1.6,2.6,3.6,4.6,5.6,6.6}{
    \pgfmathsetmacro{\r}{0.85+0.4*sin(58*\x)}
    \draw (\x,0) ellipse ({0.17} and {\r});
    \draw[dotted] (\x,-2.0) -- (\x,{-\r});
  }
  \node at (0.6,1.55) {$f(q)$};
  \draw[->] (3.0,1.25) -- (3.0,0.02);
  \node[right] at (3.0,1.35) {\small hidden circle, angle $\varphi$};
\end{tikzpicture}
\caption{Schematic of Theorem~\ref{thm:hz-cyclic}: over each configuration
$q$ of the visible system sits a hidden circle of radius $f(q)=p/\sqrt{2V(q)}$.
Free (geodesic) motion on the total space projects to motion in the potential
$V$; the hidden circle spins faster where the radius is smaller, that is,
where the potential is higher.}
\label{fig:hz-warped}
\end{figure}

\subsection{Application to the double pendulum}

\begin{corollary}[Hertzian double pendulum]\label{cor:hz-dp}
Let $V_c=V+V_0$ with $V$ given by~\eqref{eq:hz-V} and a constant
$V_0>(m_1+m_2)gl_1+m_2gl_2$, and let $M$ be as in
Proposition~\ref{prop:hz-dp-metric}. On $\Tor^3=\Tor^2\times\Sph^1$ the metric
\[
  G=M_{ab}\,\ud\theta^a\ud\theta^b+\frac{p^2}{2V_c(\theta)}\,\ud\varphi^2
\]
makes Hertz's law (geodesic motion with constant speed) equivalent to the
double pendulum~\eqref{eq:hz-dp-grav} in the sense of
Theorem~\ref{thm:hz-cyclic}, with $\tfrac12G(\dot\gamma,\dot\gamma)=E+V_0$,
where $E=\tfrac12M\dot\theta\dot\theta+V$ is the pendulum's energy.
The hidden rotor turns with angular velocity
$\dot\varphi=2V_c(\theta(t))/p$.
\end{corollary}

The constant $V_0$ is necessary only because $f=p/\sqrt{2V_c}$ requires
$V_c>0$ and $V$ attains negative values; it shifts the zero of potential
energy and so, by Theorem~\ref{thm:hz-cyclic}, only the numerical value of the
total geodesic energy. All the double-pendulum quantities in the following
section are unchanged by it.

\subsection{How literal are the hidden masses?}

Hertz demands more than a Riemannian metric: hidden \emph{masses}, with
rigid connections to the visible ones. Three observations clarify what the
representation theorem does and does not give.

\paragraph{(a) A literal rotor, with an inertia correction.}
Let a hidden particle of unit mass be constrained to the position
$\mathbf y=(r(\theta)\cos\varphi,\,r(\theta)\sin\varphi,\,0)$, a point on a circle
whose radius is a prescribed function of the pendulum angles (a cam or linkage
in a real mechanism). The total ambient metric restricted to the constraint manifold is
\[
  G=\bigl(M^{\rm bare}_{ab}+\partial_ar\,\partial_br\bigr)\ud\theta^a\ud\theta^b+r^2\ud\varphi^2 ,
\]
so the hidden mass also contributes to the inertia of the visible variables,
and by Theorem~\ref{thm:hz-cyclic} (its proof uses only $V=p^2/(2f^2)$) the visible motion has effective metric
$M^{\rm bare}+\ud r\otimes\ud r$ and potential $p^2/(2r^2)$. Choosing $r=p/\sqrt{2V_c}$ gives
$\ud r\otimes\ud r=\dfrac{p^2}{8V_c^{3}}\,\ud V\otimes\ud V$, which is $O(V_0^{-3})$ at fixed $p$.
Thus the literal rotor reproduces the double pendulum with a bare inertia that differs from
the observed one by a correction that vanishes as the hidden energy $V_c$ grows relative to
the variation of $V$. Exactness is recovered only if one declares the \emph{effective}
inertia $M$ to be the observed one. The metric $G$ of Corollary~\ref{cor:hz-dp}
encodes precisely this declaration.

\paragraph{(b) Nash embedding: realizability at the level of geometry.}
By Nash's theorem any compact Riemannian manifold embeds isometrically in
some Euclidean space \cite{hdp:nash1956}. Applied to $(\Tor^3,G)$ this exhibits
the Hertzian configuration space as the constraint manifold of a system of
unit-mass coordinates with holonomic rigid connections. But in such a
realization the pendulum angles $\theta^a$ are merely functions on the
constraint manifold, not the positions of two identifiable bobs. The theorem
shows that Hertz's \emph{geometry} is always realizable; it does not show that
the hidden masses can be chosen compatibly with the visible ones.

\paragraph{(c) Non-uniqueness.}
Nothing in the visible motion selects the hidden structure. Replacing $V_c$ by
a different positive function changes only the energy bookkeeping; adding more
hidden rotors gives $V=\sum_kp_k^2/(2f_k^2)$ with many solutions; and the Jacobi
and Eisenhart constructions of Sections~\ref{sec:hz-jacobi}
and~\ref{sec:hz-compare} reproduce the same trajectories with entirely
different hidden content. The Hertzian explanation of gravity is therefore not
testable by observing the pendulum.

\section{A second geometrization: the Jacobi metric}
\label{sec:hz-jacobi}

Hertz enlarges the configuration space. There is also a way to keep its
dimension and to absorb the potential into the metric by a conformal
rescaling, at the price of fixing the energy.

\begin{theorem}[Jacobi]\label{thm:hz-jacobi}
On the Hill region $U_E=\{q\in Q: V(q)<E\}$ let
$J_E=(E-V)\,M$. The solutions of Newton's equations $\ddot q^a+\Gamma^a_{bc}\dot q^b\dot q^c=-M^{ab}\partial_bV$
of energy $E$ are, after the reparametrization $\ud\tau=(E-V)\,\ud t$, the
geodesics of $(U_E,J_E)$.
\end{theorem}

This is the geometric content of Maupertuis's principle: the abbreviated action
\[
  \int\sqrt{2(E-V)}\;\ud s_M
\]
is, up to the factor $\sqrt2$, the $J_E$-length of the path
\cite{hdp:jacobi1866,hdp:arnold1989}. (The relation of this principle to Hamilton's and
Hertz's is treated in the chapter on variational principles.)

\subsection{Curvature of the Jacobi metric}

In two dimensions the stability of nearby geodesics is governed by the Gaussian
curvature: for a unit-speed geodesic of a surface, the normal deviation $\xi$
between neighbouring geodesics obeys $\xi''+K\,\xi=0$ (Jacobi--Levi-Civita
equation), exponentially growing where $K<0$.

\begin{proposition}[Conformal formula]\label{prop:hz-conformal}
For $J_E=\Omega\,M$ with $\Omega=E-V>0$ on a surface $(Q,M)$,
\begin{equation}\label{eq:hz-KJ}
  K_{J_E}\;=\;\frac{1}{\Omega}\Bigl(K_M-\tfrac12\,\Lap_M\ln\Omega\Bigr),
\end{equation}
where $\Lap_M$ is the Laplace--Beltrami operator of $M$. Near a regular point of the
Hill boundary $\{V=E\}$, $K_{J_E}\simeq|\nabla\Omega|_M^2/(2\Omega^3)\to+\infty$.
\end{proposition}

\begin{proof}
For $\hat g=e^{2u}g$ on a surface, $\hat K=e^{-2u}(K-\Lap_gu)$ \cite{hdp:docarmo1992}.
Put $e^{2u}=\Omega$. For the boundary behaviour,
$\Lap\ln\Omega=\Lap\Omega/\Omega-|\nabla\Omega|^2/\Omega^2$, and the second term dominates.
(Formula~\eqref{eq:hz-KJ} was also checked symbolically against the Riemann tensor of $J_E$ for the double pendulum.)
\end{proof}

For the double pendulum $K_M$ is~\eqref{eq:hz-KM}, $\Omega=E-V$ with $V$
from~\eqref{eq:hz-V}, and~\eqref{eq:hz-KJ} gives $K_{J_E}$ in closed form.

\begin{proposition}[Negative curvature is unavoidable above the separatrix energy]\label{prop:hz-gb}
Let $E>V_{\max}=(m_1+m_2)gl_1+m_2gl_2$. Then $U_E=\Tor^2$, $J_E$ is a smooth
Riemannian metric on the torus, and $\int_{\Tor^2}K_{J_E}\,\ud A_{J_E}=0$. Hence, unless
$K_{J_E}\equiv0$, it takes negative values somewhere.
\end{proposition}

\begin{proof}
$V<V_{\max}<E$ everywhere, so $\Omega>0$ on all of $\Tor^2$. Apply Gauss--Bonnet
with $\chi(\Tor^2)=0$. (For $m_1=m_2=l_1=l_2=g=1$ one finds numerically
$K_{J_4}\in[-1.0,\,0.30]$, so $K_{J_E}\not\equiv0$.)
\end{proof}

\subsection{Where the negative curvature sits}

Figure~\ref{fig:hz-curvature} shows the sign of $K_{J_E}$ for the
unit-parameter pendulum ($m_1=m_2=l_1=l_2=g=1$, so $V=-2\cos\theta_1-\cos\theta_2$,
$-3\le V\le3$) at three energies, and Table~\ref{tab:hz-curv} the fraction of
the (Jacobi-)area where it is negative.

\begin{figure}[t]
\centering
\includegraphics[width=\textwidth]{jacobi_curvature.pdf}
\caption{Sign of the Gaussian curvature $K_{J_E}$ of the Jacobi metric of the
double pendulum on the torus $(\theta_1,\theta_2)\in[-\pi,\pi)^2$, for
$m_1=m_2=l_1=l_2=g=1$. Blue: $K_{J_E}>0$; red: $K_{J_E}<0$; grey: inaccessible
region $V\ge E$. At $E=-2$ the accessible region is a disc about the stable
equilibrium and $K_{J_E}>0$ throughout. At $E=1$ a red region has opened; at
$E=4>V_{\max}=3$ the whole torus is accessible and, as
Proposition~\ref{prop:hz-gb} requires, the curvature has both signs.}
\label{fig:hz-curvature}
\end{figure}

\begin{table}[t]
\centering
\begin{tabular}{@{}rcc@{}}
\toprule
$E$ & accessible part of the torus & fraction of $J_E$-area with $K_{J_E}<0$\\
\midrule
$-2$ & $0.13$ & $0$\\
$1$  & $0.69$ & $0.20$\\
$4$  & $1$    & $0.44$\\
\bottomrule
\end{tabular}
\caption{Accessible fraction of the coordinate torus and fraction of the Jacobi area
$\sqrt{\det J_E}\,\ud\theta_1\ud\theta_2$ on which the curvature is negative, for the
unit-parameter double pendulum (grid evaluation, two significant digits).}
\label{tab:hz-curv}
\end{table}

\begin{remark}[What the picture does and does not say]\label{rem:hz-caveat}
The pattern is the expected one. At low energy the motion is a perturbed
small oscillation and the whole accessible region has $K_{J_E}>0$; as energy
rises, negative-curvature regions open and grow, and above the energy at which
the rods can pass through the upright position the topology forces them to exist.
But three cautions are required.
(i) $K<0$ is sufficient for local exponential divergence of nearby geodesics, not necessary:
oscillating positive curvature can produce instability by parametric resonance, which is the
mechanism of the Hannay-angle chapter, and this is how chaos arises in many
systems whose average curvature is positive \cite{hdp:casetti1996,hdp:pettini2007}.
(ii) Remark~\ref{rem:hz-gb-free} shows that a torus metric with negative-curvature regions can
have integrable geodesic flow; the sign of $K$ locates \emph{possible} instability,
it does not by itself establish chaos.
(iii) $J_E$ degenerates at the Hill boundary and its curvature diverges there
(Proposition~\ref{prop:hz-conformal}), so the picture below the separatrix energy is a statement
about the interior of $U_E$ only.
\end{remark}

\section{Three geometrizations side by side}
\label{sec:hz-compare}

A third construction due to Eisenhart lifts the dynamics to a \emph{Lorentzian}
space. On $Q\times\Rb_t\times\Rb_z$ take
\begin{equation}\label{eq:hz-eisenhart}
  G_{\rm E}=M_{ab}\,\ud q^a\ud q^b+2\,\ud t\,\ud z-2V(q)\,\ud t^2 .
\end{equation}
The $t$-equation of the geodesic flow gives $\ddot t=0$, so $t$ is an affine
parameter, and the $q$-equation is $\ddot q^a+\Gamma^a_{bc}\dot q^b\dot q^c=-M^{ab}\partial_bV$ (with
$\Gamma^a_{tt}=M^{ab}\partial_bV$) \cite{hdp:eisenhart1929}. Table~\ref{tab:hz-compare} compares the three.

\begin{table}[t]
\centering
\small
\begin{tabular}{@{}lccc@{}}
\toprule
 & Hertz (hidden rotor) & Jacobi & Eisenhart\\
\midrule
Signature & Riemannian & Riemannian & Lorentzian\\
Dimension & $n+1$ & $n$ & $n+2$\\
Fixed data & hidden momentum $p$ & energy $E$ & none\\
Restriction on $V$ & $V+V_0>0$ & Hill region $V<E$ & none\\
Geodesic parameter & physical time $t$ & $\tau$, $\ud\tau=(E-V)\ud t$ & physical time $t$\\
Where $V$ sits & hidden kinetic energy & conformal factor & $G_{tt}$\\
Hertzian reading & hidden masses & none & none\\
\bottomrule
\end{tabular}
\caption{Three ways to turn the potential force into geometry. Only Hertz's
construction has a mechanical reading in terms of positive masses.}
\label{tab:hz-compare}
\end{table}

% ======================================================================
%  Section insert: Eisenhart's lifts in detail  (+ comparison table)
%  Drop into hertz_chapter.tex in place of the old "Three geometrizations
%  side by side" section (i.e. before \section{Assessment...}).
%
%  Needs: amsmath, amsthm (definition/theorem/proposition/remark/example
%  environments as defined in the chapter), enumitem, booktabs.
%  Cross-references used: thm:hz-cyclic, sec:hz-hidden, sec:hz-jacobi,
%  prop:hz-dp-metric, eq:hz-V, eq:hz-KM, eq:hz-warped.
%  Bibliography keys used: hdp:eisenhart1929, hdp:duval1985,
%  hdp:duval1991, hdp:casetti2000.
% ======================================================================

\section{Eisenhart's lifts in detail}
\label{sec:hz-eisenhart}

In 1928--29 L.~P.~Eisenhart showed that the trajectories of a natural
Lagrangian system can be obtained as geodesics of a metric on a larger space
\cite{hdp:eisenhart1929}. His paper contains \emph{two} constructions,
and they stand in different relations to Hertz. The first is
Lorentzian, works for every potential (including time-dependent ones) and
needs two extra coordinates; the second is Riemannian, needs a positive
potential and one extra coordinate. The second turns out to be
the metric of Theorem~\ref{thm:hz-cyclic}. The first was rediscovered
in 1984--85 by Duval, Burdet, K\"unzle and Perrin, who called the resulting
spaces \emph{Bargmann structures} \cite{hdp:duval1985,hdp:duval1991}, and
the construction is now often called the Eisenhart--Duval lift.

\subsection{The Lorentzian (null) lift}

\begin{definition}[Eisenhart lift]\label{def:hz-eis}
Let $(Q,M)$ be a Riemannian manifold and $V\in C^\infty(Q)$ arbitrary. On
$\widehat Q=Q\times\Rb_t\times\Rb_z$ put
\begin{equation}\label{eq:hz-eis-metric}
  G_{\rm E}=M_{ab}(q)\,\ud q^a\ud q^b+2\,\ud t\,\ud z-2V(q)\,\ud t^2 .
\end{equation}
\end{definition}

The $(t,z)$ block $\bigl(\begin{smallmatrix}-2V&1\\1&0\end{smallmatrix}\bigr)$
has determinant $-1$ whatever $V$ is, so $G_{\rm E}$ is a non-degenerate metric
of Lorentzian signature $(n{+}1,1)$. Its inverse has the components
\begin{equation}\label{eq:hz-eis-inverse}
  G^{ab}=M^{ab},\qquad G^{tz}=1,\qquad G^{zz}=2V,\qquad G^{tt}=G^{ta}=G^{za}=0 .
\end{equation}
The only Christoffel symbols that differ from those of $M$ and from zero are
\begin{equation}\label{eq:hz-eis-christoffel}
  \Gamma^a_{tt}=M^{ab}\partial_bV,\qquad
  \Gamma^z_{at}=\Gamma^z_{ta}=-\partial_aV .
\end{equation}
(For the double pendulum this was checked symbolically for a flat base and
confirmed for the curved base $M$ through the Ricci tensor below.)

\begin{theorem}[Eisenhart]\label{thm:hz-eis}
A curve $\gamma(s)=(q(s),t(s),z(s))$ is a geodesic of $G_{\rm E}$ if and only if
\begin{equation}\label{eq:hz-eis-geod}
  \ddot t=0,\qquad
  \ddot q^a+\Gamma^a_{bc}\dot q^b\dot q^c=-M^{ab}\partial_bV\,\dot t^{\,2},\qquad
  \ddot z=2\,\dot t\,\partial_aV\,\dot q^a .
\end{equation}
Consequently, if $\dot t=1$ (so $s=t$ is physical time):
\begin{enumerate}[label=(\alph*),nosep]
  \item $q(t)$ solves Newton's equations $\ddot q^a+\Gamma^a_{bc}\dot q^b\dot q^c=-M^{ab}\partial_bV$;
  \item $\dot z=2V(q)-\varepsilon$ with a constant $\varepsilon$, so $p_t:=G(\partial_t,\dot\gamma)=-\varepsilon$;
  \item $G_{\rm E}(\dot\gamma,\dot\gamma)=2(E-\varepsilon)$, where $E=\tfrac12M\dot q\dot q+V$ is the energy of $q(t)$;
  \item $\gamma$ is a \emph{null} geodesic exactly when $\varepsilon=E$, and then
        $z(t)=z_0-\int_0^t L\,\ud t'=z_0-S[q]$ with $L=\tfrac12M\dot q\dot q-V$.
\end{enumerate}
Conversely every solution $q(t)$ of Newton's equations lifts, for every
choice of $\varepsilon$, $z_0$ and $t_0$, to a geodesic with $\dot t=1$.
\end{theorem}

\begin{proof}
Equations~\eqref{eq:hz-eis-geod} are the geodesic equations with the symbols
\eqref{eq:hz-eis-christoffel}: $\Gamma^t_{AB}=0$ gives $\ddot t=0$;
$\ddot q^a+\Gamma^a_{bc}\dot q^b\dot q^c+\Gamma^a_{tt}\dot t^2=0$ is the second;
$\ddot z+2\Gamma^z_{at}\dot q^a\dot t=0$ is the third. For $\dot t=1$ the third says
$\ddot z=2\,\ud V/\ud t$, hence $\dot z=2V-\varepsilon$, and $p_t=-2V\dot t+\dot z=-\varepsilon$.
Then
\[
  G(\dot\gamma,\dot\gamma)=M\dot q\dot q+2\dot z-2V
  =2T+2(2V-\varepsilon)-2V=2(T+V-\varepsilon)=2(E-\varepsilon),
\]
which gives (c) and the first half of (d). If $\varepsilon=E$ then
$\dot z=2V-E=V-T=-L$. The converse is immediate by defining
$z$ through $\dot z=2V-\varepsilon$.
\end{proof}

\begin{remark}[The extra coordinate is the action]\label{rem:hz-z-action}
Statement (d) is the geometric content of the lift: along the physical
(null) geodesics the new coordinate $z$ is minus Hamilton's principal
function evaluated along the path. This is no accident. With $\dot t=1$ the
geodesic Lagrangian is $\tfrac12G(\dot\gamma,\dot\gamma)=L+\dot z$, so
$\int\tfrac12G(\dot\gamma,\dot\gamma)\,\ud t=S[q]+z(t_2)-z(t_1)$:
the energy functional of the lift is Hamilton's action up to endpoint terms.
\end{remark}

\subsection{Hamiltonian form and conserved quantities}

The geodesic Hamiltonian $H_{\rm E}=\tfrac12G^{AB}p_Ap_B$ is, by~\eqref{eq:hz-eis-inverse},
\begin{equation}\label{eq:hz-eis-ham}
  H_{\rm E}=\tfrac12M^{ab}p_ap_b+p_tp_z+V(q)\,p_z^{\,2}.
\end{equation}
Both $t$ and $z$ are cyclic for time-independent $V$, so $p_t$ and $p_z$ are conserved.
With $p_z=G(\partial_z,\dot\gamma)=\dot t=1$ and
$H(q,p)=\tfrac12M^{ab}p_ap_b+V$ the physical Hamiltonian, \eqref{eq:hz-eis-ham} becomes
\[
  H_{\rm E}=H(q,p)+p_t ,
\]
and the null condition $H_{\rm E}=0$ is $p_t=-H=-E$: Hamilton's equations of $H_{\rm E}$
are the equations of the \emph{extended phase space}, in which energy is the momentum conjugate
to time and the dynamics is a constraint. The constant $p_z$ plays the role of a mass
(in the Bargmann picture it is the mass or the electric charge of the particle), and the
fixed value $p_z=1$ is what selects Newtonian dynamics with unit
inertia. Hamilton's equations give $\dot t=\partial H_{\rm E}/\partial p_t=p_z=1$ and
$\dot z=\partial H_{\rm E}/\partial p_z=p_t+2V=2V-E$, as found above.

\subsection{Geometry of the lift: a pp-wave}

\begin{proposition}[Curvature of the Eisenhart metric]\label{prop:hz-eis-curv}
\leavevmode
\begin{enumerate}[label=(\roman*),nosep]
  \item $\nabla\,\partial_z=0$: the null vector field $\partial_z$ is covariantly constant.
  \item The Ricci tensor of $G_{\rm E}$ has only the components
        \[
          \operatorname{Ric}_{ab}=\operatorname{Ric}^{M}_{ab},\qquad
          \operatorname{Ric}_{tt}=\Lap_MV .
        \]
\end{enumerate}
\end{proposition}

\begin{proof}
(i) is $\Gamma^A_{Bz}=0$, read off from \eqref{eq:hz-eis-christoffel}. For (ii) I computed the Ricci
tensor symbolically for a flat base with arbitrary $V(x,y)$ and, for the
unit-parameter double pendulum with its curved metric $M$, at random points, where the
only non-vanishing components were $\operatorname{Ric}_{ab}=K_MM_{ab}$ and
$\operatorname{Ric}_{tt}=\Lap_MV$ to rounding error; the general statement is the standard curvature
property of Eisenhart metrics \cite{hdp:casetti2000}.
\end{proof}

A Lorentzian manifold with a covariantly constant null vector whose metric has this
block form is a (generalized) \emph{pp-wave}, and the structure
$(\widehat Q,G_{\rm E},\partial_z)$ is a Bargmann structure \cite{hdp:duval1985,hdp:duval1991}.
Two consequences are worth stating.
\begin{itemize}[nosep]
  \item \emph{Gravity as curvature.} With $\Lap_MV=4\pi G\rho$ for a Newtonian
        potential, $\operatorname{Ric}_{tt}=\Lap_MV$ is the Poisson equation, and Ricci-flatness of the lift
        (flat $M$, harmonic $V$) is Newtonian gravity in vacuum. The lift makes Newtonian
        gravity \emph{look} like general relativity; it is a statement about a Lorentzian
        space of one dimension more than spacetime, not about spacetime.
  \item \emph{Tidal tensor.} The geodesic deviation equation of $G_{\rm E}$ along a lifted
        trajectory, restricted to the $q$-directions, is the usual variational equation of the
        dynamics,
        \[
          \frac{D^2\xi^a}{\ud t^2}+R^{M\,a}{}_{bcd}\,\dot q^b\xi^c\dot q^d
          +M^{ab}\nabla_b\nabla_cV\,\xi^c=0,
        \]
        and the trace of the tidal tensor is $\Lap_MV$ (the origin of $\operatorname{Ric}_{tt}$)
        \cite{hdp:casetti2000}. In contrast to the Jacobi metric of
        Section~\ref{sec:hz-jacobi}, no factor $E-V$ enters and there is no Hill-boundary singularity;
        but the price is that the criterion for exponential instability is the standard
        linear-stability computation, not a purely geometric sign condition.
\end{itemize}

\subsection{Examples}

\begin{example}[Harmonic oscillator]\label{ex:hz-eis-ho}
For $Q=\Rb$, $M=\ud q^2$ and $V=\tfrac12\omega^2q^2$,
$G_{\rm E}=\ud q^2+2\,\ud t\,\ud z-\omega^2q^2\ud t^2$ is a plane wave with $\operatorname{Ric}_{tt}=\omega^2$.
The null geodesics with $\dot t=1$ are $q=A\cos\theta$, $\theta=\omega t+\phi$, and
\[
  z=z_0+\tfrac14A^2\omega\,\sin2\theta ,
\]
which is $z_0-\int L\,\ud t$ with $L=-\tfrac12A^2\omega^2\cos2\theta$. The $z$-motion is a
bounded oscillation at twice the frequency.
\end{example}

\begin{example}[Double pendulum]\label{ex:hz-eis-dp}
With $M$ from Proposition~\ref{prop:hz-dp-metric} and $V$ from~\eqref{eq:hz-V},
\[
\begin{aligned}
  G_{\rm E}={}&(m_1+m_2)\,l_1^2\,\ud\theta_1^2+2m_2l_1l_2\cos(\theta_1-\theta_2)\,\ud\theta_1\ud\theta_2
  +m_2l_2^2\,\ud\theta_2^2\\
  &+2\,\ud t\,\ud z+2g\bigl[(m_1+m_2)l_1\cos\theta_1+m_2l_2\cos\theta_2\bigr]\ud t^2 .
\end{aligned}
\]
This is a four-dimensional Lorentzian manifold $\Tor^2\times\Rb^2$ with
$\operatorname{Ric}_{ab}=K_MM_{ab}$ (with $K_M$ from~\eqref{eq:hz-KM}) and $\operatorname{Ric}_{tt}=\Lap_MV$.
Its Killing fields include $\partial_z$ (null, parallel) and $\partial_t$; the rotation
$\partial_{\theta_1}+\partial_{\theta_2}$ that was a symmetry of the free pendulum is
broken by $V$. The null geodesics with $\dot t=1$ project to the chaotic trajectories of the pendulum,
and $z$ along them is $-S[\theta]$.
\end{example}

\subsection{The Riemannian lift, and its identity with the hidden rotor}

Eisenhart's second construction, for \emph{time-independent potentials that are positive} on $Q$,
uses one extra coordinate and a Riemannian metric of the form
$M+A(q)\,\ud z^2$ with $A$ built from $V$
\cite{hdp:eisenhart1929,hdp:casetti2000}. Taking $A=1/V_c$ with $V_c=V+c>0$,
\begin{equation}\label{eq:hz-eis-riem}
  G_{\rm R}=M_{ab}\,\ud q^a\ud q^b+\frac{1}{V_c(q)}\,\ud z^2 ,
\end{equation}
the Hamiltonian is $\tfrac12M^{ab}p_ap_b+\tfrac12V_c\,p_z^2$, which equals
the physical Hamiltonian $\tfrac12M^{ab}p_ap_b+V_c$ for $p_z=\sqrt2$. This is exactly
the metric~\eqref{eq:hz-warped} of Theorem~\ref{thm:hz-cyclic} with $f^2=p^2/(2V_c)$, $p=\sqrt2$,
and $\varphi=z$. The theorem therefore \emph{is} Eisenhart's Riemannian lift, with a different proof and with
a physical interpretation of the extra coordinate added: Eisenhart's $z$ is Hertz's hidden
cyclic coordinate, and $V=\tfrac12f^2\dot z^2$ is the kinetic energy of the hidden rotor.
(Hertz did not write down the warped metric; the idea that a potential is the kinetic energy of
a cyclic hidden motion is his, and the exact geometric statement is Eisenhart's.)

\subsection{Hertzian or not?}

Does the Lorentzian lift satisfy Hertz's requirements? Only in part. Its dynamics is
free in the sense that the law is geodesic, and no force appears. But Hertz's mass metric is positive
definite: kinetic energy is $\tfrac12\sum m\dot x^2\ge0$, and the straightest path is defined by minimizing
a norm of acceleration. In signature $(n{+}1,1)$ there is no such norm, the physical
trajectories are \emph{null}, and the extra coordinate $z$ is a null direction, not a hidden mass.
The Lorentzian lift is thus forceless but not Hertzian; the Riemannian lift is both.
That the two lifts are available side by side underlines the point of
Section~\ref{sec:hz-hidden}(c): the geometric content of the visible dynamics does not fix its higher-dimensional
realization.

\section{Four geometrizations side by side}
\label{sec:hz-compare}

\begin{table}[t]
\centering
\small
\setlength{\tabcolsep}{4pt}
\begin{tabular}{@{}lccc@{}}
\toprule
 & Hertz / Eisenhart--Riem. & Jacobi & Eisenhart--Lorentz\\
\midrule
Signature & Riemannian & Riemannian & Lorentzian $(n{+}1,1)$\\
Dimension & $n+1$ & $n$ & $n+2$\\
Fixed data & $p_z$ (hidden momentum) & energy $E$ & $p_z=1$, null condition\\
Restriction on $V$ & $V+c>0$ & Hill region $V<E$ & none\\
Geodesic parameter & physical time $t$ & $\tau$, $\ud\tau=(E-V)\ud t$ & physical time $t$\\
Where $V$ sits & hidden kinetic energy & conformal factor & $G_{tt}$\\
Role of extra coordinates & hidden angle $z$ & none & $t$ and null $z=-S$\\
Time-dependent $V$ & no & no & yes\\
Hertzian reading & hidden masses & none & none (null)\\
\bottomrule
\end{tabular}
\caption{Three ways (four constructions, since the first column is Hertz's idea and Eisenhart's
Riemannian metric at once) to turn the potential force into geometry. Only the first has a
mechanical reading in terms of positive masses.}
\label{tab:hz-compare}
\end{table}

Table~\ref{tab:hz-compare} summarizes the comparison. The Lorentzian lift is the only one that
treats time-dependent potentials uniformly, since $t$ is a coordinate of the lifted
space; the Jacobi construction requires conservation of energy and the Riemannian lift
requires a time-independent positive $V$.

\section{Assessment: what Hertz's program delivers}
\label{sec:hz-assess}

\paragraph{What it gains.}
Applied to the double pendulum the program does what it promises, in one
respect: it removes every force from the statement of the law. The tension in
the rods is a reaction of the rigid connection, the centrifugal-looking terms
are Christoffel symbols, and gravity is the kinetic energy of a rotor. The
geometry also yields something new that is invisible in the force picture:
the Gauss--Bonnet argument of Proposition~\ref{prop:hz-gb} explains, without any
numerics, why the Jacobi geometry of the double pendulum above the upright
energy must contain regions in which nearby motions diverge.

\paragraph{What it costs.}
The cost is the hidden structure, and the analysis of Section~\ref{sec:hz-hidden}
shows that it is both inessential and unconstrained. The visible motion
of the pendulum determines the ordinary differential equation~\eqref{eq:hz-dp-grav}
and nothing else; the hidden rotor, the Jacobi metric and the Eisenhart
lift are three unrelated ways of writing that same equation as a geodesic problem,
each in a different dimension and with different auxiliary data. To the extent
that observation of the pendulum is all there is, the hidden masses
are idle wheels. This is the pattern argued throughout this book: the equations
of motion are the primary, real-world content of mechanics, and variational and
geometric formulations are ideal-world devices for finding and organizing them.
The three geometrizations are in this sense charts of one object, each valid
on its own domain and each introducing singularities of its own: the positivity
requirement $V+V_0>0$, the Hill boundary, the Lorentzian signature.

\paragraph{What it cannot reach.}
Hertz's systems are closed and conservative: all forces, including friction,
must be traced to hidden conservative motion. Dissipation therefore enters
only as a statement that the hidden energy is out of sight, and genuine open
systems, such as the damped double pendulum, are outside the scheme as stated;
they are the subject of the chapters on damping and on open systems.
The nonholonomic extension, in which straightest paths differ from shortest
paths, is the other place where Hertz's mechanics is not simply geodesic flow
(Section~\ref{sec:hz-geometry}).

\paragraph{Historical coda.}
The reception was divided between those who found the hidden-mass hypothesis a
barrier to physical content and those who appreciated the geometric form
\cite{hdp:lutzen2005}. What survived was less the philosophy than the geometry:
the representation of constrained mechanics as geodesic flow on a Riemannian
manifold, which is the standard language of the modern geometric treatments.

\section*{Problems}
\addcontentsline{toc}{section}{Problems}

\begin{enumerate}[label=\textbf{\thechapter.\arabic*}.]
  \item Prove Proposition~\ref{prop:hz-law}(i) by Lagrange multipliers: minimize
        $g(a,a)$ subject to $\omega_\alpha(a)=c_\alpha$ and show the multipliers
        are the constraint reactions.
  \item Derive~\eqref{eq:hz-dp-free} directly from Proposition~\ref{prop:hz-law}(iii) with
        the constraints $|\mathbf x_1|=l_1$, $|\mathbf x_2-\mathbf x_1|=l_2$, eliminating the
        multipliers (the rod tensions).
  \item Identify the second integral of the free double pendulum with
        $J$ of Proposition~\ref{prop:hz-dp-curvature}, and find the dependence of the
        period of the $\delta$-motion on $J$ and the speed.
  \item In Example~\ref{ex:hz-simple} show that the surface of revolution is
        smooth, that $K=-f_{ss}/f$ does not depend on $p$, and determine where $K<0$
        (compare Gauss--Bonnet for the torus).
  \item Verify the claim of Section~\ref{sec:hz-hidden}(a) that $\ud r\otimes\ud r=O(V_0^{-3})$
        at fixed $p$, and compute the leading correction to the pendulum's small-oscillation frequency.
  \item Verify that the geodesics of~\eqref{eq:hz-eisenhart} with $\dot t=1$ project to
        solutions of~\eqref{eq:hz-dp-grav}, and find the quantity conserved by the
        null Killing vector $\partial_z$.
  \item Show that for $g\neq0$ the vector field $\partial_{\theta_1}+\partial_{\theta_2}$ is
        \emph{not} a Killing field of $J_E$, so the angular momentum $J$ of
        Proposition~\ref{prop:hz-dp-curvature} is no longer conserved. Why does this loss of
        symmetry not by itself prove that the flow is non-integrable?
\end{enumerate}

\section*{Notes and references}
\addcontentsline{toc}{section}{Notes and references}

\emph{On verification.} The kinetic metric, the geodesic-versus-Lagrange form of the equations of
motion, the explicit accelerations~\eqref{eq:hz-dp-free}--\eqref{eq:hz-dp-grav}, the Christoffel symbols
of the warped metric, the curvature~\eqref{eq:hz-KM} and the conformal
formula~\eqref{eq:hz-KJ} were checked symbolically (SymPy); the curvature map and the table were computed
numerically from~\eqref{eq:hz-KJ} on a $900\times900$ grid, and the Gauss--Bonnet integral at $E=4$
vanishes to rounding error, as it should.

\begingroup
\let\chapter\section
\renewcommand{\bibname}{References}
\begin{thebibliography}{99}
\bibitem{hdp:hertz1894} H.~Hertz, \emph{Die Prinzipien der Mechanik in neuem Zusammenhange dargestellt},
  Gesammelte Werke, Vol.~III (J.~A.~Barth, Leipzig, 1894), with a preface by H.~von Helmholtz.
  English translation: \emph{The Principles of Mechanics Presented in a New Form}, trans.\ D.~E.~Jones and
  J.~T.~Walley (Macmillan, London, 1899; repr.\ Dover, New York, 1956).
\bibitem{hdp:lutzen2005} J.~L\"utzen, \emph{Mechanistic Images in Geometric Form: Heinrich Hertz's
  ``Principles of Mechanics''} (Oxford University Press, Oxford, 2005).
\bibitem{hdp:gauss1829} C.~F.~Gauss, \"Uber ein neues allgemeines Grundgesetz der Mechanik,
  \emph{J.\ reine angew.\ Math.}\ \textbf{4} (1829) 232--235.
\bibitem{hdp:eom-hertz} V.~V.~Rumyantsev, Hertz principle, in \emph{Encyclopedia of Mathematics} (EMS Press).
\bibitem{hdp:routh1877} E.~J.~Routh, \emph{A Treatise on the Stability of a Given State of Motion}
  (Macmillan, London, 1877).
\bibitem{hdp:helmholtz1884} H.~von Helmholtz, Principien der Statik monocyklischer Systeme,
  \emph{J.\ reine angew.\ Math.}\ \textbf{97} (1884) 111--140, 317--336.
\bibitem{hdp:nash1956} J.~Nash, The imbedding problem for Riemannian manifolds,
  \emph{Ann.\ of Math.}\ \textbf{63} (1956) 20--63.
\bibitem{hdp:jacobi1866} C.~G.~J.~Jacobi, \emph{Vorlesungen \"uber Dynamik}, ed.\ A.~Clebsch
  (G.~Reimer, Berlin, 1866).
\bibitem{hdp:arnold1989} V.~I.~Arnold, \emph{Mathematical Methods of Classical Mechanics},
  2nd ed., Graduate Texts in Mathematics \textbf{60} (Springer, New York, 1989).
\bibitem{hdp:docarmo1992} M.~P.~do Carmo, \emph{Riemannian Geometry} (Birkh\"auser, Boston, 1992).
\bibitem{hdp:eisenhart1929} L.~P.~Eisenhart, Dynamical trajectories and geodesics,
  \emph{Ann.\ of Math.}\ \textbf{30} (1929) 591--606.
\bibitem{hdp:casetti1996} L.~Casetti, C.~Clementi and M.~Pettini, Riemannian theory of Hamiltonian
  chaos and Lyapunov exponents, \emph{Phys.\ Rev.\ E} \textbf{54} (1996) 5969--5984.
\bibitem{hdp:pettini2007} M.~Pettini, \emph{Geometry and Topology in Hamiltonian Dynamics and
  Statistical Mechanics} (Springer, New York, 2007).
\end{thebibliography}
\endgroup
 from a `book' preamble that loads
%    amsmath, amssymb, amsthm, mathtools, enumitem, booktabs, graphicx,
%    tikz (+ arrows.meta library), hyperref.
%  Figure file needed in the same directory:  jacobi_curvature.pdf
%  Theorem-type environments are defined below only if not yet defined.
% ======================================================================

\makeatletter
\@ifundefined{theorem}{\theoremstyle{plain}\newtheorem{theorem}{Theorem}[chapter]}{}
\@ifundefined{proposition}{\theoremstyle{plain}\newtheorem{proposition}[theorem]{Proposition}}{}
\@ifundefined{corollary}{\theoremstyle{plain}\newtheorem{corollary}[theorem]{Corollary}}{}
\@ifundefined{definition}{\theoremstyle{definition}\newtheorem{definition}[theorem]{Definition}}{}
\@ifundefined{example}{\theoremstyle{definition}\newtheorem{example}[theorem]{Example}}{}
\@ifundefined{remark}{\theoremstyle{remark}\newtheorem{remark}[theorem]{Remark}}{}
\makeatother

\providecommand{\Rb}{\mathbb{R}}
\providecommand{\Tor}{\mathbb{T}}
\providecommand{\Sph}{\mathbb{S}}
\providecommand{\ud}{\mathrm{d}}
\providecommand{\Lap}{\Delta}
\setlength{\emergencystretch}{2.5em}

\chapter{Hertz's Forceless Mechanics and the Double Pendulum}
\label{ch:hertz-double-pendulum}

\section{A mechanics without forces}
\label{sec:hz-intro}

Newton's \emph{Principia} takes force as a primitive. The concept is so
deeply built into the textbook presentation of mechanics that one forgets
how uneasy its status has always been: a force is neither a displacement
nor a mass, it cannot be seen, and the classical examples---gravitation
across empty space, the tension of a string, the reaction of a table---have
little in common except that each can be written on the right-hand side of
$m\ddot{\mathbf x}=\mathbf F$. In \emph{Die Prinzipien der Mechanik in neuem
Zusammenhange dargestellt}, which appeared in 1894 after Hertz's death
and with a preface by Helmholtz, Heinrich Hertz proposed to remove force
from the foundations altogether \cite{hdp:hertz1894}. His own ``image''
of mechanics rests on three notions only---time, space and mass---and on
a single law, which in the Jones--Walley translation of \S309 reads:
``Every free system persists in its state of rest or of uniform motion in
a straightest path.''

The sentence has two moving parts. A \emph{free} system is one subject to
no forces whatever; it is subject only to rigid connections between the
positions of its parts. A \emph{straightest path} is a purely geometric
notion, the path of least curvature compatible with those connections. What
the Newtonian calls a force is then, in Hertz's account, of one of two
kinds. It is either the reaction of a rigid connection among the visible
bodies, or it is the visible trace of \emph{hidden masses}, rigidly
connected to the visible ones and moving in ways we do not observe
\cite{hdp:hertz1894,hdp:lutzen2005}. Helmholtz's preface to the book already
voices reservations on whether hidden systems of the required kind can
actually be found for the forces of nature, and the reception of the
book was, for that reason, largely respectful and largely unconvinced.

This chapter takes Hertz's program seriously as \emph{mathematics} and
asks what it does when applied to the smallest system in which all its
ingredients meet: the planar double pendulum. The double pendulum has
rigid connections, which are Hertz's native material (the two rods); it has
a force, gravity, which Hertz must explain away; and it has chaos, which
any mechanics must be able to represent. The plan is as follows.

\begin{itemize}[nosep]
  \item Section~\ref{sec:hz-geometry} states Hertz's law precisely and shows
        that, for holonomic connections, it is the statement that motion is
        a geodesic of the mass-weighted configuration space.
  \item Section~\ref{sec:hz-free-dp} shows that the double pendulum
        \emph{without} gravity is a free Hertz system, computes the
        curvature of its configuration space, and finds a hidden Killing
        field.
  \item Section~\ref{sec:hz-hidden} switches gravity on and eliminates it in
        Hertz's way: by a representation theorem that turns any positive
        potential into the kinetic energy of one hidden cyclic coordinate.
  \item Section~\ref{sec:hz-jacobi} gives the second geometrization of the
        same dynamics, the Jacobi metric at fixed energy, and uses it to ask
        where the chaos of the double pendulum sits in the geometry.
  \item Sections~\ref{sec:hz-compare} and~\ref{sec:hz-assess} compare the
        geometrizations and assess what the Hertzian program does and does
        not deliver, in the light of the theme of this book: the equations of
        motion, not the formulations built around them, are the real-world
        content of mechanics.
\end{itemize}

\section{Hertz's law and the geometry of configuration space}
\label{sec:hz-geometry}

\subsection{Systems, connections, free motion}

\begin{definition}[Hertzian material system]\label{def:hz-system}
A \emph{material system} consists of $N$ particles with masses $m_i>0$ and
positions $\mathbf x_i\in\Rb^3$. Its \emph{mass-weighted ambient space} is
$(\Rb^{3N},g)$ with
\[
  g(\xi,\eta)\;=\;\sum_{i=1}^{N}m_i\,\xi_i\!\cdot\eta_i .
\]
A \emph{rigid connection} is a relation that is linear and homogeneous in
the differentials of the coordinates, $\omega_\alpha(x)(\ud x)=0$,
$\alpha=1,\dots,k$, where the $\omega_\alpha$ are Pfaffian forms depending
on position. Write
$D_x=\bigcap_\alpha\ker\omega_\alpha(x)\subset\Rb^{3N}$ for the space of
\emph{admissible velocities}. The connections are \emph{holonomic} if $D$
is integrable, i.e.\ if there is a submanifold $\mathcal N\subset\Rb^{3N}$
through every admissible position with $D=T\mathcal N$; otherwise they are
\emph{nonholonomic} (both terms are Hertz's). A system is \emph{free} if
it is subject to rigid connections only, with no active forces.
\end{definition}

Hertz measures path elements by the same mass-weighted expression
$g(\ud x,\ud x)=\sum m_i\,\ud x_i^2$ (he divides by the total mass, which is
immaterial here). Differentiating an admissible motion $x(t)$, with
$\dot x\in D_x$, along the connections gives, for the acceleration $a=\ddot x$,
\[
  \omega_\alpha(x)(a)\;=\;c_\alpha(x,\dot x),
  \qquad
  c_\alpha:=-\,\partial_\beta\omega_{\alpha,\gamma}\,\dot x^\beta\dot x^\gamma ,
\]
so that at a given state $(x,\dot x)$ the admissible accelerations form an
affine space $A_{(x,\dot x)}=a_0+D_x$.

\begin{definition}[Straightest path]\label{def:hz-straightest}
An admissible path is \emph{straightest} if at every point its curvature
is the least among all admissible paths having the same position and the
same tangent. For a path parametrized by constant speed, curvature is
proportional to $|a|_g$, so this means
\[
  a\;=\;\operatorname*{arg\,min}_{a'\in A_{(x,\dot x)}}\ g(a',a').
\]
\end{definition}

\begin{proposition}[Hertz's law in Newtonian language]\label{prop:hz-law}
Let a free system move along a straightest path. Then:
\begin{enumerate}[label=(\roman*),nosep]
  \item $g(a,w)=0$ for every $w\in D_x$, i.e.\ $a\in D_x^{\perp}$;
  \item the speed $g(\dot x,\dot x)^{1/2}$ is constant;
  \item in components, $m_i\ddot{\mathbf x}_i=\mathbf R_i$ with a reaction
        $R=(\mathbf R_i)$ satisfying $g^{-1}\!R\perp D_x$, that is, the
        equations of d'Alembert--Lagrange with zero active forces;
  \item if the connections are holonomic, $D=T\mathcal N$, then $x(t)$ is a
        geodesic of the induced Riemannian manifold $(\mathcal N,g|_{\mathcal N})$.
\end{enumerate}
\end{proposition}

\begin{proof}
(i) Minimizing $g(a',a')$ over the affine space $a_0+D_x$ yields the
orthogonal projection of the origin onto it, which is orthogonal to the
direction space $D_x$. (ii) $\frac{\ud}{\ud t}g(\dot x,\dot x)=2g(a,\dot x)=0$
because $\dot x\in D_x$. (iii) Write the reaction as $R:=g(a,\cdot)$; then
$R(w)=0$ on $D_x$ is (i). (iv) Let $\nabla$ be the flat connection of
$\Rb^{3N}$. For a curve in $\mathcal N$ the Gauss formula splits $a=\nabla_{\dot x}\dot x$
into a part tangent to $\mathcal N$, which equals the covariant acceleration
$\nabla^{\mathcal N}_{\dot x}\dot x$ of the induced metric, and a normal part.
By (i) the tangential part vanishes.
\end{proof}

\begin{remark}[Gauss, and what is new]
Gauss's principle of least constraint minimizes
\[
  \sum_i m_i\Bigl|\mathbf a_i-\frac{\mathbf F_i}{m_i}\Bigr|^2
\]
over admissible accelerations \cite{hdp:gauss1829}; Hertz's law is the case
$\mathbf F=0$, and for stationary connections without active forces the two
coincide \cite{hdp:eom-hertz}. The mathematical gain is therefore nil; the
gain is conceptual. In Gauss's principle forces are given data and
constraints are restrictions; in Hertz's the forces are to be \emph{derived},
and the whole of dynamics becomes a statement about the metric $g$ and the
distribution $D$.
\end{remark}

\begin{remark}[Nonholonomic connections]
For nonholonomic $D$ a straightest path is in general neither a geodesic of
any induced metric nor a minimizer of length among admissible paths; Hertz
stressed that Hamilton's principle and the principle of least action fail for
such systems \cite{hdp:lutzen2005}. The straightest-path condition (i) is the
Lagrange--d'Alembert equation, not the variational (``vakonomic'') one. The
double pendulum has holonomic connections only, so none of this affects the
rest of the chapter.
\end{remark}

\subsection{Holonomic Hertzian mechanics is geodesic flow}

For holonomic connections, part (iv) says that the entire dynamics of a free
system is the geodesic flow of $(\mathcal N,M)$, where in coordinates
$q^a$ on $\mathcal N$ the \emph{kinetic metric} is
\begin{equation}\label{eq:hz-kinetic-metric}
  M_{ab}(q)\;=\;\sum_i m_i\,\frac{\partial\mathbf x_i}{\partial q^a}\!\cdot\frac{\partial\mathbf x_i}{\partial q^b},
\end{equation}
and the equations of motion are
\begin{equation}\label{eq:hz-geodesic}
  \ddot q^a+\Gamma^a_{bc}(q)\,\dot q^b\dot q^c=0 ,
  \qquad
  \Gamma^a_{bc}=\tfrac12M^{ad}\bigl(\partial_bM_{dc}+\partial_cM_{db}-\partial_dM_{bc}\bigr).
\end{equation}
The ``centrifugal'' and ``Coriolis'' terms of textbook mechanics are the
Christoffel symbols $\Gamma$: they are not forces but the coordinate
expression of the connection of $M$. What remains outside this scheme, and what
Hertz must account for in order to keep his claim of generality, is every
term that is \emph{not} of this form, notably a potential force on the
right-hand side of~\eqref{eq:hz-geodesic}.

\section{The double pendulum without gravity is a free Hertz system}
\label{sec:hz-free-dp}

\begin{figure}[t]
\centering
\begin{tikzpicture}[scale=1.15,>=Stealth]
  % ceiling
  \draw[thick] (-1.0,0) -- (1.0,0);
  \foreach \x in {-0.9,-0.6,...,0.9}{\draw (\x,0) -- (\x-0.15,0.18);}
  \coordinate (O) at (0,0);
  \coordinate (A) at ({2*sin(28)},{-2*cos(28)});
  \coordinate (B) at ({2*sin(28)+1.6*sin(62)},{-2*cos(28)-1.6*cos(62)});
  % verticals
  \draw[dashed,gray] (O) -- (0,-3.2);
  \draw[dashed,gray] (A) -- ++(0,-1.9);
  % rods
  \draw[very thick] (O) -- (A) node[midway,above right=-2pt] {$l_1$};
  \draw[very thick] (A) -- (B) node[midway,above right=-1pt] {$l_2$};
  % bobs
  \fill (O) circle (1.6pt);
  \filldraw[fill=white,thick] (A) circle (0.13) node[right=5pt] {$m_1$};
  \filldraw[fill=white,thick] (B) circle (0.13) node[right=5pt] {$m_2$};
  % angles
  \draw[->] (0,-0.85) arc[start angle=-90,end angle=-62,radius=0.85];
  \node at (0.22,-1.12) {\small$\theta_1$};
  \draw[->] (A)++(0,-0.7) arc[start angle=-90,end angle=-28,radius=0.7];
  \node at ({2*sin(28)+0.42},{-2*cos(28)-0.98}) {\small$\theta_2$};
  % gravity
  \draw[->,thick] (3.2,-0.4) -- (3.2,-1.4) node[midway,right] {$g$};
\end{tikzpicture}
\caption{The planar double pendulum. Both angles are measured from the
downward vertical, so the configuration manifold is the torus
$\Tor^2=\Sph^1\times\Sph^1$ (the rods may swing through the upright position).}
\label{fig:hz-dp}
\end{figure}

Let the two bobs have masses $m_1,m_2$ and positions $\mathbf x_1,\mathbf x_2$
in a vertical plane, joined by massless rigid rods of lengths $l_1,l_2$ to the
pivot and to each other (Figure~\ref{fig:hz-dp}). The mass-weighted ambient
space is $\Rb^4$ with $g=\operatorname{diag}(m_1,m_1,m_2,m_2)$, and the
connections are the two holonomic relations
$|\mathbf x_1|=l_1$, $|\mathbf x_2-\mathbf x_1|=l_2$. They define
$\mathcal N=\Tor^2$, parametrized by
\[
  \mathbf x_1=l_1(\sin\theta_1,-\cos\theta_1),\qquad
  \mathbf x_2=\mathbf x_1+l_2(\sin\theta_2,-\cos\theta_2).
\]

\begin{proposition}[Kinetic metric]\label{prop:hz-dp-metric}
With $\delta:=\theta_1-\theta_2$, the induced metric~\eqref{eq:hz-kinetic-metric} is
\[
  M\;=\;\begin{pmatrix}(m_1+m_2)\,l_1^2 & m_2\,l_1l_2\cos\delta\\[2pt]
                       m_2\,l_1l_2\cos\delta & m_2\,l_2^2\end{pmatrix},
  \qquad
  \det M=m_2\,l_1^2l_2^2\,\bigl(m_1+m_2\sin^2\delta\bigr).
\]
\end{proposition}

\begin{proof}
Differentiate the parametrization: $\partial_{\theta_1}\mathbf x_1=l_1(\cos\theta_1,\sin\theta_1)$,
$\partial_{\theta_1}\mathbf x_2=\partial_{\theta_1}\mathbf x_1$,
$\partial_{\theta_2}\mathbf x_1=0$, $\partial_{\theta_2}\mathbf x_2=l_2(\cos\theta_2,\sin\theta_2)$.
Then $M_{11}=(m_1+m_2)l_1^2$, $M_{22}=m_2l_2^2$ and
$M_{12}=m_2l_1l_2(\cos\theta_1\cos\theta_2+\sin\theta_1\sin\theta_2)=m_2l_1l_2\cos\delta$.
\end{proof}

The metric depends on $\theta_1,\theta_2$ only through $\delta$. Hertz's
law for the gravity-free pendulum is therefore the geodesic
equation~\eqref{eq:hz-geodesic} of $(\Tor^2,M)$, which reads explicitly
\begin{equation}\label{eq:hz-dp-free}
\begin{aligned}
  \ddot\theta_1&=-\frac{m_2\,l_1\dot\theta_1^2\sin\delta\cos\delta+m_2\,l_2\dot\theta_2^2\sin\delta}{l_1\,(m_1+m_2\sin^2\delta)},\\[2pt]
  \ddot\theta_2&=\frac{(m_1+m_2)\,l_1\dot\theta_1^2\sin\delta+m_2\,l_2\dot\theta_2^2\sin\delta\cos\delta}{l_2\,(m_1+m_2\sin^2\delta)} .
\end{aligned}
\end{equation}

\begin{proposition}[Curvature and a hidden symmetry]\label{prop:hz-dp-curvature}
The Gaussian curvature of $(\Tor^2,M)$ is
\begin{equation}\label{eq:hz-KM}
  K_M\;=\;\frac{m_1\cos\delta}{l_1l_2\,\bigl(m_1+m_2\sin^2\delta\bigr)^{2}} .
\end{equation}
The vector field $\partial_{\theta_1}+\partial_{\theta_2}$ is Killing, with Noether
integral
\[
  J=\bigl[(m_1+m_2)l_1^2+m_2l_1l_2\cos\delta\bigr]\dot\theta_1
   +\bigl[m_2l_2^2+m_2l_1l_2\cos\delta\bigr]\dot\theta_2 ,
\]
the total angular momentum about the pivot. Together with the
speed, $J$ makes the geodesic flow of $(\Tor^2,M)$ Liouville integrable.
\end{proposition}

\begin{proof}
Formula~\eqref{eq:hz-KM} follows from $K=R_{1212}/\det M$ with the Christoffel
symbols of Proposition~\ref{prop:hz-dp-metric} (the computation was checked
symbolically). Since $M$ depends on $\delta=\theta_1-\theta_2$ only, it is
invariant under the simultaneous rotation $\theta_i\mapsto\theta_i+s$, whose
generator is $\partial_{\theta_1}+\partial_{\theta_2}$; the conserved quantity
$M(\partial_{\theta_1}+\partial_{\theta_2},\dot\theta)$ is $J$. A $2$-degree-of-freedom
system with two independent commuting integrals (here the speed and $J$) is
integrable.
\end{proof}

\begin{remark}[Negative curvature without chaos]\label{rem:hz-gb-free}
By~\eqref{eq:hz-KM}, $K_M>0$ for $|\delta|<\pi/2$ and $K_M<0$ for $|\delta|>\pi/2$.
This is forced: the Gauss--Bonnet theorem for the torus gives
$\int_{\Tor^2}K_M\,\ud A=2\pi\chi(\Tor^2)=0$, so a non-flat metric on $\Tor^2$ must
have regions of both signs (check directly: $\ud A=\sqrt{\det M}\,\ud\theta_1\ud\theta_2$
and $K_M\sqrt{\det M}\propto\cos\delta\,(m_1+m_2\sin^2\delta)^{-3/2}$ integrates to zero
over $\delta\in[0,2\pi]$). Yet the geodesic flow is integrable. Negative curvature somewhere is thus
\emph{not} an indicator of chaos; we shall need this caveat in
Section~\ref{sec:hz-jacobi}.
\end{remark}

\section{Switching on gravity: hidden masses}
\label{sec:hz-hidden}

With gravity the double pendulum is the Lagrangian system with
$L=\tfrac12M_{ab}\dot\theta^a\dot\theta^b-V$,
\begin{equation}\label{eq:hz-V}
  V(\theta_1,\theta_2)=-(m_1+m_2)\,g\,l_1\cos\theta_1-m_2\,g\,l_2\cos\theta_2 ,
\end{equation}
and its equations of motion are
\begin{equation}\label{eq:hz-dp-grav}
  \ddot\theta^a+\Gamma^a_{bc}\dot\theta^b\dot\theta^c=-M^{ab}\partial_bV ,
\end{equation}
that is, \eqref{eq:hz-dp-free} with the additional terms
\begin{align*}
  \ddot\theta_1&\ni\frac{m_2\,g\sin\theta_2\cos\delta-(m_1+m_2)\,g\sin\theta_1}{l_1\,(m_1+m_2\sin^2\delta)},\\
  \ddot\theta_2&\ni\frac{(m_1+m_2)\,g\,(\sin\theta_1\cos\delta-\sin\theta_2)}{l_2\,(m_1+m_2\sin^2\delta)} .
\end{align*}
The right-hand side of~\eqref{eq:hz-dp-grav} is a force. For a Hertzian it
must be an appearance.

\subsection{Hertz's response and a representation theorem}

Hertz's answer is that forces are the effect of motions we do not see. The
visible system is coupled by rigid connections to a hidden system; the
combined system is free; and the visible variables, which alone are
observed, move as if acted on by forces \cite{hdp:hertz1894,hdp:lutzen2005}.
In the language of later analytical mechanics the mechanism is familiar: a
hidden coordinate that is \emph{cyclic} can be eliminated by the Routh
procedure \cite{hdp:routh1877}, leaving a visible system with an extra
potential term (Helmholtz's monocyclic systems are the thermodynamic
instance \cite{hdp:helmholtz1884}). The following theorem shows that this
device is not merely available but \emph{universal} for potentials.

\begin{theorem}[Cyclic representation of a potential]\label{thm:hz-cyclic}
Let $(Q,M)$ be a Riemannian manifold, $V\in C^\infty(Q)$ with $V>0$, and
$p>0$ a constant. On $Q\times\Sph^1$ with coordinates $(q^a,\varphi)$ put
\begin{equation}\label{eq:hz-warped}
  G=M_{ab}(q)\,\ud q^a\ud q^b+f(q)^2\,\ud\varphi^2,\qquad f=\frac{p}{\sqrt{2V}} .
\end{equation}
\begin{enumerate}[label=(\alph*),nosep]
  \item If $(q(t),\varphi(t))$ is a geodesic of $G$ with $f^2\dot\varphi=p$, then
        $q(t)$ solves $\ddot q^a+\Gamma^a_{bc}\dot q^b\dot q^c=-M^{ab}\partial_bV$ and
        $\tfrac12M_{ab}\dot q^a\dot q^b+V=\tfrac12G(\dot\gamma,\dot\gamma)$ is constant.
  \item Conversely, if $q(t)$ solves the equation in (a), then
        $\varphi(t)=\varphi_0+\int_0^t 2V(q(s))/p\;\ud s$ makes
        $(q,\varphi)$ a geodesic of $G$ with $f^2\dot\varphi=p$.
\end{enumerate}
\end{theorem}

\begin{proof}
Since $G$ is block diagonal and $\varphi$ does not appear in it, the only
Christoffel symbols of $G$ that differ from those of $M$ and from zero are
\[
  \Gamma^a_{\varphi\varphi}=-\tfrac12M^{ab}\partial_b(f^2),\qquad
  \Gamma^\varphi_{a\varphi}=\Gamma^\varphi_{\varphi a}=\partial_af/f ,
\]
while $\Gamma^a_{b\varphi}=\Gamma^\varphi_{ab}=\Gamma^\varphi_{\varphi\varphi}=0$.
The geodesic equations are therefore
\[
  \ddot q^a+\Gamma^a_{bc}\dot q^b\dot q^c-\tfrac12M^{ab}\partial_b(f^2)\,\dot\varphi^2=0,
  \qquad
  \ddot\varphi+2\,\frac{\partial_af}{f}\,\dot q^a\dot\varphi=0 .
\]
The second equation is $\frac{\ud}{\ud t}(f^2\dot\varphi)=0$, the conservation of
the cyclic momentum. Insert $\dot\varphi=p/f^2$ in the first:
\[
  \tfrac12M^{ab}\partial_b(f^2)\,\frac{p^2}{f^4}
  =-\tfrac12M^{ab}\partial_b\!\Bigl(\frac{p^2}{f^2}\Bigr)
  =-M^{ab}\partial_b V ,
\]
because $p^2/(2f^2)=V$. This proves (a); the energy statement follows from
$\tfrac12f^2\dot\varphi^2=p^2/(2f^2)=V$. Part (b) reverses the steps: define
$\varphi$ by $\dot\varphi=p/f^2=2V/p$, so the second geodesic equation holds, and
the first reduces to the given equation.
\end{proof}

The theorem is Routh reduction run backwards: $R=\tfrac12M\dot q\dot q-p^2/(2f^2)$
is the Routhian of $G$, and $p^2/(2f^2)$ is its effective potential.
(A metric of the form $g+f^2\ud\varphi^2$ is also what Kaluza--Klein
theory calls a circle compactification with a scalar $f$ and no gauge field;
the potential $p^2/(2f^2)$ is the usual momentum-induced one.)
The physical reading, in Hertz's terms, is that \emph{potential energy is the
kinetic energy of the hidden rotor}: $V=\tfrac12f^2\dot\varphi^2$.

\begin{example}[The simple pendulum as a Clairaut problem]\label{ex:hz-simple}
For a pendulum of mass $m$ and length $l$, $Q=\Sph^1$ with $M=ml^2\ud\theta^2$ and
$V=mgl(1-\cos\theta)+\varepsilon$, $\varepsilon>0$ (a constant shift changes only
the value of the conserved energy). Then \eqref{eq:hz-warped} is a metric on a
torus of revolution with profile $f(\theta)=p/\sqrt{2V(\theta)}$,
and the pendulum's motion is a geodesic of that surface. The cyclic
momentum $p=f^2\dot\varphi$ is Clairaut's constant: it equals $f$ times the
component of the velocity along the parallel. In terms of arclength
$s=\sqrt m\,l\,\theta$ the Gaussian curvature of the surface is $-f_{ss}/f$.
The oscillating and rotating regimes of the pendulum are the geodesics that
oscillate about, or wind around, the neck of the surface.
\end{example}

\begin{figure}[t]
\centering
\begin{tikzpicture}[scale=1.0]
  \draw[->] (-0.3,-2.0) -- (7.3,-2.0) node[right] {$q\in Q$};
  \draw[thick,domain=0.2:6.8,samples=80,smooth] plot (\x,{0.85+0.4*sin(58*\x)});
  \draw[thick,domain=0.2:6.8,samples=80,smooth] plot (\x,{-(0.85+0.4*sin(58*\x))});
  \foreach \x in {0.6,1.6,2.6,3.6,4.6,5.6,6.6}{
    \pgfmathsetmacro{\r}{0.85+0.4*sin(58*\x)}
    \draw (\x,0) ellipse ({0.17} and {\r});
    \draw[dotted] (\x,-2.0) -- (\x,{-\r});
  }
  \node at (0.6,1.55) {$f(q)$};
  \draw[->] (3.0,1.25) -- (3.0,0.02);
  \node[right] at (3.0,1.35) {\small hidden circle, angle $\varphi$};
\end{tikzpicture}
\caption{Schematic of Theorem~\ref{thm:hz-cyclic}: over each configuration
$q$ of the visible system sits a hidden circle of radius $f(q)=p/\sqrt{2V(q)}$.
Free (geodesic) motion on the total space projects to motion in the potential
$V$; the hidden circle spins faster where the radius is smaller, that is,
where the potential is higher.}
\label{fig:hz-warped}
\end{figure}

\subsection{Application to the double pendulum}

\begin{corollary}[Hertzian double pendulum]\label{cor:hz-dp}
Let $V_c=V+V_0$ with $V$ given by~\eqref{eq:hz-V} and a constant
$V_0>(m_1+m_2)gl_1+m_2gl_2$, and let $M$ be as in
Proposition~\ref{prop:hz-dp-metric}. On $\Tor^3=\Tor^2\times\Sph^1$ the metric
\[
  G=M_{ab}\,\ud\theta^a\ud\theta^b+\frac{p^2}{2V_c(\theta)}\,\ud\varphi^2
\]
makes Hertz's law (geodesic motion with constant speed) equivalent to the
double pendulum~\eqref{eq:hz-dp-grav} in the sense of
Theorem~\ref{thm:hz-cyclic}, with $\tfrac12G(\dot\gamma,\dot\gamma)=E+V_0$,
where $E=\tfrac12M\dot\theta\dot\theta+V$ is the pendulum's energy.
The hidden rotor turns with angular velocity
$\dot\varphi=2V_c(\theta(t))/p$.
\end{corollary}

The constant $V_0$ is necessary only because $f=p/\sqrt{2V_c}$ requires
$V_c>0$ and $V$ attains negative values; it shifts the zero of potential
energy and so, by Theorem~\ref{thm:hz-cyclic}, only the numerical value of the
total geodesic energy. All the double-pendulum quantities in the following
section are unchanged by it.

\subsection{How literal are the hidden masses?}

Hertz demands more than a Riemannian metric: hidden \emph{masses}, with
rigid connections to the visible ones. Three observations clarify what the
representation theorem does and does not give.

\paragraph{(a) A literal rotor, with an inertia correction.}
Let a hidden particle of unit mass be constrained to the position
$\mathbf y=(r(\theta)\cos\varphi,\,r(\theta)\sin\varphi,\,0)$, a point on a circle
whose radius is a prescribed function of the pendulum angles (a cam or linkage
in a real mechanism). The total ambient metric restricted to the constraint manifold is
\[
  G=\bigl(M^{\rm bare}_{ab}+\partial_ar\,\partial_br\bigr)\ud\theta^a\ud\theta^b+r^2\ud\varphi^2 ,
\]
so the hidden mass also contributes to the inertia of the visible variables,
and by Theorem~\ref{thm:hz-cyclic} (its proof uses only $V=p^2/(2f^2)$) the visible motion has effective metric
$M^{\rm bare}+\ud r\otimes\ud r$ and potential $p^2/(2r^2)$. Choosing $r=p/\sqrt{2V_c}$ gives
$\ud r\otimes\ud r=\dfrac{p^2}{8V_c^{3}}\,\ud V\otimes\ud V$, which is $O(V_0^{-3})$ at fixed $p$.
Thus the literal rotor reproduces the double pendulum with a bare inertia that differs from
the observed one by a correction that vanishes as the hidden energy $V_c$ grows relative to
the variation of $V$. Exactness is recovered only if one declares the \emph{effective}
inertia $M$ to be the observed one. The metric $G$ of Corollary~\ref{cor:hz-dp}
encodes precisely this declaration.

\paragraph{(b) Nash embedding: realizability at the level of geometry.}
By Nash's theorem any compact Riemannian manifold embeds isometrically in
some Euclidean space \cite{hdp:nash1956}. Applied to $(\Tor^3,G)$ this exhibits
the Hertzian configuration space as the constraint manifold of a system of
unit-mass coordinates with holonomic rigid connections. But in such a
realization the pendulum angles $\theta^a$ are merely functions on the
constraint manifold, not the positions of two identifiable bobs. The theorem
shows that Hertz's \emph{geometry} is always realizable; it does not show that
the hidden masses can be chosen compatibly with the visible ones.

\paragraph{(c) Non-uniqueness.}
Nothing in the visible motion selects the hidden structure. Replacing $V_c$ by
a different positive function changes only the energy bookkeeping; adding more
hidden rotors gives $V=\sum_kp_k^2/(2f_k^2)$ with many solutions; and the Jacobi
and Eisenhart constructions of Sections~\ref{sec:hz-jacobi}
and~\ref{sec:hz-compare} reproduce the same trajectories with entirely
different hidden content. The Hertzian explanation of gravity is therefore not
testable by observing the pendulum.

\section{A second geometrization: the Jacobi metric}
\label{sec:hz-jacobi}

Hertz enlarges the configuration space. There is also a way to keep its
dimension and to absorb the potential into the metric by a conformal
rescaling, at the price of fixing the energy.

\begin{theorem}[Jacobi]\label{thm:hz-jacobi}
On the Hill region $U_E=\{q\in Q: V(q)<E\}$ let
$J_E=(E-V)\,M$. The solutions of Newton's equations $\ddot q^a+\Gamma^a_{bc}\dot q^b\dot q^c=-M^{ab}\partial_bV$
of energy $E$ are, after the reparametrization $\ud\tau=(E-V)\,\ud t$, the
geodesics of $(U_E,J_E)$.
\end{theorem}

This is the geometric content of Maupertuis's principle: the abbreviated action
\[
  \int\sqrt{2(E-V)}\;\ud s_M
\]
is, up to the factor $\sqrt2$, the $J_E$-length of the path
\cite{hdp:jacobi1866,hdp:arnold1989}. (The relation of this principle to Hamilton's and
Hertz's is treated in the chapter on variational principles.)

\subsection{Curvature of the Jacobi metric}

In two dimensions the stability of nearby geodesics is governed by the Gaussian
curvature: for a unit-speed geodesic of a surface, the normal deviation $\xi$
between neighbouring geodesics obeys $\xi''+K\,\xi=0$ (Jacobi--Levi-Civita
equation), exponentially growing where $K<0$.

\begin{proposition}[Conformal formula]\label{prop:hz-conformal}
For $J_E=\Omega\,M$ with $\Omega=E-V>0$ on a surface $(Q,M)$,
\begin{equation}\label{eq:hz-KJ}
  K_{J_E}\;=\;\frac{1}{\Omega}\Bigl(K_M-\tfrac12\,\Lap_M\ln\Omega\Bigr),
\end{equation}
where $\Lap_M$ is the Laplace--Beltrami operator of $M$. Near a regular point of the
Hill boundary $\{V=E\}$, $K_{J_E}\simeq|\nabla\Omega|_M^2/(2\Omega^3)\to+\infty$.
\end{proposition}

\begin{proof}
For $\hat g=e^{2u}g$ on a surface, $\hat K=e^{-2u}(K-\Lap_gu)$ \cite{hdp:docarmo1992}.
Put $e^{2u}=\Omega$. For the boundary behaviour,
$\Lap\ln\Omega=\Lap\Omega/\Omega-|\nabla\Omega|^2/\Omega^2$, and the second term dominates.
(Formula~\eqref{eq:hz-KJ} was also checked symbolically against the Riemann tensor of $J_E$ for the double pendulum.)
\end{proof}

For the double pendulum $K_M$ is~\eqref{eq:hz-KM}, $\Omega=E-V$ with $V$
from~\eqref{eq:hz-V}, and~\eqref{eq:hz-KJ} gives $K_{J_E}$ in closed form.

\begin{proposition}[Negative curvature is unavoidable above the separatrix energy]\label{prop:hz-gb}
Let $E>V_{\max}=(m_1+m_2)gl_1+m_2gl_2$. Then $U_E=\Tor^2$, $J_E$ is a smooth
Riemannian metric on the torus, and $\int_{\Tor^2}K_{J_E}\,\ud A_{J_E}=0$. Hence, unless
$K_{J_E}\equiv0$, it takes negative values somewhere.
\end{proposition}

\begin{proof}
$V<V_{\max}<E$ everywhere, so $\Omega>0$ on all of $\Tor^2$. Apply Gauss--Bonnet
with $\chi(\Tor^2)=0$. (For $m_1=m_2=l_1=l_2=g=1$ one finds numerically
$K_{J_4}\in[-1.0,\,0.30]$, so $K_{J_E}\not\equiv0$.)
\end{proof}

\subsection{Where the negative curvature sits}

Figure~\ref{fig:hz-curvature} shows the sign of $K_{J_E}$ for the
unit-parameter pendulum ($m_1=m_2=l_1=l_2=g=1$, so $V=-2\cos\theta_1-\cos\theta_2$,
$-3\le V\le3$) at three energies, and Table~\ref{tab:hz-curv} the fraction of
the (Jacobi-)area where it is negative.

\begin{figure}[t]
\centering
\includegraphics[width=\textwidth]{jacobi_curvature.pdf}
\caption{Sign of the Gaussian curvature $K_{J_E}$ of the Jacobi metric of the
double pendulum on the torus $(\theta_1,\theta_2)\in[-\pi,\pi)^2$, for
$m_1=m_2=l_1=l_2=g=1$. Blue: $K_{J_E}>0$; red: $K_{J_E}<0$; grey: inaccessible
region $V\ge E$. At $E=-2$ the accessible region is a disc about the stable
equilibrium and $K_{J_E}>0$ throughout. At $E=1$ a red region has opened; at
$E=4>V_{\max}=3$ the whole torus is accessible and, as
Proposition~\ref{prop:hz-gb} requires, the curvature has both signs.}
\label{fig:hz-curvature}
\end{figure}

\begin{table}[t]
\centering
\begin{tabular}{@{}rcc@{}}
\toprule
$E$ & accessible part of the torus & fraction of $J_E$-area with $K_{J_E}<0$\\
\midrule
$-2$ & $0.13$ & $0$\\
$1$  & $0.69$ & $0.20$\\
$4$  & $1$    & $0.44$\\
\bottomrule
\end{tabular}
\caption{Accessible fraction of the coordinate torus and fraction of the Jacobi area
$\sqrt{\det J_E}\,\ud\theta_1\ud\theta_2$ on which the curvature is negative, for the
unit-parameter double pendulum (grid evaluation, two significant digits).}
\label{tab:hz-curv}
\end{table}

\begin{remark}[What the picture does and does not say]\label{rem:hz-caveat}
The pattern is the expected one. At low energy the motion is a perturbed
small oscillation and the whole accessible region has $K_{J_E}>0$; as energy
rises, negative-curvature regions open and grow, and above the energy at which
the rods can pass through the upright position the topology forces them to exist.
But three cautions are required.
(i) $K<0$ is sufficient for local exponential divergence of nearby geodesics, not necessary:
oscillating positive curvature can produce instability by parametric resonance, which is the
mechanism of the Hannay-angle chapter, and this is how chaos arises in many
systems whose average curvature is positive \cite{hdp:casetti1996,hdp:pettini2007}.
(ii) Remark~\ref{rem:hz-gb-free} shows that a torus metric with negative-curvature regions can
have integrable geodesic flow; the sign of $K$ locates \emph{possible} instability,
it does not by itself establish chaos.
(iii) $J_E$ degenerates at the Hill boundary and its curvature diverges there
(Proposition~\ref{prop:hz-conformal}), so the picture below the separatrix energy is a statement
about the interior of $U_E$ only.
\end{remark}

\section{Three geometrizations side by side}
\label{sec:hz-compare}

A third construction due to Eisenhart lifts the dynamics to a \emph{Lorentzian}
space. On $Q\times\Rb_t\times\Rb_z$ take
\begin{equation}\label{eq:hz-eisenhart}
  G_{\rm E}=M_{ab}\,\ud q^a\ud q^b+2\,\ud t\,\ud z-2V(q)\,\ud t^2 .
\end{equation}
The $t$-equation of the geodesic flow gives $\ddot t=0$, so $t$ is an affine
parameter, and the $q$-equation is $\ddot q^a+\Gamma^a_{bc}\dot q^b\dot q^c=-M^{ab}\partial_bV$ (with
$\Gamma^a_{tt}=M^{ab}\partial_bV$) \cite{hdp:eisenhart1929}. Table~\ref{tab:hz-compare} compares the three.

\begin{table}[t]
\centering
\small
\begin{tabular}{@{}lccc@{}}
\toprule
 & Hertz (hidden rotor) & Jacobi & Eisenhart\\
\midrule
Signature & Riemannian & Riemannian & Lorentzian\\
Dimension & $n+1$ & $n$ & $n+2$\\
Fixed data & hidden momentum $p$ & energy $E$ & none\\
Restriction on $V$ & $V+V_0>0$ & Hill region $V<E$ & none\\
Geodesic parameter & physical time $t$ & $\tau$, $\ud\tau=(E-V)\ud t$ & physical time $t$\\
Where $V$ sits & hidden kinetic energy & conformal factor & $G_{tt}$\\
Hertzian reading & hidden masses & none & none\\
\bottomrule
\end{tabular}
\caption{Three ways to turn the potential force into geometry. Only Hertz's
construction has a mechanical reading in terms of positive masses.}
\label{tab:hz-compare}
\end{table}

% ======================================================================
%  Section insert: Eisenhart's lifts in detail  (+ comparison table)
%  Drop into hertz_chapter.tex in place of the old "Three geometrizations
%  side by side" section (i.e. before \section{Assessment...}).
%
%  Needs: amsmath, amsthm (definition/theorem/proposition/remark/example
%  environments as defined in the chapter), enumitem, booktabs.
%  Cross-references used: thm:hz-cyclic, sec:hz-hidden, sec:hz-jacobi,
%  prop:hz-dp-metric, eq:hz-V, eq:hz-KM, eq:hz-warped.
%  Bibliography keys used: hdp:eisenhart1929, hdp:duval1985,
%  hdp:duval1991, hdp:casetti2000.
% ======================================================================

\section{Eisenhart's lifts in detail}
\label{sec:hz-eisenhart}

In 1928--29 L.~P.~Eisenhart showed that the trajectories of a natural
Lagrangian system can be obtained as geodesics of a metric on a larger space
\cite{hdp:eisenhart1929}. His paper contains \emph{two} constructions,
and they stand in different relations to Hertz. The first is
Lorentzian, works for every potential (including time-dependent ones) and
needs two extra coordinates; the second is Riemannian, needs a positive
potential and one extra coordinate. The second turns out to be
the metric of Theorem~\ref{thm:hz-cyclic}. The first was rediscovered
in 1984--85 by Duval, Burdet, K\"unzle and Perrin, who called the resulting
spaces \emph{Bargmann structures} \cite{hdp:duval1985,hdp:duval1991}, and
the construction is now often called the Eisenhart--Duval lift.

\subsection{The Lorentzian (null) lift}

\begin{definition}[Eisenhart lift]\label{def:hz-eis}
Let $(Q,M)$ be a Riemannian manifold and $V\in C^\infty(Q)$ arbitrary. On
$\widehat Q=Q\times\Rb_t\times\Rb_z$ put
\begin{equation}\label{eq:hz-eis-metric}
  G_{\rm E}=M_{ab}(q)\,\ud q^a\ud q^b+2\,\ud t\,\ud z-2V(q)\,\ud t^2 .
\end{equation}
\end{definition}

The $(t,z)$ block $\bigl(\begin{smallmatrix}-2V&1\\1&0\end{smallmatrix}\bigr)$
has determinant $-1$ whatever $V$ is, so $G_{\rm E}$ is a non-degenerate metric
of Lorentzian signature $(n{+}1,1)$. Its inverse has the components
\begin{equation}\label{eq:hz-eis-inverse}
  G^{ab}=M^{ab},\qquad G^{tz}=1,\qquad G^{zz}=2V,\qquad G^{tt}=G^{ta}=G^{za}=0 .
\end{equation}
The only Christoffel symbols that differ from those of $M$ and from zero are
\begin{equation}\label{eq:hz-eis-christoffel}
  \Gamma^a_{tt}=M^{ab}\partial_bV,\qquad
  \Gamma^z_{at}=\Gamma^z_{ta}=-\partial_aV .
\end{equation}
(For the double pendulum this was checked symbolically for a flat base and
confirmed for the curved base $M$ through the Ricci tensor below.)

\begin{theorem}[Eisenhart]\label{thm:hz-eis}
A curve $\gamma(s)=(q(s),t(s),z(s))$ is a geodesic of $G_{\rm E}$ if and only if
\begin{equation}\label{eq:hz-eis-geod}
  \ddot t=0,\qquad
  \ddot q^a+\Gamma^a_{bc}\dot q^b\dot q^c=-M^{ab}\partial_bV\,\dot t^{\,2},\qquad
  \ddot z=2\,\dot t\,\partial_aV\,\dot q^a .
\end{equation}
Consequently, if $\dot t=1$ (so $s=t$ is physical time):
\begin{enumerate}[label=(\alph*),nosep]
  \item $q(t)$ solves Newton's equations $\ddot q^a+\Gamma^a_{bc}\dot q^b\dot q^c=-M^{ab}\partial_bV$;
  \item $\dot z=2V(q)-\varepsilon$ with a constant $\varepsilon$, so $p_t:=G(\partial_t,\dot\gamma)=-\varepsilon$;
  \item $G_{\rm E}(\dot\gamma,\dot\gamma)=2(E-\varepsilon)$, where $E=\tfrac12M\dot q\dot q+V$ is the energy of $q(t)$;
  \item $\gamma$ is a \emph{null} geodesic exactly when $\varepsilon=E$, and then
        $z(t)=z_0-\int_0^t L\,\ud t'=z_0-S[q]$ with $L=\tfrac12M\dot q\dot q-V$.
\end{enumerate}
Conversely every solution $q(t)$ of Newton's equations lifts, for every
choice of $\varepsilon$, $z_0$ and $t_0$, to a geodesic with $\dot t=1$.
\end{theorem}

\begin{proof}
Equations~\eqref{eq:hz-eis-geod} are the geodesic equations with the symbols
\eqref{eq:hz-eis-christoffel}: $\Gamma^t_{AB}=0$ gives $\ddot t=0$;
$\ddot q^a+\Gamma^a_{bc}\dot q^b\dot q^c+\Gamma^a_{tt}\dot t^2=0$ is the second;
$\ddot z+2\Gamma^z_{at}\dot q^a\dot t=0$ is the third. For $\dot t=1$ the third says
$\ddot z=2\,\ud V/\ud t$, hence $\dot z=2V-\varepsilon$, and $p_t=-2V\dot t+\dot z=-\varepsilon$.
Then
\[
  G(\dot\gamma,\dot\gamma)=M\dot q\dot q+2\dot z-2V
  =2T+2(2V-\varepsilon)-2V=2(T+V-\varepsilon)=2(E-\varepsilon),
\]
which gives (c) and the first half of (d). If $\varepsilon=E$ then
$\dot z=2V-E=V-T=-L$. The converse is immediate by defining
$z$ through $\dot z=2V-\varepsilon$.
\end{proof}

\begin{remark}[The extra coordinate is the action]\label{rem:hz-z-action}
Statement (d) is the geometric content of the lift: along the physical
(null) geodesics the new coordinate $z$ is minus Hamilton's principal
function evaluated along the path. This is no accident. With $\dot t=1$ the
geodesic Lagrangian is $\tfrac12G(\dot\gamma,\dot\gamma)=L+\dot z$, so
$\int\tfrac12G(\dot\gamma,\dot\gamma)\,\ud t=S[q]+z(t_2)-z(t_1)$:
the energy functional of the lift is Hamilton's action up to endpoint terms.
\end{remark}

\subsection{Hamiltonian form and conserved quantities}

The geodesic Hamiltonian $H_{\rm E}=\tfrac12G^{AB}p_Ap_B$ is, by~\eqref{eq:hz-eis-inverse},
\begin{equation}\label{eq:hz-eis-ham}
  H_{\rm E}=\tfrac12M^{ab}p_ap_b+p_tp_z+V(q)\,p_z^{\,2}.
\end{equation}
Both $t$ and $z$ are cyclic for time-independent $V$, so $p_t$ and $p_z$ are conserved.
With $p_z=G(\partial_z,\dot\gamma)=\dot t=1$ and
$H(q,p)=\tfrac12M^{ab}p_ap_b+V$ the physical Hamiltonian, \eqref{eq:hz-eis-ham} becomes
\[
  H_{\rm E}=H(q,p)+p_t ,
\]
and the null condition $H_{\rm E}=0$ is $p_t=-H=-E$: Hamilton's equations of $H_{\rm E}$
are the equations of the \emph{extended phase space}, in which energy is the momentum conjugate
to time and the dynamics is a constraint. The constant $p_z$ plays the role of a mass
(in the Bargmann picture it is the mass or the electric charge of the particle), and the
fixed value $p_z=1$ is what selects Newtonian dynamics with unit
inertia. Hamilton's equations give $\dot t=\partial H_{\rm E}/\partial p_t=p_z=1$ and
$\dot z=\partial H_{\rm E}/\partial p_z=p_t+2V=2V-E$, as found above.

\subsection{Geometry of the lift: a pp-wave}

\begin{proposition}[Curvature of the Eisenhart metric]\label{prop:hz-eis-curv}
\leavevmode
\begin{enumerate}[label=(\roman*),nosep]
  \item $\nabla\,\partial_z=0$: the null vector field $\partial_z$ is covariantly constant.
  \item The Ricci tensor of $G_{\rm E}$ has only the components
        \[
          \operatorname{Ric}_{ab}=\operatorname{Ric}^{M}_{ab},\qquad
          \operatorname{Ric}_{tt}=\Lap_MV .
        \]
\end{enumerate}
\end{proposition}

\begin{proof}
(i) is $\Gamma^A_{Bz}=0$, read off from \eqref{eq:hz-eis-christoffel}. For (ii) I computed the Ricci
tensor symbolically for a flat base with arbitrary $V(x,y)$ and, for the
unit-parameter double pendulum with its curved metric $M$, at random points, where the
only non-vanishing components were $\operatorname{Ric}_{ab}=K_MM_{ab}$ and
$\operatorname{Ric}_{tt}=\Lap_MV$ to rounding error; the general statement is the standard curvature
property of Eisenhart metrics \cite{hdp:casetti2000}.
\end{proof}

A Lorentzian manifold with a covariantly constant null vector whose metric has this
block form is a (generalized) \emph{pp-wave}, and the structure
$(\widehat Q,G_{\rm E},\partial_z)$ is a Bargmann structure \cite{hdp:duval1985,hdp:duval1991}.
Two consequences are worth stating.
\begin{itemize}[nosep]
  \item \emph{Gravity as curvature.} With $\Lap_MV=4\pi G\rho$ for a Newtonian
        potential, $\operatorname{Ric}_{tt}=\Lap_MV$ is the Poisson equation, and Ricci-flatness of the lift
        (flat $M$, harmonic $V$) is Newtonian gravity in vacuum. The lift makes Newtonian
        gravity \emph{look} like general relativity; it is a statement about a Lorentzian
        space of one dimension more than spacetime, not about spacetime.
  \item \emph{Tidal tensor.} The geodesic deviation equation of $G_{\rm E}$ along a lifted
        trajectory, restricted to the $q$-directions, is the usual variational equation of the
        dynamics,
        \[
          \frac{D^2\xi^a}{\ud t^2}+R^{M\,a}{}_{bcd}\,\dot q^b\xi^c\dot q^d
          +M^{ab}\nabla_b\nabla_cV\,\xi^c=0,
        \]
        and the trace of the tidal tensor is $\Lap_MV$ (the origin of $\operatorname{Ric}_{tt}$)
        \cite{hdp:casetti2000}. In contrast to the Jacobi metric of
        Section~\ref{sec:hz-jacobi}, no factor $E-V$ enters and there is no Hill-boundary singularity;
        but the price is that the criterion for exponential instability is the standard
        linear-stability computation, not a purely geometric sign condition.
\end{itemize}

\subsection{Examples}

\begin{example}[Harmonic oscillator]\label{ex:hz-eis-ho}
For $Q=\Rb$, $M=\ud q^2$ and $V=\tfrac12\omega^2q^2$,
$G_{\rm E}=\ud q^2+2\,\ud t\,\ud z-\omega^2q^2\ud t^2$ is a plane wave with $\operatorname{Ric}_{tt}=\omega^2$.
The null geodesics with $\dot t=1$ are $q=A\cos\theta$, $\theta=\omega t+\phi$, and
\[
  z=z_0+\tfrac14A^2\omega\,\sin2\theta ,
\]
which is $z_0-\int L\,\ud t$ with $L=-\tfrac12A^2\omega^2\cos2\theta$. The $z$-motion is a
bounded oscillation at twice the frequency.
\end{example}

\begin{example}[Double pendulum]\label{ex:hz-eis-dp}
With $M$ from Proposition~\ref{prop:hz-dp-metric} and $V$ from~\eqref{eq:hz-V},
\[
\begin{aligned}
  G_{\rm E}={}&(m_1+m_2)\,l_1^2\,\ud\theta_1^2+2m_2l_1l_2\cos(\theta_1-\theta_2)\,\ud\theta_1\ud\theta_2
  +m_2l_2^2\,\ud\theta_2^2\\
  &+2\,\ud t\,\ud z+2g\bigl[(m_1+m_2)l_1\cos\theta_1+m_2l_2\cos\theta_2\bigr]\ud t^2 .
\end{aligned}
\]
This is a four-dimensional Lorentzian manifold $\Tor^2\times\Rb^2$ with
$\operatorname{Ric}_{ab}=K_MM_{ab}$ (with $K_M$ from~\eqref{eq:hz-KM}) and $\operatorname{Ric}_{tt}=\Lap_MV$.
Its Killing fields include $\partial_z$ (null, parallel) and $\partial_t$; the rotation
$\partial_{\theta_1}+\partial_{\theta_2}$ that was a symmetry of the free pendulum is
broken by $V$. The null geodesics with $\dot t=1$ project to the chaotic trajectories of the pendulum,
and $z$ along them is $-S[\theta]$.
\end{example}

\subsection{The Riemannian lift, and its identity with the hidden rotor}

Eisenhart's second construction, for \emph{time-independent potentials that are positive} on $Q$,
uses one extra coordinate and a Riemannian metric of the form
$M+A(q)\,\ud z^2$ with $A$ built from $V$
\cite{hdp:eisenhart1929,hdp:casetti2000}. Taking $A=1/V_c$ with $V_c=V+c>0$,
\begin{equation}\label{eq:hz-eis-riem}
  G_{\rm R}=M_{ab}\,\ud q^a\ud q^b+\frac{1}{V_c(q)}\,\ud z^2 ,
\end{equation}
the Hamiltonian is $\tfrac12M^{ab}p_ap_b+\tfrac12V_c\,p_z^2$, which equals
the physical Hamiltonian $\tfrac12M^{ab}p_ap_b+V_c$ for $p_z=\sqrt2$. This is exactly
the metric~\eqref{eq:hz-warped} of Theorem~\ref{thm:hz-cyclic} with $f^2=p^2/(2V_c)$, $p=\sqrt2$,
and $\varphi=z$. The theorem therefore \emph{is} Eisenhart's Riemannian lift, with a different proof and with
a physical interpretation of the extra coordinate added: Eisenhart's $z$ is Hertz's hidden
cyclic coordinate, and $V=\tfrac12f^2\dot z^2$ is the kinetic energy of the hidden rotor.
(Hertz did not write down the warped metric; the idea that a potential is the kinetic energy of
a cyclic hidden motion is his, and the exact geometric statement is Eisenhart's.)

\subsection{Hertzian or not?}

Does the Lorentzian lift satisfy Hertz's requirements? Only in part. Its dynamics is
free in the sense that the law is geodesic, and no force appears. But Hertz's mass metric is positive
definite: kinetic energy is $\tfrac12\sum m\dot x^2\ge0$, and the straightest path is defined by minimizing
a norm of acceleration. In signature $(n{+}1,1)$ there is no such norm, the physical
trajectories are \emph{null}, and the extra coordinate $z$ is a null direction, not a hidden mass.
The Lorentzian lift is thus forceless but not Hertzian; the Riemannian lift is both.
That the two lifts are available side by side underlines the point of
Section~\ref{sec:hz-hidden}(c): the geometric content of the visible dynamics does not fix its higher-dimensional
realization.

\section{Four geometrizations side by side}
\label{sec:hz-compare}

\begin{table}[t]
\centering
\small
\setlength{\tabcolsep}{4pt}
\begin{tabular}{@{}lccc@{}}
\toprule
 & Hertz / Eisenhart--Riem. & Jacobi & Eisenhart--Lorentz\\
\midrule
Signature & Riemannian & Riemannian & Lorentzian $(n{+}1,1)$\\
Dimension & $n+1$ & $n$ & $n+2$\\
Fixed data & $p_z$ (hidden momentum) & energy $E$ & $p_z=1$, null condition\\
Restriction on $V$ & $V+c>0$ & Hill region $V<E$ & none\\
Geodesic parameter & physical time $t$ & $\tau$, $\ud\tau=(E-V)\ud t$ & physical time $t$\\
Where $V$ sits & hidden kinetic energy & conformal factor & $G_{tt}$\\
Role of extra coordinates & hidden angle $z$ & none & $t$ and null $z=-S$\\
Time-dependent $V$ & no & no & yes\\
Hertzian reading & hidden masses & none & none (null)\\
\bottomrule
\end{tabular}
\caption{Three ways (four constructions, since the first column is Hertz's idea and Eisenhart's
Riemannian metric at once) to turn the potential force into geometry. Only the first has a
mechanical reading in terms of positive masses.}
\label{tab:hz-compare}
\end{table}

Table~\ref{tab:hz-compare} summarizes the comparison. The Lorentzian lift is the only one that
treats time-dependent potentials uniformly, since $t$ is a coordinate of the lifted
space; the Jacobi construction requires conservation of energy and the Riemannian lift
requires a time-independent positive $V$.

\section{Assessment: what Hertz's program delivers}
\label{sec:hz-assess}

\paragraph{What it gains.}
Applied to the double pendulum the program does what it promises, in one
respect: it removes every force from the statement of the law. The tension in
the rods is a reaction of the rigid connection, the centrifugal-looking terms
are Christoffel symbols, and gravity is the kinetic energy of a rotor. The
geometry also yields something new that is invisible in the force picture:
the Gauss--Bonnet argument of Proposition~\ref{prop:hz-gb} explains, without any
numerics, why the Jacobi geometry of the double pendulum above the upright
energy must contain regions in which nearby motions diverge.

\paragraph{What it costs.}
The cost is the hidden structure, and the analysis of Section~\ref{sec:hz-hidden}
shows that it is both inessential and unconstrained. The visible motion
of the pendulum determines the ordinary differential equation~\eqref{eq:hz-dp-grav}
and nothing else; the hidden rotor, the Jacobi metric and the Eisenhart
lift are three unrelated ways of writing that same equation as a geodesic problem,
each in a different dimension and with different auxiliary data. To the extent
that observation of the pendulum is all there is, the hidden masses
are idle wheels. This is the pattern argued throughout this book: the equations
of motion are the primary, real-world content of mechanics, and variational and
geometric formulations are ideal-world devices for finding and organizing them.
The three geometrizations are in this sense charts of one object, each valid
on its own domain and each introducing singularities of its own: the positivity
requirement $V+V_0>0$, the Hill boundary, the Lorentzian signature.

\paragraph{What it cannot reach.}
Hertz's systems are closed and conservative: all forces, including friction,
must be traced to hidden conservative motion. Dissipation therefore enters
only as a statement that the hidden energy is out of sight, and genuine open
systems, such as the damped double pendulum, are outside the scheme as stated;
they are the subject of the chapters on damping and on open systems.
The nonholonomic extension, in which straightest paths differ from shortest
paths, is the other place where Hertz's mechanics is not simply geodesic flow
(Section~\ref{sec:hz-geometry}).

\paragraph{Historical coda.}
The reception was divided between those who found the hidden-mass hypothesis a
barrier to physical content and those who appreciated the geometric form
\cite{hdp:lutzen2005}. What survived was less the philosophy than the geometry:
the representation of constrained mechanics as geodesic flow on a Riemannian
manifold, which is the standard language of the modern geometric treatments.

\section*{Problems}
\addcontentsline{toc}{section}{Problems}

\begin{enumerate}[label=\textbf{\thechapter.\arabic*}.]
  \item Prove Proposition~\ref{prop:hz-law}(i) by Lagrange multipliers: minimize
        $g(a,a)$ subject to $\omega_\alpha(a)=c_\alpha$ and show the multipliers
        are the constraint reactions.
  \item Derive~\eqref{eq:hz-dp-free} directly from Proposition~\ref{prop:hz-law}(iii) with
        the constraints $|\mathbf x_1|=l_1$, $|\mathbf x_2-\mathbf x_1|=l_2$, eliminating the
        multipliers (the rod tensions).
  \item Identify the second integral of the free double pendulum with
        $J$ of Proposition~\ref{prop:hz-dp-curvature}, and find the dependence of the
        period of the $\delta$-motion on $J$ and the speed.
  \item In Example~\ref{ex:hz-simple} show that the surface of revolution is
        smooth, that $K=-f_{ss}/f$ does not depend on $p$, and determine where $K<0$
        (compare Gauss--Bonnet for the torus).
  \item Verify the claim of Section~\ref{sec:hz-hidden}(a) that $\ud r\otimes\ud r=O(V_0^{-3})$
        at fixed $p$, and compute the leading correction to the pendulum's small-oscillation frequency.
  \item Verify that the geodesics of~\eqref{eq:hz-eisenhart} with $\dot t=1$ project to
        solutions of~\eqref{eq:hz-dp-grav}, and find the quantity conserved by the
        null Killing vector $\partial_z$.
  \item Show that for $g\neq0$ the vector field $\partial_{\theta_1}+\partial_{\theta_2}$ is
        \emph{not} a Killing field of $J_E$, so the angular momentum $J$ of
        Proposition~\ref{prop:hz-dp-curvature} is no longer conserved. Why does this loss of
        symmetry not by itself prove that the flow is non-integrable?
\end{enumerate}

\section*{Notes and references}
\addcontentsline{toc}{section}{Notes and references}

\emph{On verification.} The kinetic metric, the geodesic-versus-Lagrange form of the equations of
motion, the explicit accelerations~\eqref{eq:hz-dp-free}--\eqref{eq:hz-dp-grav}, the Christoffel symbols
of the warped metric, the curvature~\eqref{eq:hz-KM} and the conformal
formula~\eqref{eq:hz-KJ} were checked symbolically (SymPy); the curvature map and the table were computed
numerically from~\eqref{eq:hz-KJ} on a $900\times900$ grid, and the Gauss--Bonnet integral at $E=4$
vanishes to rounding error, as it should.

\begingroup
\let\chapter\section
\renewcommand{\bibname}{References}
\begin{thebibliography}{99}
\bibitem{hdp:hertz1894} H.~Hertz, \emph{Die Prinzipien der Mechanik in neuem Zusammenhange dargestellt},
  Gesammelte Werke, Vol.~III (J.~A.~Barth, Leipzig, 1894), with a preface by H.~von Helmholtz.
  English translation: \emph{The Principles of Mechanics Presented in a New Form}, trans.\ D.~E.~Jones and
  J.~T.~Walley (Macmillan, London, 1899; repr.\ Dover, New York, 1956).
\bibitem{hdp:lutzen2005} J.~L\"utzen, \emph{Mechanistic Images in Geometric Form: Heinrich Hertz's
  ``Principles of Mechanics''} (Oxford University Press, Oxford, 2005).
\bibitem{hdp:gauss1829} C.~F.~Gauss, \"Uber ein neues allgemeines Grundgesetz der Mechanik,
  \emph{J.\ reine angew.\ Math.}\ \textbf{4} (1829) 232--235.
\bibitem{hdp:eom-hertz} V.~V.~Rumyantsev, Hertz principle, in \emph{Encyclopedia of Mathematics} (EMS Press).
\bibitem{hdp:routh1877} E.~J.~Routh, \emph{A Treatise on the Stability of a Given State of Motion}
  (Macmillan, London, 1877).
\bibitem{hdp:helmholtz1884} H.~von Helmholtz, Principien der Statik monocyklischer Systeme,
  \emph{J.\ reine angew.\ Math.}\ \textbf{97} (1884) 111--140, 317--336.
\bibitem{hdp:nash1956} J.~Nash, The imbedding problem for Riemannian manifolds,
  \emph{Ann.\ of Math.}\ \textbf{63} (1956) 20--63.
\bibitem{hdp:jacobi1866} C.~G.~J.~Jacobi, \emph{Vorlesungen \"uber Dynamik}, ed.\ A.~Clebsch
  (G.~Reimer, Berlin, 1866).
\bibitem{hdp:arnold1989} V.~I.~Arnold, \emph{Mathematical Methods of Classical Mechanics},
  2nd ed., Graduate Texts in Mathematics \textbf{60} (Springer, New York, 1989).
\bibitem{hdp:docarmo1992} M.~P.~do Carmo, \emph{Riemannian Geometry} (Birkh\"auser, Boston, 1992).
\bibitem{hdp:eisenhart1929} L.~P.~Eisenhart, Dynamical trajectories and geodesics,
  \emph{Ann.\ of Math.}\ \textbf{30} (1929) 591--606.
\bibitem{hdp:casetti1996} L.~Casetti, C.~Clementi and M.~Pettini, Riemannian theory of Hamiltonian
  chaos and Lyapunov exponents, \emph{Phys.\ Rev.\ E} \textbf{54} (1996) 5969--5984.
\bibitem{hdp:pettini2007} M.~Pettini, \emph{Geometry and Topology in Hamiltonian Dynamics and
  Statistical Mechanics} (Springer, New York, 2007).
\end{thebibliography}
\endgroup
 from a `book' preamble that loads
%    amsmath, amssymb, amsthm, mathtools, enumitem, booktabs, graphicx,
%    tikz (+ arrows.meta library), hyperref.
%  Figure file needed in the same directory:  jacobi_curvature.pdf
%  Theorem-type environments are defined below only if not yet defined.
% ======================================================================

\makeatletter
\@ifundefined{theorem}{\theoremstyle{plain}\newtheorem{theorem}{Theorem}[chapter]}{}
\@ifundefined{proposition}{\theoremstyle{plain}\newtheorem{proposition}[theorem]{Proposition}}{}
\@ifundefined{corollary}{\theoremstyle{plain}\newtheorem{corollary}[theorem]{Corollary}}{}
\@ifundefined{definition}{\theoremstyle{definition}\newtheorem{definition}[theorem]{Definition}}{}
\@ifundefined{example}{\theoremstyle{definition}\newtheorem{example}[theorem]{Example}}{}
\@ifundefined{remark}{\theoremstyle{remark}\newtheorem{remark}[theorem]{Remark}}{}
\makeatother

\providecommand{\Rb}{\mathbb{R}}
\providecommand{\Tor}{\mathbb{T}}
\providecommand{\Sph}{\mathbb{S}}
\providecommand{\ud}{\mathrm{d}}
\providecommand{\Lap}{\Delta}
\setlength{\emergencystretch}{2.5em}

\chapter{Hertz's Forceless Mechanics and the Double Pendulum}
\label{ch:hertz-double-pendulum}

\section{A mechanics without forces}
\label{sec:hz-intro}

Newton's \emph{Principia} takes force as a primitive. The concept is so
deeply built into the textbook presentation of mechanics that one forgets
how uneasy its status has always been: a force is neither a displacement
nor a mass, it cannot be seen, and the classical examples---gravitation
across empty space, the tension of a string, the reaction of a table---have
little in common except that each can be written on the right-hand side of
$m\ddot{\mathbf x}=\mathbf F$. In \emph{Die Prinzipien der Mechanik in neuem
Zusammenhange dargestellt}, which appeared in 1894 after Hertz's death
and with a preface by Helmholtz, Heinrich Hertz proposed to remove force
from the foundations altogether \cite{hdp:hertz1894}. His own ``image''
of mechanics rests on three notions only---time, space and mass---and on
a single law, which in the Jones--Walley translation of \S309 reads:
``Every free system persists in its state of rest or of uniform motion in
a straightest path.''

The sentence has two moving parts. A \emph{free} system is one subject to
no forces whatever; it is subject only to rigid connections between the
positions of its parts. A \emph{straightest path} is a purely geometric
notion, the path of least curvature compatible with those connections. What
the Newtonian calls a force is then, in Hertz's account, of one of two
kinds. It is either the reaction of a rigid connection among the visible
bodies, or it is the visible trace of \emph{hidden masses}, rigidly
connected to the visible ones and moving in ways we do not observe
\cite{hdp:hertz1894,hdp:lutzen2005}. Helmholtz's preface to the book already
voices reservations on whether hidden systems of the required kind can
actually be found for the forces of nature, and the reception of the
book was, for that reason, largely respectful and largely unconvinced.

This chapter takes Hertz's program seriously as \emph{mathematics} and
asks what it does when applied to the smallest system in which all its
ingredients meet: the planar double pendulum. The double pendulum has
rigid connections, which are Hertz's native material (the two rods); it has
a force, gravity, which Hertz must explain away; and it has chaos, which
any mechanics must be able to represent. The plan is as follows.

\begin{itemize}[nosep]
  \item Section~\ref{sec:hz-geometry} states Hertz's law precisely and shows
        that, for holonomic connections, it is the statement that motion is
        a geodesic of the mass-weighted configuration space.
  \item Section~\ref{sec:hz-free-dp} shows that the double pendulum
        \emph{without} gravity is a free Hertz system, computes the
        curvature of its configuration space, and finds a hidden Killing
        field.
  \item Section~\ref{sec:hz-hidden} switches gravity on and eliminates it in
        Hertz's way: by a representation theorem that turns any positive
        potential into the kinetic energy of one hidden cyclic coordinate.
  \item Section~\ref{sec:hz-jacobi} gives the second geometrization of the
        same dynamics, the Jacobi metric at fixed energy, and uses it to ask
        where the chaos of the double pendulum sits in the geometry.
  \item Sections~\ref{sec:hz-compare} and~\ref{sec:hz-assess} compare the
        geometrizations and assess what the Hertzian program does and does
        not deliver, in the light of the theme of this book: the equations of
        motion, not the formulations built around them, are the real-world
        content of mechanics.
\end{itemize}

\section{Hertz's law and the geometry of configuration space}
\label{sec:hz-geometry}

\subsection{Systems, connections, free motion}

\begin{definition}[Hertzian material system]\label{def:hz-system}
A \emph{material system} consists of $N$ particles with masses $m_i>0$ and
positions $\mathbf x_i\in\Rb^3$. Its \emph{mass-weighted ambient space} is
$(\Rb^{3N},g)$ with
\[
  g(\xi,\eta)\;=\;\sum_{i=1}^{N}m_i\,\xi_i\!\cdot\eta_i .
\]
A \emph{rigid connection} is a relation that is linear and homogeneous in
the differentials of the coordinates, $\omega_\alpha(x)(\ud x)=0$,
$\alpha=1,\dots,k$, where the $\omega_\alpha$ are Pfaffian forms depending
on position. Write
$D_x=\bigcap_\alpha\ker\omega_\alpha(x)\subset\Rb^{3N}$ for the space of
\emph{admissible velocities}. The connections are \emph{holonomic} if $D$
is integrable, i.e.\ if there is a submanifold $\mathcal N\subset\Rb^{3N}$
through every admissible position with $D=T\mathcal N$; otherwise they are
\emph{nonholonomic} (both terms are Hertz's). A system is \emph{free} if
it is subject to rigid connections only, with no active forces.
\end{definition}

Hertz measures path elements by the same mass-weighted expression
$g(\ud x,\ud x)=\sum m_i\,\ud x_i^2$ (he divides by the total mass, which is
immaterial here). Differentiating an admissible motion $x(t)$, with
$\dot x\in D_x$, along the connections gives, for the acceleration $a=\ddot x$,
\[
  \omega_\alpha(x)(a)\;=\;c_\alpha(x,\dot x),
  \qquad
  c_\alpha:=-\,\partial_\beta\omega_{\alpha,\gamma}\,\dot x^\beta\dot x^\gamma ,
\]
so that at a given state $(x,\dot x)$ the admissible accelerations form an
affine space $A_{(x,\dot x)}=a_0+D_x$.

\begin{definition}[Straightest path]\label{def:hz-straightest}
An admissible path is \emph{straightest} if at every point its curvature
is the least among all admissible paths having the same position and the
same tangent. For a path parametrized by constant speed, curvature is
proportional to $|a|_g$, so this means
\[
  a\;=\;\operatorname*{arg\,min}_{a'\in A_{(x,\dot x)}}\ g(a',a').
\]
\end{definition}

\begin{proposition}[Hertz's law in Newtonian language]\label{prop:hz-law}
Let a free system move along a straightest path. Then:
\begin{enumerate}[label=(\roman*),nosep]
  \item $g(a,w)=0$ for every $w\in D_x$, i.e.\ $a\in D_x^{\perp}$;
  \item the speed $g(\dot x,\dot x)^{1/2}$ is constant;
  \item in components, $m_i\ddot{\mathbf x}_i=\mathbf R_i$ with a reaction
        $R=(\mathbf R_i)$ satisfying $g^{-1}\!R\perp D_x$, that is, the
        equations of d'Alembert--Lagrange with zero active forces;
  \item if the connections are holonomic, $D=T\mathcal N$, then $x(t)$ is a
        geodesic of the induced Riemannian manifold $(\mathcal N,g|_{\mathcal N})$.
\end{enumerate}
\end{proposition}

\begin{proof}
(i) Minimizing $g(a',a')$ over the affine space $a_0+D_x$ yields the
orthogonal projection of the origin onto it, which is orthogonal to the
direction space $D_x$. (ii) $\frac{\ud}{\ud t}g(\dot x,\dot x)=2g(a,\dot x)=0$
because $\dot x\in D_x$. (iii) Write the reaction as $R:=g(a,\cdot)$; then
$R(w)=0$ on $D_x$ is (i). (iv) Let $\nabla$ be the flat connection of
$\Rb^{3N}$. For a curve in $\mathcal N$ the Gauss formula splits $a=\nabla_{\dot x}\dot x$
into a part tangent to $\mathcal N$, which equals the covariant acceleration
$\nabla^{\mathcal N}_{\dot x}\dot x$ of the induced metric, and a normal part.
By (i) the tangential part vanishes.
\end{proof}

\begin{remark}[Gauss, and what is new]
Gauss's principle of least constraint minimizes
\[
  \sum_i m_i\Bigl|\mathbf a_i-\frac{\mathbf F_i}{m_i}\Bigr|^2
\]
over admissible accelerations \cite{hdp:gauss1829}; Hertz's law is the case
$\mathbf F=0$, and for stationary connections without active forces the two
coincide \cite{hdp:eom-hertz}. The mathematical gain is therefore nil; the
gain is conceptual. In Gauss's principle forces are given data and
constraints are restrictions; in Hertz's the forces are to be \emph{derived},
and the whole of dynamics becomes a statement about the metric $g$ and the
distribution $D$.
\end{remark}

\begin{remark}[Nonholonomic connections]
For nonholonomic $D$ a straightest path is in general neither a geodesic of
any induced metric nor a minimizer of length among admissible paths; Hertz
stressed that Hamilton's principle and the principle of least action fail for
such systems \cite{hdp:lutzen2005}. The straightest-path condition (i) is the
Lagrange--d'Alembert equation, not the variational (``vakonomic'') one. The
double pendulum has holonomic connections only, so none of this affects the
rest of the chapter.
\end{remark}

\subsection{Holonomic Hertzian mechanics is geodesic flow}

For holonomic connections, part (iv) says that the entire dynamics of a free
system is the geodesic flow of $(\mathcal N,M)$, where in coordinates
$q^a$ on $\mathcal N$ the \emph{kinetic metric} is
\begin{equation}\label{eq:hz-kinetic-metric}
  M_{ab}(q)\;=\;\sum_i m_i\,\frac{\partial\mathbf x_i}{\partial q^a}\!\cdot\frac{\partial\mathbf x_i}{\partial q^b},
\end{equation}
and the equations of motion are
\begin{equation}\label{eq:hz-geodesic}
  \ddot q^a+\Gamma^a_{bc}(q)\,\dot q^b\dot q^c=0 ,
  \qquad
  \Gamma^a_{bc}=\tfrac12M^{ad}\bigl(\partial_bM_{dc}+\partial_cM_{db}-\partial_dM_{bc}\bigr).
\end{equation}
The ``centrifugal'' and ``Coriolis'' terms of textbook mechanics are the
Christoffel symbols $\Gamma$: they are not forces but the coordinate
expression of the connection of $M$. What remains outside this scheme, and what
Hertz must account for in order to keep his claim of generality, is every
term that is \emph{not} of this form, notably a potential force on the
right-hand side of~\eqref{eq:hz-geodesic}.

\section{The double pendulum without gravity is a free Hertz system}
\label{sec:hz-free-dp}

\begin{figure}[t]
\centering
\begin{tikzpicture}[scale=1.15,>=Stealth]
  % ceiling
  \draw[thick] (-1.0,0) -- (1.0,0);
  \foreach \x in {-0.9,-0.6,...,0.9}{\draw (\x,0) -- (\x-0.15,0.18);}
  \coordinate (O) at (0,0);
  \coordinate (A) at ({2*sin(28)},{-2*cos(28)});
  \coordinate (B) at ({2*sin(28)+1.6*sin(62)},{-2*cos(28)-1.6*cos(62)});
  % verticals
  \draw[dashed,gray] (O) -- (0,-3.2);
  \draw[dashed,gray] (A) -- ++(0,-1.9);
  % rods
  \draw[very thick] (O) -- (A) node[midway,above right=-2pt] {$l_1$};
  \draw[very thick] (A) -- (B) node[midway,above right=-1pt] {$l_2$};
  % bobs
  \fill (O) circle (1.6pt);
  \filldraw[fill=white,thick] (A) circle (0.13) node[right=5pt] {$m_1$};
  \filldraw[fill=white,thick] (B) circle (0.13) node[right=5pt] {$m_2$};
  % angles
  \draw[->] (0,-0.85) arc[start angle=-90,end angle=-62,radius=0.85];
  \node at (0.22,-1.12) {\small$\theta_1$};
  \draw[->] (A)++(0,-0.7) arc[start angle=-90,end angle=-28,radius=0.7];
  \node at ({2*sin(28)+0.42},{-2*cos(28)-0.98}) {\small$\theta_2$};
  % gravity
  \draw[->,thick] (3.2,-0.4) -- (3.2,-1.4) node[midway,right] {$g$};
\end{tikzpicture}
\caption{The planar double pendulum. Both angles are measured from the
downward vertical, so the configuration manifold is the torus
$\Tor^2=\Sph^1\times\Sph^1$ (the rods may swing through the upright position).}
\label{fig:hz-dp}
\end{figure}

Let the two bobs have masses $m_1,m_2$ and positions $\mathbf x_1,\mathbf x_2$
in a vertical plane, joined by massless rigid rods of lengths $l_1,l_2$ to the
pivot and to each other (Figure~\ref{fig:hz-dp}). The mass-weighted ambient
space is $\Rb^4$ with $g=\operatorname{diag}(m_1,m_1,m_2,m_2)$, and the
connections are the two holonomic relations
$|\mathbf x_1|=l_1$, $|\mathbf x_2-\mathbf x_1|=l_2$. They define
$\mathcal N=\Tor^2$, parametrized by
\[
  \mathbf x_1=l_1(\sin\theta_1,-\cos\theta_1),\qquad
  \mathbf x_2=\mathbf x_1+l_2(\sin\theta_2,-\cos\theta_2).
\]

\begin{proposition}[Kinetic metric]\label{prop:hz-dp-metric}
With $\delta:=\theta_1-\theta_2$, the induced metric~\eqref{eq:hz-kinetic-metric} is
\[
  M\;=\;\begin{pmatrix}(m_1+m_2)\,l_1^2 & m_2\,l_1l_2\cos\delta\\[2pt]
                       m_2\,l_1l_2\cos\delta & m_2\,l_2^2\end{pmatrix},
  \qquad
  \det M=m_2\,l_1^2l_2^2\,\bigl(m_1+m_2\sin^2\delta\bigr).
\]
\end{proposition}

\begin{proof}
Differentiate the parametrization: $\partial_{\theta_1}\mathbf x_1=l_1(\cos\theta_1,\sin\theta_1)$,
$\partial_{\theta_1}\mathbf x_2=\partial_{\theta_1}\mathbf x_1$,
$\partial_{\theta_2}\mathbf x_1=0$, $\partial_{\theta_2}\mathbf x_2=l_2(\cos\theta_2,\sin\theta_2)$.
Then $M_{11}=(m_1+m_2)l_1^2$, $M_{22}=m_2l_2^2$ and
$M_{12}=m_2l_1l_2(\cos\theta_1\cos\theta_2+\sin\theta_1\sin\theta_2)=m_2l_1l_2\cos\delta$.
\end{proof}

The metric depends on $\theta_1,\theta_2$ only through $\delta$. Hertz's
law for the gravity-free pendulum is therefore the geodesic
equation~\eqref{eq:hz-geodesic} of $(\Tor^2,M)$, which reads explicitly
\begin{equation}\label{eq:hz-dp-free}
\begin{aligned}
  \ddot\theta_1&=-\frac{m_2\,l_1\dot\theta_1^2\sin\delta\cos\delta+m_2\,l_2\dot\theta_2^2\sin\delta}{l_1\,(m_1+m_2\sin^2\delta)},\\[2pt]
  \ddot\theta_2&=\frac{(m_1+m_2)\,l_1\dot\theta_1^2\sin\delta+m_2\,l_2\dot\theta_2^2\sin\delta\cos\delta}{l_2\,(m_1+m_2\sin^2\delta)} .
\end{aligned}
\end{equation}

\begin{proposition}[Curvature and a hidden symmetry]\label{prop:hz-dp-curvature}
The Gaussian curvature of $(\Tor^2,M)$ is
\begin{equation}\label{eq:hz-KM}
  K_M\;=\;\frac{m_1\cos\delta}{l_1l_2\,\bigl(m_1+m_2\sin^2\delta\bigr)^{2}} .
\end{equation}
The vector field $\partial_{\theta_1}+\partial_{\theta_2}$ is Killing, with Noether
integral
\[
  J=\bigl[(m_1+m_2)l_1^2+m_2l_1l_2\cos\delta\bigr]\dot\theta_1
   +\bigl[m_2l_2^2+m_2l_1l_2\cos\delta\bigr]\dot\theta_2 ,
\]
the total angular momentum about the pivot. Together with the
speed, $J$ makes the geodesic flow of $(\Tor^2,M)$ Liouville integrable.
\end{proposition}

\begin{proof}
Formula~\eqref{eq:hz-KM} follows from $K=R_{1212}/\det M$ with the Christoffel
symbols of Proposition~\ref{prop:hz-dp-metric} (the computation was checked
symbolically). Since $M$ depends on $\delta=\theta_1-\theta_2$ only, it is
invariant under the simultaneous rotation $\theta_i\mapsto\theta_i+s$, whose
generator is $\partial_{\theta_1}+\partial_{\theta_2}$; the conserved quantity
$M(\partial_{\theta_1}+\partial_{\theta_2},\dot\theta)$ is $J$. A $2$-degree-of-freedom
system with two independent commuting integrals (here the speed and $J$) is
integrable.
\end{proof}

\begin{remark}[Negative curvature without chaos]\label{rem:hz-gb-free}
By~\eqref{eq:hz-KM}, $K_M>0$ for $|\delta|<\pi/2$ and $K_M<0$ for $|\delta|>\pi/2$.
This is forced: the Gauss--Bonnet theorem for the torus gives
$\int_{\Tor^2}K_M\,\ud A=2\pi\chi(\Tor^2)=0$, so a non-flat metric on $\Tor^2$ must
have regions of both signs (check directly: $\ud A=\sqrt{\det M}\,\ud\theta_1\ud\theta_2$
and $K_M\sqrt{\det M}\propto\cos\delta\,(m_1+m_2\sin^2\delta)^{-3/2}$ integrates to zero
over $\delta\in[0,2\pi]$). Yet the geodesic flow is integrable. Negative curvature somewhere is thus
\emph{not} an indicator of chaos; we shall need this caveat in
Section~\ref{sec:hz-jacobi}.
\end{remark}

\section{Switching on gravity: hidden masses}
\label{sec:hz-hidden}

With gravity the double pendulum is the Lagrangian system with
$L=\tfrac12M_{ab}\dot\theta^a\dot\theta^b-V$,
\begin{equation}\label{eq:hz-V}
  V(\theta_1,\theta_2)=-(m_1+m_2)\,g\,l_1\cos\theta_1-m_2\,g\,l_2\cos\theta_2 ,
\end{equation}
and its equations of motion are
\begin{equation}\label{eq:hz-dp-grav}
  \ddot\theta^a+\Gamma^a_{bc}\dot\theta^b\dot\theta^c=-M^{ab}\partial_bV ,
\end{equation}
that is, \eqref{eq:hz-dp-free} with the additional terms
\begin{align*}
  \ddot\theta_1&\ni\frac{m_2\,g\sin\theta_2\cos\delta-(m_1+m_2)\,g\sin\theta_1}{l_1\,(m_1+m_2\sin^2\delta)},\\
  \ddot\theta_2&\ni\frac{(m_1+m_2)\,g\,(\sin\theta_1\cos\delta-\sin\theta_2)}{l_2\,(m_1+m_2\sin^2\delta)} .
\end{align*}
The right-hand side of~\eqref{eq:hz-dp-grav} is a force. For a Hertzian it
must be an appearance.

\subsection{Hertz's response and a representation theorem}

Hertz's answer is that forces are the effect of motions we do not see. The
visible system is coupled by rigid connections to a hidden system; the
combined system is free; and the visible variables, which alone are
observed, move as if acted on by forces \cite{hdp:hertz1894,hdp:lutzen2005}.
In the language of later analytical mechanics the mechanism is familiar: a
hidden coordinate that is \emph{cyclic} can be eliminated by the Routh
procedure \cite{hdp:routh1877}, leaving a visible system with an extra
potential term (Helmholtz's monocyclic systems are the thermodynamic
instance \cite{hdp:helmholtz1884}). The following theorem shows that this
device is not merely available but \emph{universal} for potentials.

\begin{theorem}[Cyclic representation of a potential]\label{thm:hz-cyclic}
Let $(Q,M)$ be a Riemannian manifold, $V\in C^\infty(Q)$ with $V>0$, and
$p>0$ a constant. On $Q\times\Sph^1$ with coordinates $(q^a,\varphi)$ put
\begin{equation}\label{eq:hz-warped}
  G=M_{ab}(q)\,\ud q^a\ud q^b+f(q)^2\,\ud\varphi^2,\qquad f=\frac{p}{\sqrt{2V}} .
\end{equation}
\begin{enumerate}[label=(\alph*),nosep]
  \item If $(q(t),\varphi(t))$ is a geodesic of $G$ with $f^2\dot\varphi=p$, then
        $q(t)$ solves $\ddot q^a+\Gamma^a_{bc}\dot q^b\dot q^c=-M^{ab}\partial_bV$ and
        $\tfrac12M_{ab}\dot q^a\dot q^b+V=\tfrac12G(\dot\gamma,\dot\gamma)$ is constant.
  \item Conversely, if $q(t)$ solves the equation in (a), then
        $\varphi(t)=\varphi_0+\int_0^t 2V(q(s))/p\;\ud s$ makes
        $(q,\varphi)$ a geodesic of $G$ with $f^2\dot\varphi=p$.
\end{enumerate}
\end{theorem}

\begin{proof}
Since $G$ is block diagonal and $\varphi$ does not appear in it, the only
Christoffel symbols of $G$ that differ from those of $M$ and from zero are
\[
  \Gamma^a_{\varphi\varphi}=-\tfrac12M^{ab}\partial_b(f^2),\qquad
  \Gamma^\varphi_{a\varphi}=\Gamma^\varphi_{\varphi a}=\partial_af/f ,
\]
while $\Gamma^a_{b\varphi}=\Gamma^\varphi_{ab}=\Gamma^\varphi_{\varphi\varphi}=0$.
The geodesic equations are therefore
\[
  \ddot q^a+\Gamma^a_{bc}\dot q^b\dot q^c-\tfrac12M^{ab}\partial_b(f^2)\,\dot\varphi^2=0,
  \qquad
  \ddot\varphi+2\,\frac{\partial_af}{f}\,\dot q^a\dot\varphi=0 .
\]
The second equation is $\frac{\ud}{\ud t}(f^2\dot\varphi)=0$, the conservation of
the cyclic momentum. Insert $\dot\varphi=p/f^2$ in the first:
\[
  \tfrac12M^{ab}\partial_b(f^2)\,\frac{p^2}{f^4}
  =-\tfrac12M^{ab}\partial_b\!\Bigl(\frac{p^2}{f^2}\Bigr)
  =-M^{ab}\partial_b V ,
\]
because $p^2/(2f^2)=V$. This proves (a); the energy statement follows from
$\tfrac12f^2\dot\varphi^2=p^2/(2f^2)=V$. Part (b) reverses the steps: define
$\varphi$ by $\dot\varphi=p/f^2=2V/p$, so the second geodesic equation holds, and
the first reduces to the given equation.
\end{proof}

The theorem is Routh reduction run backwards: $R=\tfrac12M\dot q\dot q-p^2/(2f^2)$
is the Routhian of $G$, and $p^2/(2f^2)$ is its effective potential.
(A metric of the form $g+f^2\ud\varphi^2$ is also what Kaluza--Klein
theory calls a circle compactification with a scalar $f$ and no gauge field;
the potential $p^2/(2f^2)$ is the usual momentum-induced one.)
The physical reading, in Hertz's terms, is that \emph{potential energy is the
kinetic energy of the hidden rotor}: $V=\tfrac12f^2\dot\varphi^2$.

\begin{example}[The simple pendulum as a Clairaut problem]\label{ex:hz-simple}
For a pendulum of mass $m$ and length $l$, $Q=\Sph^1$ with $M=ml^2\ud\theta^2$ and
$V=mgl(1-\cos\theta)+\varepsilon$, $\varepsilon>0$ (a constant shift changes only
the value of the conserved energy). Then \eqref{eq:hz-warped} is a metric on a
torus of revolution with profile $f(\theta)=p/\sqrt{2V(\theta)}$,
and the pendulum's motion is a geodesic of that surface. The cyclic
momentum $p=f^2\dot\varphi$ is Clairaut's constant: it equals $f$ times the
component of the velocity along the parallel. In terms of arclength
$s=\sqrt m\,l\,\theta$ the Gaussian curvature of the surface is $-f_{ss}/f$.
The oscillating and rotating regimes of the pendulum are the geodesics that
oscillate about, or wind around, the neck of the surface.
\end{example}

\begin{figure}[t]
\centering
\begin{tikzpicture}[scale=1.0]
  \draw[->] (-0.3,-2.0) -- (7.3,-2.0) node[right] {$q\in Q$};
  \draw[thick,domain=0.2:6.8,samples=80,smooth] plot (\x,{0.85+0.4*sin(58*\x)});
  \draw[thick,domain=0.2:6.8,samples=80,smooth] plot (\x,{-(0.85+0.4*sin(58*\x))});
  \foreach \x in {0.6,1.6,2.6,3.6,4.6,5.6,6.6}{
    \pgfmathsetmacro{\r}{0.85+0.4*sin(58*\x)}
    \draw (\x,0) ellipse ({0.17} and {\r});
    \draw[dotted] (\x,-2.0) -- (\x,{-\r});
  }
  \node at (0.6,1.55) {$f(q)$};
  \draw[->] (3.0,1.25) -- (3.0,0.02);
  \node[right] at (3.0,1.35) {\small hidden circle, angle $\varphi$};
\end{tikzpicture}
\caption{Schematic of Theorem~\ref{thm:hz-cyclic}: over each configuration
$q$ of the visible system sits a hidden circle of radius $f(q)=p/\sqrt{2V(q)}$.
Free (geodesic) motion on the total space projects to motion in the potential
$V$; the hidden circle spins faster where the radius is smaller, that is,
where the potential is higher.}
\label{fig:hz-warped}
\end{figure}

\subsection{Application to the double pendulum}

\begin{corollary}[Hertzian double pendulum]\label{cor:hz-dp}
Let $V_c=V+V_0$ with $V$ given by~\eqref{eq:hz-V} and a constant
$V_0>(m_1+m_2)gl_1+m_2gl_2$, and let $M$ be as in
Proposition~\ref{prop:hz-dp-metric}. On $\Tor^3=\Tor^2\times\Sph^1$ the metric
\[
  G=M_{ab}\,\ud\theta^a\ud\theta^b+\frac{p^2}{2V_c(\theta)}\,\ud\varphi^2
\]
makes Hertz's law (geodesic motion with constant speed) equivalent to the
double pendulum~\eqref{eq:hz-dp-grav} in the sense of
Theorem~\ref{thm:hz-cyclic}, with $\tfrac12G(\dot\gamma,\dot\gamma)=E+V_0$,
where $E=\tfrac12M\dot\theta\dot\theta+V$ is the pendulum's energy.
The hidden rotor turns with angular velocity
$\dot\varphi=2V_c(\theta(t))/p$.
\end{corollary}

The constant $V_0$ is necessary only because $f=p/\sqrt{2V_c}$ requires
$V_c>0$ and $V$ attains negative values; it shifts the zero of potential
energy and so, by Theorem~\ref{thm:hz-cyclic}, only the numerical value of the
total geodesic energy. All the double-pendulum quantities in the following
section are unchanged by it.

\subsection{How literal are the hidden masses?}

Hertz demands more than a Riemannian metric: hidden \emph{masses}, with
rigid connections to the visible ones. Three observations clarify what the
representation theorem does and does not give.

\paragraph{(a) A literal rotor, with an inertia correction.}
Let a hidden particle of unit mass be constrained to the position
$\mathbf y=(r(\theta)\cos\varphi,\,r(\theta)\sin\varphi,\,0)$, a point on a circle
whose radius is a prescribed function of the pendulum angles (a cam or linkage
in a real mechanism). The total ambient metric restricted to the constraint manifold is
\[
  G=\bigl(M^{\rm bare}_{ab}+\partial_ar\,\partial_br\bigr)\ud\theta^a\ud\theta^b+r^2\ud\varphi^2 ,
\]
so the hidden mass also contributes to the inertia of the visible variables,
and by Theorem~\ref{thm:hz-cyclic} (its proof uses only $V=p^2/(2f^2)$) the visible motion has effective metric
$M^{\rm bare}+\ud r\otimes\ud r$ and potential $p^2/(2r^2)$. Choosing $r=p/\sqrt{2V_c}$ gives
$\ud r\otimes\ud r=\dfrac{p^2}{8V_c^{3}}\,\ud V\otimes\ud V$, which is $O(V_0^{-3})$ at fixed $p$.
Thus the literal rotor reproduces the double pendulum with a bare inertia that differs from
the observed one by a correction that vanishes as the hidden energy $V_c$ grows relative to
the variation of $V$. Exactness is recovered only if one declares the \emph{effective}
inertia $M$ to be the observed one. The metric $G$ of Corollary~\ref{cor:hz-dp}
encodes precisely this declaration.

\paragraph{(b) Nash embedding: realizability at the level of geometry.}
By Nash's theorem any compact Riemannian manifold embeds isometrically in
some Euclidean space \cite{hdp:nash1956}. Applied to $(\Tor^3,G)$ this exhibits
the Hertzian configuration space as the constraint manifold of a system of
unit-mass coordinates with holonomic rigid connections. But in such a
realization the pendulum angles $\theta^a$ are merely functions on the
constraint manifold, not the positions of two identifiable bobs. The theorem
shows that Hertz's \emph{geometry} is always realizable; it does not show that
the hidden masses can be chosen compatibly with the visible ones.

\paragraph{(c) Non-uniqueness.}
Nothing in the visible motion selects the hidden structure. Replacing $V_c$ by
a different positive function changes only the energy bookkeeping; adding more
hidden rotors gives $V=\sum_kp_k^2/(2f_k^2)$ with many solutions; and the Jacobi
and Eisenhart constructions of Sections~\ref{sec:hz-jacobi}
and~\ref{sec:hz-compare} reproduce the same trajectories with entirely
different hidden content. The Hertzian explanation of gravity is therefore not
testable by observing the pendulum.

\section{A second geometrization: the Jacobi metric}
\label{sec:hz-jacobi}

Hertz enlarges the configuration space. There is also a way to keep its
dimension and to absorb the potential into the metric by a conformal
rescaling, at the price of fixing the energy.

\begin{theorem}[Jacobi]\label{thm:hz-jacobi}
On the Hill region $U_E=\{q\in Q: V(q)<E\}$ let
$J_E=(E-V)\,M$. The solutions of Newton's equations $\ddot q^a+\Gamma^a_{bc}\dot q^b\dot q^c=-M^{ab}\partial_bV$
of energy $E$ are, after the reparametrization $\ud\tau=(E-V)\,\ud t$, the
geodesics of $(U_E,J_E)$.
\end{theorem}

This is the geometric content of Maupertuis's principle: the abbreviated action
\[
  \int\sqrt{2(E-V)}\;\ud s_M
\]
is, up to the factor $\sqrt2$, the $J_E$-length of the path
\cite{hdp:jacobi1866,hdp:arnold1989}. (The relation of this principle to Hamilton's and
Hertz's is treated in the chapter on variational principles.)

\subsection{Curvature of the Jacobi metric}

In two dimensions the stability of nearby geodesics is governed by the Gaussian
curvature: for a unit-speed geodesic of a surface, the normal deviation $\xi$
between neighbouring geodesics obeys $\xi''+K\,\xi=0$ (Jacobi--Levi-Civita
equation), exponentially growing where $K<0$.

\begin{proposition}[Conformal formula]\label{prop:hz-conformal}
For $J_E=\Omega\,M$ with $\Omega=E-V>0$ on a surface $(Q,M)$,
\begin{equation}\label{eq:hz-KJ}
  K_{J_E}\;=\;\frac{1}{\Omega}\Bigl(K_M-\tfrac12\,\Lap_M\ln\Omega\Bigr),
\end{equation}
where $\Lap_M$ is the Laplace--Beltrami operator of $M$. Near a regular point of the
Hill boundary $\{V=E\}$, $K_{J_E}\simeq|\nabla\Omega|_M^2/(2\Omega^3)\to+\infty$.
\end{proposition}

\begin{proof}
For $\hat g=e^{2u}g$ on a surface, $\hat K=e^{-2u}(K-\Lap_gu)$ \cite{hdp:docarmo1992}.
Put $e^{2u}=\Omega$. For the boundary behaviour,
$\Lap\ln\Omega=\Lap\Omega/\Omega-|\nabla\Omega|^2/\Omega^2$, and the second term dominates.
(Formula~\eqref{eq:hz-KJ} was also checked symbolically against the Riemann tensor of $J_E$ for the double pendulum.)
\end{proof}

For the double pendulum $K_M$ is~\eqref{eq:hz-KM}, $\Omega=E-V$ with $V$
from~\eqref{eq:hz-V}, and~\eqref{eq:hz-KJ} gives $K_{J_E}$ in closed form.

\begin{proposition}[Negative curvature is unavoidable above the separatrix energy]\label{prop:hz-gb}
Let $E>V_{\max}=(m_1+m_2)gl_1+m_2gl_2$. Then $U_E=\Tor^2$, $J_E$ is a smooth
Riemannian metric on the torus, and $\int_{\Tor^2}K_{J_E}\,\ud A_{J_E}=0$. Hence, unless
$K_{J_E}\equiv0$, it takes negative values somewhere.
\end{proposition}

\begin{proof}
$V<V_{\max}<E$ everywhere, so $\Omega>0$ on all of $\Tor^2$. Apply Gauss--Bonnet
with $\chi(\Tor^2)=0$. (For $m_1=m_2=l_1=l_2=g=1$ one finds numerically
$K_{J_4}\in[-1.0,\,0.30]$, so $K_{J_E}\not\equiv0$.)
\end{proof}

\subsection{Where the negative curvature sits}

Figure~\ref{fig:hz-curvature} shows the sign of $K_{J_E}$ for the
unit-parameter pendulum ($m_1=m_2=l_1=l_2=g=1$, so $V=-2\cos\theta_1-\cos\theta_2$,
$-3\le V\le3$) at three energies, and Table~\ref{tab:hz-curv} the fraction of
the (Jacobi-)area where it is negative.

\begin{figure}[t]
\centering
\includegraphics[width=\textwidth]{jacobi_curvature.pdf}
\caption{Sign of the Gaussian curvature $K_{J_E}$ of the Jacobi metric of the
double pendulum on the torus $(\theta_1,\theta_2)\in[-\pi,\pi)^2$, for
$m_1=m_2=l_1=l_2=g=1$. Blue: $K_{J_E}>0$; red: $K_{J_E}<0$; grey: inaccessible
region $V\ge E$. At $E=-2$ the accessible region is a disc about the stable
equilibrium and $K_{J_E}>0$ throughout. At $E=1$ a red region has opened; at
$E=4>V_{\max}=3$ the whole torus is accessible and, as
Proposition~\ref{prop:hz-gb} requires, the curvature has both signs.}
\label{fig:hz-curvature}
\end{figure}

\begin{table}[t]
\centering
\begin{tabular}{@{}rcc@{}}
\toprule
$E$ & accessible part of the torus & fraction of $J_E$-area with $K_{J_E}<0$\\
\midrule
$-2$ & $0.13$ & $0$\\
$1$  & $0.69$ & $0.20$\\
$4$  & $1$    & $0.44$\\
\bottomrule
\end{tabular}
\caption{Accessible fraction of the coordinate torus and fraction of the Jacobi area
$\sqrt{\det J_E}\,\ud\theta_1\ud\theta_2$ on which the curvature is negative, for the
unit-parameter double pendulum (grid evaluation, two significant digits).}
\label{tab:hz-curv}
\end{table}

\begin{remark}[What the picture does and does not say]\label{rem:hz-caveat}
The pattern is the expected one. At low energy the motion is a perturbed
small oscillation and the whole accessible region has $K_{J_E}>0$; as energy
rises, negative-curvature regions open and grow, and above the energy at which
the rods can pass through the upright position the topology forces them to exist.
But three cautions are required.
(i) $K<0$ is sufficient for local exponential divergence of nearby geodesics, not necessary:
oscillating positive curvature can produce instability by parametric resonance, which is the
mechanism of the Hannay-angle chapter, and this is how chaos arises in many
systems whose average curvature is positive \cite{hdp:casetti1996,hdp:pettini2007}.
(ii) Remark~\ref{rem:hz-gb-free} shows that a torus metric with negative-curvature regions can
have integrable geodesic flow; the sign of $K$ locates \emph{possible} instability,
it does not by itself establish chaos.
(iii) $J_E$ degenerates at the Hill boundary and its curvature diverges there
(Proposition~\ref{prop:hz-conformal}), so the picture below the separatrix energy is a statement
about the interior of $U_E$ only.
\end{remark}

\section{Three geometrizations side by side}
\label{sec:hz-compare}

A third construction due to Eisenhart lifts the dynamics to a \emph{Lorentzian}
space. On $Q\times\Rb_t\times\Rb_z$ take
\begin{equation}\label{eq:hz-eisenhart}
  G_{\rm E}=M_{ab}\,\ud q^a\ud q^b+2\,\ud t\,\ud z-2V(q)\,\ud t^2 .
\end{equation}
The $t$-equation of the geodesic flow gives $\ddot t=0$, so $t$ is an affine
parameter, and the $q$-equation is $\ddot q^a+\Gamma^a_{bc}\dot q^b\dot q^c=-M^{ab}\partial_bV$ (with
$\Gamma^a_{tt}=M^{ab}\partial_bV$) \cite{hdp:eisenhart1929}. Table~\ref{tab:hz-compare} compares the three.

\begin{table}[t]
\centering
\small
\begin{tabular}{@{}lccc@{}}
\toprule
 & Hertz (hidden rotor) & Jacobi & Eisenhart\\
\midrule
Signature & Riemannian & Riemannian & Lorentzian\\
Dimension & $n+1$ & $n$ & $n+2$\\
Fixed data & hidden momentum $p$ & energy $E$ & none\\
Restriction on $V$ & $V+V_0>0$ & Hill region $V<E$ & none\\
Geodesic parameter & physical time $t$ & $\tau$, $\ud\tau=(E-V)\ud t$ & physical time $t$\\
Where $V$ sits & hidden kinetic energy & conformal factor & $G_{tt}$\\
Hertzian reading & hidden masses & none & none\\
\bottomrule
\end{tabular}
\caption{Three ways to turn the potential force into geometry. Only Hertz's
construction has a mechanical reading in terms of positive masses.}
\label{tab:hz-compare}
\end{table}

% ======================================================================
%  Section insert: Eisenhart's lifts in detail  (+ comparison table)
%  Drop into hertz_chapter.tex in place of the old "Three geometrizations
%  side by side" section (i.e. before \section{Assessment...}).
%
%  Needs: amsmath, amsthm (definition/theorem/proposition/remark/example
%  environments as defined in the chapter), enumitem, booktabs.
%  Cross-references used: thm:hz-cyclic, sec:hz-hidden, sec:hz-jacobi,
%  prop:hz-dp-metric, eq:hz-V, eq:hz-KM, eq:hz-warped.
%  Bibliography keys used: hdp:eisenhart1929, hdp:duval1985,
%  hdp:duval1991, hdp:casetti2000.
% ======================================================================

\section{Eisenhart's lifts in detail}
\label{sec:hz-eisenhart}

In 1928--29 L.~P.~Eisenhart showed that the trajectories of a natural
Lagrangian system can be obtained as geodesics of a metric on a larger space
\cite{hdp:eisenhart1929}. His paper contains \emph{two} constructions,
and they stand in different relations to Hertz. The first is
Lorentzian, works for every potential (including time-dependent ones) and
needs two extra coordinates; the second is Riemannian, needs a positive
potential and one extra coordinate. The second turns out to be
the metric of Theorem~\ref{thm:hz-cyclic}. The first was rediscovered
in 1984--85 by Duval, Burdet, K\"unzle and Perrin, who called the resulting
spaces \emph{Bargmann structures} \cite{hdp:duval1985,hdp:duval1991}, and
the construction is now often called the Eisenhart--Duval lift.

\subsection{The Lorentzian (null) lift}

\begin{definition}[Eisenhart lift]\label{def:hz-eis}
Let $(Q,M)$ be a Riemannian manifold and $V\in C^\infty(Q)$ arbitrary. On
$\widehat Q=Q\times\Rb_t\times\Rb_z$ put
\begin{equation}\label{eq:hz-eis-metric}
  G_{\rm E}=M_{ab}(q)\,\ud q^a\ud q^b+2\,\ud t\,\ud z-2V(q)\,\ud t^2 .
\end{equation}
\end{definition}

The $(t,z)$ block $\bigl(\begin{smallmatrix}-2V&1\\1&0\end{smallmatrix}\bigr)$
has determinant $-1$ whatever $V$ is, so $G_{\rm E}$ is a non-degenerate metric
of Lorentzian signature $(n{+}1,1)$. Its inverse has the components
\begin{equation}\label{eq:hz-eis-inverse}
  G^{ab}=M^{ab},\qquad G^{tz}=1,\qquad G^{zz}=2V,\qquad G^{tt}=G^{ta}=G^{za}=0 .
\end{equation}
The only Christoffel symbols that differ from those of $M$ and from zero are
\begin{equation}\label{eq:hz-eis-christoffel}
  \Gamma^a_{tt}=M^{ab}\partial_bV,\qquad
  \Gamma^z_{at}=\Gamma^z_{ta}=-\partial_aV .
\end{equation}
(For the double pendulum this was checked symbolically for a flat base and
confirmed for the curved base $M$ through the Ricci tensor below.)

\begin{theorem}[Eisenhart]\label{thm:hz-eis}
A curve $\gamma(s)=(q(s),t(s),z(s))$ is a geodesic of $G_{\rm E}$ if and only if
\begin{equation}\label{eq:hz-eis-geod}
  \ddot t=0,\qquad
  \ddot q^a+\Gamma^a_{bc}\dot q^b\dot q^c=-M^{ab}\partial_bV\,\dot t^{\,2},\qquad
  \ddot z=2\,\dot t\,\partial_aV\,\dot q^a .
\end{equation}
Consequently, if $\dot t=1$ (so $s=t$ is physical time):
\begin{enumerate}[label=(\alph*),nosep]
  \item $q(t)$ solves Newton's equations $\ddot q^a+\Gamma^a_{bc}\dot q^b\dot q^c=-M^{ab}\partial_bV$;
  \item $\dot z=2V(q)-\varepsilon$ with a constant $\varepsilon$, so $p_t:=G(\partial_t,\dot\gamma)=-\varepsilon$;
  \item $G_{\rm E}(\dot\gamma,\dot\gamma)=2(E-\varepsilon)$, where $E=\tfrac12M\dot q\dot q+V$ is the energy of $q(t)$;
  \item $\gamma$ is a \emph{null} geodesic exactly when $\varepsilon=E$, and then
        $z(t)=z_0-\int_0^t L\,\ud t'=z_0-S[q]$ with $L=\tfrac12M\dot q\dot q-V$.
\end{enumerate}
Conversely every solution $q(t)$ of Newton's equations lifts, for every
choice of $\varepsilon$, $z_0$ and $t_0$, to a geodesic with $\dot t=1$.
\end{theorem}

\begin{proof}
Equations~\eqref{eq:hz-eis-geod} are the geodesic equations with the symbols
\eqref{eq:hz-eis-christoffel}: $\Gamma^t_{AB}=0$ gives $\ddot t=0$;
$\ddot q^a+\Gamma^a_{bc}\dot q^b\dot q^c+\Gamma^a_{tt}\dot t^2=0$ is the second;
$\ddot z+2\Gamma^z_{at}\dot q^a\dot t=0$ is the third. For $\dot t=1$ the third says
$\ddot z=2\,\ud V/\ud t$, hence $\dot z=2V-\varepsilon$, and $p_t=-2V\dot t+\dot z=-\varepsilon$.
Then
\[
  G(\dot\gamma,\dot\gamma)=M\dot q\dot q+2\dot z-2V
  =2T+2(2V-\varepsilon)-2V=2(T+V-\varepsilon)=2(E-\varepsilon),
\]
which gives (c) and the first half of (d). If $\varepsilon=E$ then
$\dot z=2V-E=V-T=-L$. The converse is immediate by defining
$z$ through $\dot z=2V-\varepsilon$.
\end{proof}

\begin{remark}[The extra coordinate is the action]\label{rem:hz-z-action}
Statement (d) is the geometric content of the lift: along the physical
(null) geodesics the new coordinate $z$ is minus Hamilton's principal
function evaluated along the path. This is no accident. With $\dot t=1$ the
geodesic Lagrangian is $\tfrac12G(\dot\gamma,\dot\gamma)=L+\dot z$, so
$\int\tfrac12G(\dot\gamma,\dot\gamma)\,\ud t=S[q]+z(t_2)-z(t_1)$:
the energy functional of the lift is Hamilton's action up to endpoint terms.
\end{remark}

\subsection{Hamiltonian form and conserved quantities}

The geodesic Hamiltonian $H_{\rm E}=\tfrac12G^{AB}p_Ap_B$ is, by~\eqref{eq:hz-eis-inverse},
\begin{equation}\label{eq:hz-eis-ham}
  H_{\rm E}=\tfrac12M^{ab}p_ap_b+p_tp_z+V(q)\,p_z^{\,2}.
\end{equation}
Both $t$ and $z$ are cyclic for time-independent $V$, so $p_t$ and $p_z$ are conserved.
With $p_z=G(\partial_z,\dot\gamma)=\dot t=1$ and
$H(q,p)=\tfrac12M^{ab}p_ap_b+V$ the physical Hamiltonian, \eqref{eq:hz-eis-ham} becomes
\[
  H_{\rm E}=H(q,p)+p_t ,
\]
and the null condition $H_{\rm E}=0$ is $p_t=-H=-E$: Hamilton's equations of $H_{\rm E}$
are the equations of the \emph{extended phase space}, in which energy is the momentum conjugate
to time and the dynamics is a constraint. The constant $p_z$ plays the role of a mass
(in the Bargmann picture it is the mass or the electric charge of the particle), and the
fixed value $p_z=1$ is what selects Newtonian dynamics with unit
inertia. Hamilton's equations give $\dot t=\partial H_{\rm E}/\partial p_t=p_z=1$ and
$\dot z=\partial H_{\rm E}/\partial p_z=p_t+2V=2V-E$, as found above.

\subsection{Geometry of the lift: a pp-wave}

\begin{proposition}[Curvature of the Eisenhart metric]\label{prop:hz-eis-curv}
\leavevmode
\begin{enumerate}[label=(\roman*),nosep]
  \item $\nabla\,\partial_z=0$: the null vector field $\partial_z$ is covariantly constant.
  \item The Ricci tensor of $G_{\rm E}$ has only the components
        \[
          \operatorname{Ric}_{ab}=\operatorname{Ric}^{M}_{ab},\qquad
          \operatorname{Ric}_{tt}=\Lap_MV .
        \]
\end{enumerate}
\end{proposition}

\begin{proof}
(i) is $\Gamma^A_{Bz}=0$, read off from \eqref{eq:hz-eis-christoffel}. For (ii) I computed the Ricci
tensor symbolically for a flat base with arbitrary $V(x,y)$ and, for the
unit-parameter double pendulum with its curved metric $M$, at random points, where the
only non-vanishing components were $\operatorname{Ric}_{ab}=K_MM_{ab}$ and
$\operatorname{Ric}_{tt}=\Lap_MV$ to rounding error; the general statement is the standard curvature
property of Eisenhart metrics \cite{hdp:casetti2000}.
\end{proof}

A Lorentzian manifold with a covariantly constant null vector whose metric has this
block form is a (generalized) \emph{pp-wave}, and the structure
$(\widehat Q,G_{\rm E},\partial_z)$ is a Bargmann structure \cite{hdp:duval1985,hdp:duval1991}.
Two consequences are worth stating.
\begin{itemize}[nosep]
  \item \emph{Gravity as curvature.} With $\Lap_MV=4\pi G\rho$ for a Newtonian
        potential, $\operatorname{Ric}_{tt}=\Lap_MV$ is the Poisson equation, and Ricci-flatness of the lift
        (flat $M$, harmonic $V$) is Newtonian gravity in vacuum. The lift makes Newtonian
        gravity \emph{look} like general relativity; it is a statement about a Lorentzian
        space of one dimension more than spacetime, not about spacetime.
  \item \emph{Tidal tensor.} The geodesic deviation equation of $G_{\rm E}$ along a lifted
        trajectory, restricted to the $q$-directions, is the usual variational equation of the
        dynamics,
        \[
          \frac{D^2\xi^a}{\ud t^2}+R^{M\,a}{}_{bcd}\,\dot q^b\xi^c\dot q^d
          +M^{ab}\nabla_b\nabla_cV\,\xi^c=0,
        \]
        and the trace of the tidal tensor is $\Lap_MV$ (the origin of $\operatorname{Ric}_{tt}$)
        \cite{hdp:casetti2000}. In contrast to the Jacobi metric of
        Section~\ref{sec:hz-jacobi}, no factor $E-V$ enters and there is no Hill-boundary singularity;
        but the price is that the criterion for exponential instability is the standard
        linear-stability computation, not a purely geometric sign condition.
\end{itemize}

\subsection{Examples}

\begin{example}[Harmonic oscillator]\label{ex:hz-eis-ho}
For $Q=\Rb$, $M=\ud q^2$ and $V=\tfrac12\omega^2q^2$,
$G_{\rm E}=\ud q^2+2\,\ud t\,\ud z-\omega^2q^2\ud t^2$ is a plane wave with $\operatorname{Ric}_{tt}=\omega^2$.
The null geodesics with $\dot t=1$ are $q=A\cos\theta$, $\theta=\omega t+\phi$, and
\[
  z=z_0+\tfrac14A^2\omega\,\sin2\theta ,
\]
which is $z_0-\int L\,\ud t$ with $L=-\tfrac12A^2\omega^2\cos2\theta$. The $z$-motion is a
bounded oscillation at twice the frequency.
\end{example}

\begin{example}[Double pendulum]\label{ex:hz-eis-dp}
With $M$ from Proposition~\ref{prop:hz-dp-metric} and $V$ from~\eqref{eq:hz-V},
\[
\begin{aligned}
  G_{\rm E}={}&(m_1+m_2)\,l_1^2\,\ud\theta_1^2+2m_2l_1l_2\cos(\theta_1-\theta_2)\,\ud\theta_1\ud\theta_2
  +m_2l_2^2\,\ud\theta_2^2\\
  &+2\,\ud t\,\ud z+2g\bigl[(m_1+m_2)l_1\cos\theta_1+m_2l_2\cos\theta_2\bigr]\ud t^2 .
\end{aligned}
\]
This is a four-dimensional Lorentzian manifold $\Tor^2\times\Rb^2$ with
$\operatorname{Ric}_{ab}=K_MM_{ab}$ (with $K_M$ from~\eqref{eq:hz-KM}) and $\operatorname{Ric}_{tt}=\Lap_MV$.
Its Killing fields include $\partial_z$ (null, parallel) and $\partial_t$; the rotation
$\partial_{\theta_1}+\partial_{\theta_2}$ that was a symmetry of the free pendulum is
broken by $V$. The null geodesics with $\dot t=1$ project to the chaotic trajectories of the pendulum,
and $z$ along them is $-S[\theta]$.
\end{example}

\subsection{The Riemannian lift, and its identity with the hidden rotor}

Eisenhart's second construction, for \emph{time-independent potentials that are positive} on $Q$,
uses one extra coordinate and a Riemannian metric of the form
$M+A(q)\,\ud z^2$ with $A$ built from $V$
\cite{hdp:eisenhart1929,hdp:casetti2000}. Taking $A=1/V_c$ with $V_c=V+c>0$,
\begin{equation}\label{eq:hz-eis-riem}
  G_{\rm R}=M_{ab}\,\ud q^a\ud q^b+\frac{1}{V_c(q)}\,\ud z^2 ,
\end{equation}
the Hamiltonian is $\tfrac12M^{ab}p_ap_b+\tfrac12V_c\,p_z^2$, which equals
the physical Hamiltonian $\tfrac12M^{ab}p_ap_b+V_c$ for $p_z=\sqrt2$. This is exactly
the metric~\eqref{eq:hz-warped} of Theorem~\ref{thm:hz-cyclic} with $f^2=p^2/(2V_c)$, $p=\sqrt2$,
and $\varphi=z$. The theorem therefore \emph{is} Eisenhart's Riemannian lift, with a different proof and with
a physical interpretation of the extra coordinate added: Eisenhart's $z$ is Hertz's hidden
cyclic coordinate, and $V=\tfrac12f^2\dot z^2$ is the kinetic energy of the hidden rotor.
(Hertz did not write down the warped metric; the idea that a potential is the kinetic energy of
a cyclic hidden motion is his, and the exact geometric statement is Eisenhart's.)

\subsection{Hertzian or not?}

Does the Lorentzian lift satisfy Hertz's requirements? Only in part. Its dynamics is
free in the sense that the law is geodesic, and no force appears. But Hertz's mass metric is positive
definite: kinetic energy is $\tfrac12\sum m\dot x^2\ge0$, and the straightest path is defined by minimizing
a norm of acceleration. In signature $(n{+}1,1)$ there is no such norm, the physical
trajectories are \emph{null}, and the extra coordinate $z$ is a null direction, not a hidden mass.
The Lorentzian lift is thus forceless but not Hertzian; the Riemannian lift is both.
That the two lifts are available side by side underlines the point of
Section~\ref{sec:hz-hidden}(c): the geometric content of the visible dynamics does not fix its higher-dimensional
realization.

\section{Four geometrizations side by side}
\label{sec:hz-compare}

\begin{table}[t]
\centering
\small
\setlength{\tabcolsep}{4pt}
\begin{tabular}{@{}lccc@{}}
\toprule
 & Hertz / Eisenhart--Riem. & Jacobi & Eisenhart--Lorentz\\
\midrule
Signature & Riemannian & Riemannian & Lorentzian $(n{+}1,1)$\\
Dimension & $n+1$ & $n$ & $n+2$\\
Fixed data & $p_z$ (hidden momentum) & energy $E$ & $p_z=1$, null condition\\
Restriction on $V$ & $V+c>0$ & Hill region $V<E$ & none\\
Geodesic parameter & physical time $t$ & $\tau$, $\ud\tau=(E-V)\ud t$ & physical time $t$\\
Where $V$ sits & hidden kinetic energy & conformal factor & $G_{tt}$\\
Role of extra coordinates & hidden angle $z$ & none & $t$ and null $z=-S$\\
Time-dependent $V$ & no & no & yes\\
Hertzian reading & hidden masses & none & none (null)\\
\bottomrule
\end{tabular}
\caption{Three ways (four constructions, since the first column is Hertz's idea and Eisenhart's
Riemannian metric at once) to turn the potential force into geometry. Only the first has a
mechanical reading in terms of positive masses.}
\label{tab:hz-compare}
\end{table}

Table~\ref{tab:hz-compare} summarizes the comparison. The Lorentzian lift is the only one that
treats time-dependent potentials uniformly, since $t$ is a coordinate of the lifted
space; the Jacobi construction requires conservation of energy and the Riemannian lift
requires a time-independent positive $V$.

\section{Assessment: what Hertz's program delivers}
\label{sec:hz-assess}

\paragraph{What it gains.}
Applied to the double pendulum the program does what it promises, in one
respect: it removes every force from the statement of the law. The tension in
the rods is a reaction of the rigid connection, the centrifugal-looking terms
are Christoffel symbols, and gravity is the kinetic energy of a rotor. The
geometry also yields something new that is invisible in the force picture:
the Gauss--Bonnet argument of Proposition~\ref{prop:hz-gb} explains, without any
numerics, why the Jacobi geometry of the double pendulum above the upright
energy must contain regions in which nearby motions diverge.

\paragraph{What it costs.}
The cost is the hidden structure, and the analysis of Section~\ref{sec:hz-hidden}
shows that it is both inessential and unconstrained. The visible motion
of the pendulum determines the ordinary differential equation~\eqref{eq:hz-dp-grav}
and nothing else; the hidden rotor, the Jacobi metric and the Eisenhart
lift are three unrelated ways of writing that same equation as a geodesic problem,
each in a different dimension and with different auxiliary data. To the extent
that observation of the pendulum is all there is, the hidden masses
are idle wheels. This is the pattern argued throughout this book: the equations
of motion are the primary, real-world content of mechanics, and variational and
geometric formulations are ideal-world devices for finding and organizing them.
The three geometrizations are in this sense charts of one object, each valid
on its own domain and each introducing singularities of its own: the positivity
requirement $V+V_0>0$, the Hill boundary, the Lorentzian signature.

\paragraph{What it cannot reach.}
Hertz's systems are closed and conservative: all forces, including friction,
must be traced to hidden conservative motion. Dissipation therefore enters
only as a statement that the hidden energy is out of sight, and genuine open
systems, such as the damped double pendulum, are outside the scheme as stated;
they are the subject of the chapters on damping and on open systems.
The nonholonomic extension, in which straightest paths differ from shortest
paths, is the other place where Hertz's mechanics is not simply geodesic flow
(Section~\ref{sec:hz-geometry}).

\paragraph{Historical coda.}
The reception was divided between those who found the hidden-mass hypothesis a
barrier to physical content and those who appreciated the geometric form
\cite{hdp:lutzen2005}. What survived was less the philosophy than the geometry:
the representation of constrained mechanics as geodesic flow on a Riemannian
manifold, which is the standard language of the modern geometric treatments.

\section*{Problems}
\addcontentsline{toc}{section}{Problems}

\begin{enumerate}[label=\textbf{\thechapter.\arabic*}.]
  \item Prove Proposition~\ref{prop:hz-law}(i) by Lagrange multipliers: minimize
        $g(a,a)$ subject to $\omega_\alpha(a)=c_\alpha$ and show the multipliers
        are the constraint reactions.
  \item Derive~\eqref{eq:hz-dp-free} directly from Proposition~\ref{prop:hz-law}(iii) with
        the constraints $|\mathbf x_1|=l_1$, $|\mathbf x_2-\mathbf x_1|=l_2$, eliminating the
        multipliers (the rod tensions).
  \item Identify the second integral of the free double pendulum with
        $J$ of Proposition~\ref{prop:hz-dp-curvature}, and find the dependence of the
        period of the $\delta$-motion on $J$ and the speed.
  \item In Example~\ref{ex:hz-simple} show that the surface of revolution is
        smooth, that $K=-f_{ss}/f$ does not depend on $p$, and determine where $K<0$
        (compare Gauss--Bonnet for the torus).
  \item Verify the claim of Section~\ref{sec:hz-hidden}(a) that $\ud r\otimes\ud r=O(V_0^{-3})$
        at fixed $p$, and compute the leading correction to the pendulum's small-oscillation frequency.
  \item Verify that the geodesics of~\eqref{eq:hz-eisenhart} with $\dot t=1$ project to
        solutions of~\eqref{eq:hz-dp-grav}, and find the quantity conserved by the
        null Killing vector $\partial_z$.
  \item Show that for $g\neq0$ the vector field $\partial_{\theta_1}+\partial_{\theta_2}$ is
        \emph{not} a Killing field of $J_E$, so the angular momentum $J$ of
        Proposition~\ref{prop:hz-dp-curvature} is no longer conserved. Why does this loss of
        symmetry not by itself prove that the flow is non-integrable?
\end{enumerate}

\section*{Notes and references}
\addcontentsline{toc}{section}{Notes and references}

\emph{On verification.} The kinetic metric, the geodesic-versus-Lagrange form of the equations of
motion, the explicit accelerations~\eqref{eq:hz-dp-free}--\eqref{eq:hz-dp-grav}, the Christoffel symbols
of the warped metric, the curvature~\eqref{eq:hz-KM} and the conformal
formula~\eqref{eq:hz-KJ} were checked symbolically (SymPy); the curvature map and the table were computed
numerically from~\eqref{eq:hz-KJ} on a $900\times900$ grid, and the Gauss--Bonnet integral at $E=4$
vanishes to rounding error, as it should.

\begingroup
\let\chapter\section
\renewcommand{\bibname}{References}
\begin{thebibliography}{99}
\bibitem{hdp:hertz1894} H.~Hertz, \emph{Die Prinzipien der Mechanik in neuem Zusammenhange dargestellt},
  Gesammelte Werke, Vol.~III (J.~A.~Barth, Leipzig, 1894), with a preface by H.~von Helmholtz.
  English translation: \emph{The Principles of Mechanics Presented in a New Form}, trans.\ D.~E.~Jones and
  J.~T.~Walley (Macmillan, London, 1899; repr.\ Dover, New York, 1956).
\bibitem{hdp:lutzen2005} J.~L\"utzen, \emph{Mechanistic Images in Geometric Form: Heinrich Hertz's
  ``Principles of Mechanics''} (Oxford University Press, Oxford, 2005).
\bibitem{hdp:gauss1829} C.~F.~Gauss, \"Uber ein neues allgemeines Grundgesetz der Mechanik,
  \emph{J.\ reine angew.\ Math.}\ \textbf{4} (1829) 232--235.
\bibitem{hdp:eom-hertz} V.~V.~Rumyantsev, Hertz principle, in \emph{Encyclopedia of Mathematics} (EMS Press).
\bibitem{hdp:routh1877} E.~J.~Routh, \emph{A Treatise on the Stability of a Given State of Motion}
  (Macmillan, London, 1877).
\bibitem{hdp:helmholtz1884} H.~von Helmholtz, Principien der Statik monocyklischer Systeme,
  \emph{J.\ reine angew.\ Math.}\ \textbf{97} (1884) 111--140, 317--336.
\bibitem{hdp:nash1956} J.~Nash, The imbedding problem for Riemannian manifolds,
  \emph{Ann.\ of Math.}\ \textbf{63} (1956) 20--63.
\bibitem{hdp:jacobi1866} C.~G.~J.~Jacobi, \emph{Vorlesungen \"uber Dynamik}, ed.\ A.~Clebsch
  (G.~Reimer, Berlin, 1866).
\bibitem{hdp:arnold1989} V.~I.~Arnold, \emph{Mathematical Methods of Classical Mechanics},
  2nd ed., Graduate Texts in Mathematics \textbf{60} (Springer, New York, 1989).
\bibitem{hdp:docarmo1992} M.~P.~do Carmo, \emph{Riemannian Geometry} (Birkh\"auser, Boston, 1992).
\bibitem{hdp:eisenhart1929} L.~P.~Eisenhart, Dynamical trajectories and geodesics,
  \emph{Ann.\ of Math.}\ \textbf{30} (1929) 591--606.
\bibitem{hdp:casetti1996} L.~Casetti, C.~Clementi and M.~Pettini, Riemannian theory of Hamiltonian
  chaos and Lyapunov exponents, \emph{Phys.\ Rev.\ E} \textbf{54} (1996) 5969--5984.
\bibitem{hdp:pettini2007} M.~Pettini, \emph{Geometry and Topology in Hamiltonian Dynamics and
  Statistical Mechanics} (Springer, New York, 2007).
\end{thebibliography}
\endgroup
